% =============================================================================
% ПЛИС на практике: Sipeed Tang Nano 20K
% Сборка: xelatex main.tex  (запустить дважды)
% =============================================================================
\documentclass[
    fontsize=12pt,
    headings=big,
    numbers=noenddot,
    toc=listof,
    headinclude=true,
    footinclude=false,
    open=any,
    oneside
]{scrreprt}

% =============================================================================
% Преамбула учебника «ПЛИС на практике: Sipeed Tang Nano 20K»
% Компиляция: xelatex main.tex (дважды, для оглавления и ссылок)
% =============================================================================

% --- Языки и шрифты (XeLaTeX) ---
\usepackage{fontspec}
\usepackage{polyglossia}
\setdefaultlanguage{russian}
\setotherlanguage{english}

% На Windows используются Times New Roman / Arial / Consolas,
% на Linux -- метрически совместимые Liberation и DejaVu Sans Mono.
\IfFontExistsTF{Times New Roman}
  {\def\MainFontName{Times New Roman}}{\def\MainFontName{Liberation Serif}}
\IfFontExistsTF{Arial}
  {\def\SansFontName{Arial}}{\def\SansFontName{Liberation Sans}}
\IfFontExistsTF{Consolas}
  {\def\MonoFontName{Consolas}}{\def\MonoFontName{DejaVu Sans Mono}}
\setmainfont{\MainFontName}[Ligatures=TeX]
\setsansfont{\SansFontName}[Ligatures=TeX]
\setmonofont{\MonoFontName}[Scale=MatchLowercase]
% polyglossia требует явно указать шрифты для кириллицы
\newfontfamily\cyrillicfont{\MainFontName}[Ligatures=TeX]
\newfontfamily\cyrillicfontsf{\SansFontName}[Ligatures=TeX]
\newfontfamily\cyrillicfonttt{\MonoFontName}[Scale=MatchLowercase]

\usepackage{microtype}
\emergencystretch=3em

% --- Абзацы ---
\usepackage{indentfirst}
\usepackage{setspace}
\setstretch{1.2}
\setlength{\parindent}{1.25cm}
\setlength{\parskip}{3pt plus 1pt}

% --- Математика ---
\usepackage{amsmath}
\usepackage{amssymb}

% --- Поля ---
\usepackage{geometry}
\geometry{a4paper, top=22mm, bottom=22mm, left=25mm, right=20mm,
          headheight=15mm, footskip=12mm}

% --- Цвета ---
\usepackage[table]{xcolor}
\definecolor{FpgaBlue}{RGB}{0, 91, 187}
\definecolor{FpgaDark}{RGB}{22, 43, 58}
\definecolor{FpgaOrange}{RGB}{255, 140, 0}
\definecolor{FpgaGreen}{RGB}{0, 140, 70}
\definecolor{FpgaRed}{RGB}{200, 30, 45}
\definecolor{FpgaPurple}{RGB}{120, 40, 160}
\definecolor{FpgaLight}{RGB}{242, 246, 252}
\definecolor{FpgaGray}{RGB}{120, 120, 125}
\definecolor{BoardGreen}{RGB}{20, 90, 60}
\definecolor{ChipBlack}{RGB}{35, 35, 40}
\definecolor{CodeBack}{RGB}{248, 249, 251}
\definecolor{CodeComment}{RGB}{0, 128, 0}
\definecolor{CodeKeyword}{RGB}{0, 60, 170}
\definecolor{CodeString}{RGB}{170, 60, 0}

% --- Графика ---
\usepackage{etoolbox}
\usepackage{xstring}
\usepackage{graphicx}
\usepackage{float}
\usepackage{tikz}
\usetikzlibrary{arrows.meta, calc, positioning, shapes.geometric, shapes.misc,
  shapes.gates.logic.US, shapes.gates.logic.IEC, shapes.symbols,
  fit, backgrounds, decorations.pathreplacing, patterns, shadows, matrix}
\usepackage{tikz-timing}
\usepackage{pgfplots}
\pgfplotsset{compat=1.18}
\usepgfplotslibrary{groupplots}
\usepackage[siunitx, RPvoltages]{circuitikz}

% --- Заголовки, колонтитулы ---
\usepackage[automark, headsepline]{scrlayer-scrpage}
\clearpairofpagestyles
\ihead{\headmark}
\ohead{\textbf{\textcolor{FpgaBlue}{TANG NANO 20K}}}
\cfoot{\pagemark}
\setkomafont{pageheadfoot}{\sffamily\small\color{FpgaDark}}
\setkomafont{pagenumber}{\sffamily\bfseries\color{FpgaBlue}}
\setkomafont{chapter}{\sffamily\huge\bfseries\color{FpgaBlue}}
\setkomafont{section}{\sffamily\Large\bfseries\color{FpgaBlue}}
\setkomafont{subsection}{\sffamily\large\bfseries\color{FpgaDark}}
\setkomafont{subsubsection}{\sffamily\normalsize\bfseries\color{FpgaOrange}}
\setkomafont{paragraph}{\sffamily\bfseries\color{FpgaDark}}
\setkomafont{part}{\sffamily\Huge\bfseries\color{FpgaBlue}}
\renewcommand*{\chapterformat}{%
  \colorbox{FpgaBlue}{\color{white}\,\thechapter\,}\enskip}

% --- Списки и таблицы ---
\usepackage{enumitem}
\setlist{itemsep=1pt, parsep=0pt, topsep=3pt}
\setlist[itemize]{label=\textcolor{FpgaBlue}{\textbullet}, leftmargin=*}
\setlist[enumerate]{label=\textcolor{FpgaBlue}{\bfseries\arabic*.}, leftmargin=*}
\setlist[description]{font=\bfseries\color{FpgaDark}}
\usepackage{booktabs}
\usepackage{tabularx}
\usepackage{multirow}
\usepackage{longtable}
\usepackage{caption}
\captionsetup{font=small, labelfont={bf,color=FpgaBlue}}
\usepackage{subcaption}

% --- Листинги кода (minted: нужен Pygments и ключ -shell-escape) ---
\usepackage[outputdir=../build]{minted}
\usemintedstyle{friendly}
\setminted{fontsize=\small, breaklines=true, tabsize=4, baselinestretch=1.0}

% --- Цветные блоки ---
\usepackage[most]{tcolorbox}
\tcbuselibrary{minted}
\tcbset{codestyle/.style={enhanced, breakable, listing only, listing engine=minted,
  colback=CodeBack, colframe=FpgaBlue, boxrule=0pt, leftrule=2.2pt, arc=0pt,
  left=6pt, right=4pt, top=2pt, bottom=2pt, before skip=6pt, after skip=8pt,
  fonttitle=\ttfamily\small\bfseries, coltitle=FpgaDark, colbacktitle=FpgaLight,
  attach boxed title to top left={yshift=-1pt}, boxed title style={boxrule=0pt, arc=0pt}}}
% Фрагмент кода: \begin{vcode} ... \end{vcode}, язык по умолчанию Verilog
\newtcblisting{vcode}[1][verilog]{codestyle, minted language=#1,
  minted options={fontsize=\small, breaklines=true, tabsize=4}}
% Файл целиком с номерами строк: \vfile{путь}{подпись}
\newtcbinputlisting{\vfile}[2]{codestyle, minted language=verilog,
  minted options={fontsize=\small, breaklines=true, tabsize=4, linenos, numbersep=5pt},
  listing file={#1}, title={#2}, left=20pt}
\newtcolorbox{infobox}[1][Справка]{
  enhanced, breakable, colback=FpgaLight, colframe=FpgaBlue,
  fonttitle=\sffamily\bfseries, title={#1}, arc=2mm, boxrule=0.8pt}
\newtcolorbox{warnbox}[1][Внимание!]{
  enhanced, breakable, colback=orange!6, colframe=FpgaOrange,
  fonttitle=\sffamily\bfseries, title={#1}, arc=2mm, boxrule=0.8pt}
\newtcolorbox{taskbox}[1][Задание]{
  enhanced, breakable, colback=green!4, colframe=FpgaGreen,
  fonttitle=\sffamily\bfseries, title={#1}, arc=2mm, boxrule=0.8pt}
\newtcolorbox{ideabox}[1][Главная мысль]{
  enhanced, breakable, colback=purple!4, colframe=FpgaPurple,
  fonttitle=\sffamily\bfseries, title={#1}, arc=2mm, boxrule=0.8pt}

% --- Прочее ---
\usepackage{siunitx}
\sisetup{output-decimal-marker={,}, group-separator={\,}, locale=FR}
\usepackage{hyperref}
\hypersetup{colorlinks=true, linkcolor=FpgaBlue, urlcolor=FpgaBlue,
  citecolor=FpgaBlue, pdftitle={ПЛИС на практике: Sipeed Tang Nano 20K},
  pdfauthor={Корягин С.А.}}

% Удобные команды
\newcommand{\code}[1]{\texttt{\small #1}}
\newcommand{\term}[1]{\textbf{\textcolor{FpgaDark}{#1}}}
\newcommand{\codepath}{../code}

% Общие стили TikZ
\tikzset{
  block/.style={draw=FpgaBlue, fill=FpgaLight, thick, rounded corners=2pt,
                minimum height=9mm, minimum width=22mm, align=center,
                font=\sffamily\small},
  hblock/.style={block, fill=FpgaBlue, text=white},
  oblock/.style={block, draw=FpgaOrange, fill=orange!10},
  gblock/.style={block, draw=FpgaGreen, fill=green!8},
  arr/.style={-{Stealth[length=2.5mm]}, thick, color=FpgaDark},
  wire/.style={thick, color=FpgaDark},
  bus/.style={line width=2pt, color=FpgaBlue},
  lbl/.style={font=\sffamily\scriptsize, color=FpgaDark},
}
\tikzset{
  gbase/.style={draw=FpgaDark, thick, fill=FpgaLight, logic gate inputs=nn,
                minimum width=12mm, minimum height=9mm, inner sep=0pt},
  gAND/.style={and gate US, gbase},
  gOR/.style={or gate US, gbase},
  gXOR/.style={xor gate US, gbase},
  gNAND/.style={nand gate US, gbase},
  gNOR/.style={nor gate US, gbase},
  gXNOR/.style={xnor gate US, gbase},
  ibase/.style={draw=FpgaDark, thick, fill=FpgaLight, logic gate inputs=nn,
                minimum width=8mm, minimum height=10mm, font=\sffamily\small},
  iAND/.style={and gate IEC, ibase},
  iOR/.style={or gate IEC, ibase},
  iXOR/.style={xor gate IEC, ibase},
  iNAND/.style={nand gate IEC, ibase},
  iNOR/.style={nor gate IEC, ibase},
  iXNOR/.style={xnor gate IEC, ibase},
  iNOT/.style={not gate IEC, draw=FpgaDark, thick, fill=FpgaLight,
               minimum width=8mm, minimum height=7mm, font=\sffamily\small},
  gNOT/.style={not gate US, draw=FpgaDark, thick, fill=FpgaLight,
               minimum width=9mm, minimum height=6mm, inner sep=0pt},
}


\begin{document}

% ------------------------------ Титульный лист ------------------------------
\begin{titlepage}
\begin{tikzpicture}[remember picture, overlay]
  \fill[FpgaDark] (current page.north west) rectangle ([yshift=-11cm]current page.north east);
  % декоративная «матрица» логических блоков
  \foreach \x in {0,...,15} \foreach \y in {0,...,5} {
    \pgfmathsetmacro{\op}{mod(\x*7+\y*13,10)/12+0.1}
    \fill[FpgaBlue, opacity=0.35*\op] ([xshift=\x*1.35cm+0.4cm, yshift=-\y*1.35cm-2.8cm]current page.north west)
      rectangle ++(0.95,0.95);
  }
  \foreach \x in {0,...,14} \foreach \y in {0,...,5}
    \draw[FpgaOrange!70, opacity=0.35] ([xshift=\x*1.35cm+1.35cm, yshift=-\y*1.35cm-2.33cm]current page.north west) -- ++(0.4,0);
  \node[anchor=north west, text=white, font=\sffamily\bfseries\fontsize{34}{40}\selectfont, align=left]
    at ([xshift=2cm, yshift=-3.4cm]current page.north west)
    {ПЛИС на практике};
  \node[anchor=north west, text=FpgaOrange, font=\sffamily\bfseries\LARGE, align=left]
    at ([xshift=2cm, yshift=-5.4cm]current page.north west)
    {Sipeed Tang Nano 20K};
  \node[anchor=north west, text=white, font=\sffamily\large, align=left]
    at ([xshift=2cm, yshift=-6.7cm]current page.north west)
    {Учебное пособие для начинающих\\разработчиков цифровой аппаратуры};
  \node[anchor=south west, text=white!80, font=\sffamily\small, align=left]
    at ([xshift=2cm, yshift=-10.4cm]current page.north west)
    {GW2AR-18 $\cdot$ 20\,736 LUT4 $\cdot$ 64 Мбит SDRAM $\cdot$ Verilog};
\end{tikzpicture}

\vspace*{11cm}
\begin{center}
\begin{tikzpicture}[scale=0.9]
  % Упрощённое изображение платы
  \fill[BoardGreen, rounded corners=3pt] (0,0) rectangle (9,3.4);
  \foreach \i in {0,...,19} {
    \fill[yellow!70!orange] (0.3+\i*0.42,0.12) circle (0.08);
    \fill[yellow!70!orange] (0.3+\i*0.42,3.28) circle (0.08);
  }
  \fill[gray!70] (-0.35,1.2) rectangle (0.6,2.2);        % USB-C
  \node[font=\tiny\sffamily, white] at (0.12,0.95) {USB-C};
  \fill[ChipBlack] (3.2,0.8) rectangle (5.2,2.8);         % ПЛИС
  \node[white, font=\tiny\sffamily, align=center] at (4.2,1.8) {GOWIN\\GW2AR-18};
  \fill[ChipBlack] (1.2,1.2) rectangle (2.3,2.3);         % BL616
  \node[white, font=\tiny\sffamily] at (1.75,1.75) {BL616};
  \fill[gray!50] (8.1,1.0) rectangle (9.3,2.4);            % HDMI
  \node[font=\tiny\sffamily] at (8.7,1.7) {HDMI};
  \foreach \i in {0,...,5} \fill[FpgaOrange] (5.8+\i*0.28,2.7) rectangle ++(0.18,0.12);
  \fill[gray!30] (6,0.6) rectangle (6.6,1.2); \fill[gray!30] (7,0.6) rectangle (7.6,1.2);
  \node[font=\tiny\sffamily, white] at (6.3,0.4) {S1};
  \node[font=\tiny\sffamily, white] at (7.3,0.4) {S2};
\end{tikzpicture}
\end{center}

\vfill
\begin{center}
{\sffamily\large Корягин Сергей Андреевич}\\[2mm]
{\sffamily\small асп. СПбГЭТУ «ЛЭТИ», педагог ЦО «Кудрово»}\\[8mm]
{\sffamily г. Кудрово, \the\year}
\end{center}
\end{titlepage}

\tableofcontents

\chapter*{Как читать это пособие}
\addcontentsline{toc}{chapter}{Как читать это пособие}

Вы уже знаете, что такое цифровой сигнал, умеете читать таблицы истинности
и собирать схемы из элементов И, ИЛИ, НЕ, И-НЕ, ИЛИ-НЕ. Это пособие
продолжает ту же историю: мы научимся собирать \emph{любые} цифровые
схемы не паяльником и не на макетной плате, а \emph{описывая} их текстом, после чего
специальная микросхема (ПЛИС) превращается в ту самую схему, которую
мы описали.

\begin{figure}[H]
\centering
\begin{tikzpicture}[node distance=6mm and 9mm]
  \node[block] (a) {Логические\\элементы};
  \node[block, right=of a] (b) {Триггеры\\и регистры};
  \node[block, right=of b] (c) {Что такое\\ПЛИС};
  \node[hblock, right=of c] (d) {Verilog\\и примеры};
  \node[oblock, below=of d] (e) {Свои\\проекты};
  \node[block, left=of e] (f) {Память\\SDRAM};
  \draw[arr] (a) -- (b); \draw[arr] (b) -- (c); \draw[arr] (c) -- (d);
  \draw[arr] (d) -- (e); \draw[arr] (d) -- (f); \draw[arr] (f) -- (e);
  \node[lbl, above=1mm of a] {вы уже здесь};
\end{tikzpicture}
\caption{Путь, который мы пройдём}
\end{figure}

Пособие построено так:
\begin{enumerate}
  \item \textbf{Главы 1--2} объясняют, что такое ПЛИС, зачем она нужна и что
        находится внутри кристалла.
  \item \textbf{Глава 3} знакомит с платой Sipeed Tang Nano 20K.
  \item \textbf{Глава 4} показывает весь путь от текста до работающей схемы.
  \item \textbf{Глава 5} учит основам языка Verilog.
  \item \textbf{Главы 6--7} содержат шесть законченных проектов и учат
        проверять схему в симуляторе.
  \item \textbf{Глава 8} рассказывает о встроенной памяти SDRAM на 64 Мбит.
  \item \textbf{Глава 9} состоит из заданий для самостоятельной работы.
\end{enumerate}

Все исходные тексты примеров лежат рядом с книгой, в папке
\nolinkurl{personal/fpga/tang_nano_20k/code}. Их можно открыть, скопировать и
сразу загрузить в плату.

\begin{infobox}[Обозначения]
\begin{itemize}
  \item Синие блоки «Справка» содержат полезные подробности.
  \item Оранжевые блоки «Внимание!» предупреждают о типичных ошибках.
  \item Фиолетовые блоки «Главная мысль» подводят итог раздела.
  \item Зелёные блоки «Задание» предлагают что-то сделать самостоятельно.
\end{itemize}
\end{infobox}


% ============================================================
% ЧАСТЬ I. ОСНОВЫ ЦИФРОВОЙ СХЕМОТЕХНИКИ
% ============================================================
\part{Основы цифровой схемотехники}
\chapter{Сигналы}
\label{ch:signals}

\section{Что такое сигнал}

\term{Сигнал} --- это физическая величина, которая меняется во времени и
несёт информацию. В электронике чаще всего это \emph{напряжение} на проводе.
Микрофон превращает звук в изменяющееся напряжение, датчик температуры ---
температуру в напряжение, кнопка --- нажатие в напряжение.

Сигналы бывают двух видов.

\begin{description}
  \item[Аналоговый сигнал] может принимать \emph{любое} значение в некотором
    диапазоне: 1,2~В, 1,21~В, 1,215~В\ldots\ Так выглядит звук, сигнал с
    микрофона или датчика.
  \item[Цифровой сигнал] принимает только несколько разрешённых значений. В
    двоичной цифровой технике их два: \textbf{логический 0} (низкое
    напряжение) и \textbf{логическая 1} (высокое напряжение).
\end{description}

\begin{figure}[H]
\centering
\begin{tikzpicture}
\begin{groupplot}[
  group style={group size=2 by 1, horizontal sep=18mm},
  width=0.5\textwidth, height=4.6cm, xmin=0, xmax=10,
  xlabel={Время}, axis lines*=left, xtick=\empty,
  label style={font=\sffamily\small}, tick label style={font=\sffamily\small},
  title style={font=\sffamily\small\bfseries},
]
\nextgroupplot[title={Аналоговый сигнал}, ylabel={$U$, В}, ymin=0, ymax=3.5]
  \addplot[very thick, FpgaBlue, domain=0:10, samples=200] {1.65+1.2*sin(deg(x*0.9))+0.35*sin(deg(x*2.7))};
\nextgroupplot[title={Цифровой сигнал}, ymin=-0.2, ymax=1.4, ytick={0,1}, yticklabels={0,1}, ylabel={логический уровень}]
  \addplot[very thick, FpgaOrange, const plot] coordinates {(0,0) (1,1) (2.5,0) (3,1) (5,1) (5.5,0) (7,1) (8,0) (9,1) (10,1)};
\end{groupplot}
\end{tikzpicture}
\caption{Аналоговый сигнал меняется плавно, цифровой --- скачками между двумя уровнями}
\end{figure}

\subsection{Почему цифровой}

Почему вся современная вычислительная техника цифровая? Главная причина ---
\textbf{устойчивость к помехам}. Любой провод ловит помехи: от соседних
проводов, от моторов, от радиоволн. В аналоговом сигнале помеха искажает
саму информацию, и отделить её от полезного сигнала уже нельзя. В цифровом
сигнале нас интересует только одно: \emph{выше или ниже порога} напряжение.
Небольшая помеха ничего не меняет.

\begin{figure}[H]
\centering
\begin{tikzpicture}
\begin{axis}[
  width=0.9\textwidth, height=5.4cm, xmin=0, xmax=10, ymin=-0.4, ymax=3.8,
  xlabel={Время}, ylabel={$U$, В}, xtick=\empty, ytick={0,0.8,2.0,3.3},
  yticklabels={0, 0{,}8, 2{,}0, 3{,}3},
  axis lines*=left, label style={font=\sffamily\small}, tick label style={font=\sffamily\small},
  clip=false,
]
  \fill[green!12] (axis cs:0,2.0) rectangle (axis cs:10,3.8);
  \fill[red!10] (axis cs:0,0.8) rectangle (axis cs:10,2.0);
  \fill[blue!8] (axis cs:0,-0.4) rectangle (axis cs:10,0.8);
  \addplot[thick, FpgaBlue, samples=400, domain=0:10]
    {(x>1 && x<4 || x>6 && x<8.5 ? 3.1 : 0.2) + 0.25*sin(deg(x*37)) + 0.15*sin(deg(x*91))};
  \node[font=\sffamily\small, anchor=west] at (axis cs:10.1,2.9) {«1»};
  \node[font=\sffamily\small, anchor=west, text=FpgaRed] at (axis cs:10.1,1.4) {запрещено};
  \node[font=\sffamily\small, anchor=west] at (axis cs:10.1,0.2) {«0»};
\end{axis}
\end{tikzpicture}
\caption{Сигнал с помехами. Пока напряжение не заходит в запрещённую зону,
приёмник безошибочно отличает 0 от 1 (пороги --- для стандарта LVCMOS 3,3~В)}
\label{fig:levels}
\end{figure}

\section{Логические уровни}

Какое напряжение считать нулём, а какое --- единицей, определяет
\term{стандарт логических уровней}. ПЛИС на плате Tang Nano 20K использует
стандарт \textbf{LVCMOS33} (питание 3,3~В):

\begin{table}[H]
\centering
\caption{Уровни LVCMOS 3,3 В}
\begin{tabular}{l l l}
\toprule
\textbf{Параметр} & \textbf{Значение} & \textbf{Смысл} \\
\midrule
$U_\text{IL,max}$ & 0,8 В & всё, что ниже, вход считает нулём \\
$U_\text{IH,min}$ & 2,0 В & всё, что выше, вход считает единицей \\
$U_\text{OL,max}$ & 0,4 В & выход выдаёт «0» не выше этого напряжения \\
$U_\text{OH,min}$ & 2,4 В & выход выдаёт «1» не ниже этого напряжения \\
\bottomrule
\end{tabular}
\end{table}

\begin{figure}[H]
\centering
\begin{tikzpicture}[yscale=1.3]
  \fill[blue!10] (0,0) rectangle (2.4,0.4);   \fill[green!15] (0,2.4) rectangle (2.4,3.3);
  \fill[gray!15] (0,0.4) rectangle (2.4,2.4);
  \fill[blue!10] (5,0) rectangle (7.4,0.8);   \fill[green!15] (5,2.0) rectangle (7.4,3.3);
  \fill[red!12] (5,0.8) rectangle (7.4,2.0);
  \draw[thick] (0,0) rectangle (2.4,3.3); \draw[thick] (5,0) rectangle (7.4,3.3);
  \node[font=\sffamily\small\bfseries] at (1.2,3.65) {Выход};
  \node[font=\sffamily\small\bfseries] at (6.2,3.65) {Вход};
  \node[font=\sffamily\small] at (1.2,2.85) {«1»}; \node[font=\sffamily\small] at (1.2,0.2) {«0»};
  \node[font=\sffamily\small] at (6.2,2.65) {«1»}; \node[font=\sffamily\small] at (6.2,0.4) {«0»};
  \node[font=\sffamily\scriptsize, text=FpgaRed] at (6.2,1.4) {неопределённость};
  \foreach \v/\l in {0.4/0{,}4, 2.4/2{,}4, 3.3/3{,}3} \node[font=\sffamily\scriptsize, anchor=east] at (0,\v) {\l~В};
  \foreach \v/\l in {0.8/0{,}8, 2.0/2{,}0} \node[font=\sffamily\scriptsize, anchor=west] at (7.4,\v) {\l~В};
  \draw[FpgaOrange, <->, thick] (3.7,2.0) -- (3.7,2.4) node[midway, right, font=\sffamily\scriptsize, align=left] {запас\\помехо-\\устойчивости};
  \draw[FpgaOrange, <->, thick] (3.7,0.4) -- (3.7,0.8);
  \draw[dashed, FpgaGray] (2.4,2.4) -- (5,2.4); \draw[dashed, FpgaGray] (2.4,2.0) -- (5,2.0);
  \draw[dashed, FpgaGray] (2.4,0.4) -- (5,0.4); \draw[dashed, FpgaGray] (2.4,0.8) -- (5,0.8);
\end{tikzpicture}
\caption{Выход гарантирует уровни «с запасом», поэтому помеха в 0,4 В не
приводит к ошибке}
\end{figure}

\begin{infobox}[Активный высокий и активный низкий уровень]
Обычно «действие» происходит при единице: 1 --- светодиод горит, 1 --- мотор
включён. Это \term{активный высокий уровень}. Но часто бывает наоборот: сигнал
сброса \emph{срабатывает при нуле}. Такие сигналы называют сигналами с
\term{активным низким уровнем} и помечают чертой сверху ($\overline{\text{RESET}}$),
решёткой (RESET\#) или окончанием \code{\_n} (\code{rst\_n}). На плате Tang
Nano 20K светодиоды именно такие --- они горят от нуля.
\end{infobox}

\section{Двоичная система счисления}

Один цифровой провод передаёт один \term{бит} --- 0 или 1. Чтобы передать
число больше единицы, берут несколько проводов (битов) и используют
\term{двоичную систему счисления}. Каждый разряд имеет вес --- степень
двойки:

\begin{figure}[H]
\centering
\begin{tikzpicture}
  \foreach \i/\b in {0/1,1/0,2/1,3/1} {
    \pgfmathtruncatemacro{\p}{3-\i}
    \pgfmathtruncatemacro{\w}{2^(3-\i)}
    \node[draw=FpgaBlue, fill=FpgaLight, thick, minimum size=11mm, font=\ttfamily\Large] (b\i) at (\i*1.6,0) {\b};
    \node[font=\sffamily\small, above=1mm of b\i] {$2^{\p}=\w$};
    \node[font=\sffamily\small, below=1mm of b\i] {$\b\cdot\w$};
  }
  \node[font=\sffamily\small, anchor=west] at (6.2,-0.85) {$= 8+0+2+1 = 11$};
  \node[font=\sffamily\small, anchor=west] at (6.2,0) {$1011_2 = 11_{10}$};
\end{tikzpicture}
\caption{Перевод двоичного числа в десятичное}
\end{figure}

С помощью $n$ битов можно записать $2^n$ разных чисел: от $0$ до $2^n-1$.
Четыре бита дают 16 комбинаций, восемь (\term{байт}) --- 256. Длинные
двоичные числа неудобно читать, поэтому их записывают в
\term{шестнадцатеричной} системе: каждые 4 бита заменяются одной цифрой от 0
до F.

\begin{table}[H]
\centering
\caption{Числа от 0 до 15 в трёх системах счисления}
\small
\begin{tabular}{c c c @{\hspace{12mm}} c c c}
\toprule
\textbf{Дес.} & \textbf{Двоич.} & \textbf{Шестн.} & \textbf{Дес.} & \textbf{Двоич.} & \textbf{Шестн.} \\
\midrule
0 & \code{0000} & 0 & 8  & \code{1000} & 8 \\
1 & \code{0001} & 1 & 9  & \code{1001} & 9 \\
2 & \code{0010} & 2 & 10 & \code{1010} & A \\
3 & \code{0011} & 3 & 11 & \code{1011} & B \\
4 & \code{0100} & 4 & 12 & \code{1100} & C \\
5 & \code{0101} & 5 & 13 & \code{1101} & D \\
6 & \code{0110} & 6 & 14 & \code{1110} & E \\
7 & \code{0111} & 7 & 15 & \code{1111} & F \\
\bottomrule
\end{tabular}
\end{table}

\begin{figure}[H]
\centering
\begin{tikzpicture}
\begin{axis}[
  width=0.75\textwidth, height=5cm, ybar, bar width=6mm,
  xlabel={Число битов $n$}, ylabel={Число комбинаций $2^n$},
  xtick={1,...,8}, ymin=0, ymax=300,
  nodes near coords, nodes near coords style={font=\sffamily\scriptsize},
  label style={font=\sffamily\small}, tick label style={font=\sffamily\small},
  axis lines*=left, grid=major, grid style={FpgaGray!20},
]
  \addplot[fill=FpgaBlue!70, draw=FpgaBlue] coordinates {(1,2) (2,4) (3,8) (4,16) (5,32) (6,64) (7,128) (8,256)};
\end{axis}
\end{tikzpicture}
\caption{Каждый новый бит удваивает число возможных значений}
\end{figure}

\section{Временные диаграммы}

Чтобы показать, как сигналы меняются во времени, рисуют
\term{временные диаграммы}: по горизонтали --- время, по вертикали --- уровень.
Переход из 0 в 1 называют \term{передним фронтом} (или просто фронтом),
переход из 1 в 0 --- \term{задним фронтом} (или спадом).

\begin{figure}[H]
\centering
\begin{tikzpicture}[x=1cm, y=1cm]
  \draw[very thick, FpgaBlue] (0,0) -- (2,0) -- (2.2,1.2) -- (5,1.2) -- (5.2,0) -- (8,0);
  \draw[-{Stealth}, FpgaGray] (0,-0.5) -- (8.5,-0.5) node[right, font=\sffamily\small] {$t$};
  \node[font=\sffamily\small, anchor=east] at (0,0) {0};
  \node[font=\sffamily\small, anchor=east] at (0,1.2) {1};
  \draw[FpgaOrange, thick, -{Stealth}] (1.3,1.9) node[above, font=\sffamily\small] {передний фронт ($0\to1$)} -- (2.05,0.8);
  \draw[FpgaOrange, thick, -{Stealth}] (6.3,1.9) node[above, font=\sffamily\small] {задний фронт, спад ($1\to0$)} -- (5.15,0.8);
  \draw[decorate, decoration={brace, amplitude=5pt, mirror}] (2.2,-0.1) -- (5,-0.1) node[midway, below=5pt, font=\sffamily\scriptsize] {длительность импульса};
\end{tikzpicture}
\caption{Импульс на временной диаграмме}
\end{figure}

Если на диаграмме несколько проводов объединены в число, их рисуют «шиной»:
две линии, между которыми написано значение, а крестик обозначает момент
смены значения.

\begin{figure}[H]
\centering
\begin{tikztimingtable}[timing/slope=0.1, xscale=2.2, yscale=1.4, timing/font=\sffamily\small, timing/rowdist=3.2ex]
  a        & 2L 4H 2L 2H \\
  b        & 4L 4H 2L \\
  data[3:0] & 2D{0} 2D{5} 2D{A} 2D{F} 2D{3} \\
  z (отключён) & 3Z 4H 3Z \\
\end{tikztimingtable}
\caption{Одиночные сигналы, шина и сигнал в третьем (отключённом) состоянии Z}
\end{figure}

\section{Тактовый сигнал}

Особую роль в цифровой технике играет \term{тактовый сигнал} (clock, \code{clk}) ---
последовательность импульсов, которые идут строго через равные промежутки
времени. Он работает как дирижёр: по его фронтам все элементы памяти схемы
одновременно обновляют свои значения.

\begin{figure}[H]
\centering
\begin{tikzpicture}[x=1cm, y=1cm]
  \foreach \i in {0,1,2,3} {
    \draw[very thick, FpgaBlue] (\i*3,0) -- (\i*3,1.2) -- (\i*3+1.5,1.2) -- (\i*3+1.5,0) -- (\i*3+3,0);
    \draw[FpgaOrange, -{Stealth}, thick] (\i*3,1.9) -- (\i*3,1.3);
  }
  \draw[-{Stealth}, FpgaGray] (-0.3,-0.4) -- (12.5,-0.4) node[right, font=\sffamily\small] {$t$};
  \draw[<->, thick] (3,-0.9) -- (6,-0.9) node[midway, below, font=\sffamily\small] {период $T$};
  \draw[<->, thick] (6,1.55) -- (7.5,1.55) node[midway, above, font=\sffamily\small] {$t_\text{и}$};
  \node[font=\sffamily\small, text=FpgaOrange, anchor=west] at (9.3,2.1) {передние фронты};
\end{tikzpicture}
\caption{Тактовый сигнал: период $T$, длительность импульса $t_\text{и}$}
\end{figure}

Основные характеристики тактового сигнала:
\[
  f = \frac{1}{T}, \qquad D = \frac{t_\text{и}}{T}\cdot 100\,\%.
\]
Здесь $f$ --- \term{частота} (сколько периодов в секунду, измеряется в герцах),
$D$ --- \term{коэффициент заполнения} (обычно 50\,\%). На плате Tang Nano 20K
стоит генератор на 27~МГц:
\[
  T = \frac{1}{27\,000\,000\ \text{Гц}} \approx 37 \cdot 10^{-9}\ \text{с} = 37\ \text{нс}.
\]
За одну секунду происходит 27 миллионов тактов. За время одного моргания глаза
(около 0,1~с) --- почти три миллиона.

\begin{table}[H]
\centering
\caption{Приставки для частоты и времени}
\begin{tabular}{l l l l}
\toprule
\textbf{Частота} & \textbf{Период} & \textbf{Пример} \\
\midrule
1 Гц  & 1 с   & мигание светодиода \\
1 кГц = $10^3$ Гц & 1 мс = $10^{-3}$ с & звуковой сигнал «писк» \\
1 МГц = $10^6$ Гц & 1 мкс = $10^{-6}$ с & простые микроконтроллеры \\
27 МГц & 37 нс & генератор Tang Nano 20K \\
1 ГГц = $10^9$ Гц & 1 нс = $10^{-9}$ с & процессоры компьютеров \\
\bottomrule
\end{tabular}
\end{table}

\begin{ideabox}
Цифровой сигнал имеет два уровня: 0 и 1. Несколько битов образуют двоичное
число. Изменения сигналов во времени изображают временными диаграммами, а
тактовый сигнал задаёт ритм, в котором работает вся схема.
\end{ideabox}

\chapter{Логические элементы и булева алгебра}
\label{ch:gates}

\section{Логические функции}

\term{Логическая функция} ставит в соответствие каждой комбинации входных
битов значение выхода: 0 или 1. Функцию удобно задавать \term{таблицей
истинности} --- перечнем всех комбинаций входов и значений выхода. У функции от
$n$ входов в таблице $2^n$ строк.

Схема, которая вычисляет логическую функцию, называется
\term{логическим элементом} (вентилем). Логические элементы изображают
условными знаками. В нашей стране принят стандарт ГОСТ (прямоугольник с
символом внутри), в зарубежной литературе и в программах часто используют
американские (ANSI) знаки. Полезно узнавать оба.

\section{Основные элементы}

\newcommand{\gaterow}[5]{%
  \begin{tikzpicture}[baseline=-1mm]
    \node[#1, minimum width=10mm, minimum height=7mm] (g) at (0,0) {};
    \ifstrequal{#1}{gNOT}{\draw[wire] (g.input) -- ++(-0.35,0);}{%
      \draw[wire] (g.input 1) -- ++(-0.35,0); \draw[wire] (g.input 2) -- ++(-0.35,0);}
    \draw[wire] (g.output) -- ++(0.35,0);
  \end{tikzpicture} &
  \begin{tikzpicture}[baseline=-1mm]
    \node[#2] (h) at (0,0) {};
    \ifstrequal{#2}{iNOT}{\draw[wire] (h.input) -- ++(-0.35,0);}{%
      \draw[wire] (h.input 1) -- ++(-0.35,0); \draw[wire] (h.input 2) -- ++(-0.35,0);}
    \draw[wire] (h.output) -- ++(0.35,0);
  \end{tikzpicture} & #3 & #4 & #5 \\[3mm]}

\begin{table}[H]
\centering
\caption{Основные логические элементы. Столбец «таблица» --- значения выхода
для входов $ab = 00, 01, 10, 11$}
\label{tab:gates_all}
\small
\begin{tabular}{c c l l c}
\toprule
\textbf{ANSI} & \textbf{ГОСТ / IEC} & \textbf{Название} & \textbf{Формула} & \textbf{Таблица} \\
\midrule
\gaterow{gNOT}{iNOT}{НЕ (инверсия)}{$y=\overline{a}$}{\code{10} (для $a=0,1$)}
\gaterow{gAND}{iAND}{И (конъюнкция)}{$y=a\cdot b$}{\code{0001}}
\gaterow{gOR}{iOR}{ИЛИ (дизъюнкция)}{$y=a+b$}{\code{0111}}
\gaterow{gNAND}{iNAND}{И-НЕ (штрих Шеффера)}{$y=\overline{a\cdot b}$}{\code{1110}}
\gaterow{gNOR}{iNOR}{ИЛИ-НЕ (стрелка Пирса)}{$y=\overline{a+b}$}{\code{1000}}
\gaterow{gXOR}{iXOR}{Исключающее ИЛИ}{$y=a\oplus b$}{\code{0110}}
\gaterow{gXNOR}{iXNOR}{Исключающее ИЛИ-НЕ}{$y=\overline{a\oplus b}$}{\code{1001}}
\bottomrule
\end{tabular}
\end{table}

Кружок на выходе элемента обозначает инверсию. Поэтому И-НЕ рисуют как
элемент И с кружком, ИЛИ-НЕ --- как ИЛИ с кружком.

Словами:
\begin{itemize}
  \item \textbf{НЕ}: выход противоположен входу;
  \item \textbf{И}: выход равен 1, только если \emph{все} входы равны 1;
  \item \textbf{ИЛИ}: выход равен 1, если \emph{хотя бы один} вход равен 1;
  \item \textbf{Исключающее ИЛИ}: выход равен 1, если входы \emph{различны}
        (для многих входов --- если число единиц нечётно).
\end{itemize}

\begin{figure}[H]
\centering
\begin{tikztimingtable}[timing/slope=0.05, xscale=2.4, yscale=1.4, timing/font=\sffamily\small, timing/rowdist=3.2ex]
  $a$              & 2L 2L 2H 2H \\
  $b$              & 2L 2H 2L 2H \\
  $a\cdot b$       & 6L 2H \\
  $a+b$            & 2L 6H \\
  $\overline{a\cdot b}$ & 6H 2L \\
  $a\oplus b$      & 2L 4H 2L \\
\end{tikztimingtable}
\caption{Работа элементов на временной диаграмме: перебираются все четыре
комбинации входов}
\end{figure}

\section{Как устроен элемент внутри}

Внутри микросхем логические элементы строятся из транзисторов. В
современной технологии КМОП (CMOS) используются пары транзисторов двух типов:
$p$-канальный открыт при нуле на затворе, $n$-канальный --- при единице.
Самый простой элемент --- инвертор --- состоит всего из двух транзисторов.

\begin{figure}[H]
\centering
\begin{circuitikz}[scale=0.95, transform shape]
  \draw (0,0) node[pmos, anchor=D] (P) {};
  \draw (P.S) node[vcc] {$+3{,}3$ В};
  \draw (P.D) -- ++(0,-0.3) coordinate (mid) -- ++(0,-0.3) node[nmos, anchor=D] (N) {};
  \draw (N.S) node[ground] {};
  \draw (P.G) -- ++(-0.5,0) coordinate (g1);
  \draw (N.G) -- ++(-0.5,0) coordinate (g2);
  \draw (g1) -- (g2);
  \draw ($(g1)!0.5!(g2)$) to[short, -o] ++(-1,0) node[left] {$a$};
  \draw (mid) to[short, *-o] ++(1.6,0) node[right] {$y=\overline{a}$};
  \node[font=\sffamily\small, align=left, anchor=west] at (4,0.2) {$a=0$: верхний открыт, $y$ соединён с $+3{,}3$ В $\to$ 1};
  \node[font=\sffamily\small, align=left, anchor=west] at (4,-1.9) {$a=1$: нижний открыт, $y$ соединён с землёй $\to$ 0};
\end{circuitikz}
\caption{КМОП-инвертор: два транзистора работают как переключатели}
\end{figure}

Элементы И-НЕ и ИЛИ-НЕ получаются из четырёх транзисторов, а И и ИЛИ ---
это И-НЕ и ИЛИ-НЕ с инвертором на выходе (шесть транзисторов). Поэтому в
микросхемах самыми «дешёвыми» оказываются именно инвертирующие элементы.

\section{Булева алгебра}

С логическими выражениями можно обращаться почти как с обычными
алгебраическими: раскрывать скобки, выносить общий множитель, сокращать.
Этими правилами занимается \term{булева алгебра}, названная в честь
английского математика Джорджа Буля. Операцию И записывают как умножение
($a\cdot b$ или $ab$), ИЛИ --- как сложение ($a+b$), НЕ --- чертой сверху
($\overline{a}$).

\begin{table}[H]
\centering
\caption{Основные законы булевой алгебры}
\label{tab:bool}
\begin{tabular}{l l l}
\toprule
\textbf{Закон} & \textbf{Для И} & \textbf{Для ИЛИ} \\
\midrule
Действия с константами & $a\cdot 0 = 0$, \ $a\cdot 1 = a$ & $a+1 = 1$, \ $a+0 = a$ \\
Повторение             & $a\cdot a = a$ & $a+a = a$ \\
Дополнение             & $a\cdot\overline{a} = 0$ & $a+\overline{a} = 1$ \\
Двойное отрицание      & \multicolumn{2}{c}{$\overline{\overline{a}} = a$} \\
Переместительный       & $ab = ba$ & $a+b = b+a$ \\
Сочетательный          & $(ab)c = a(bc)$ & $(a+b)+c = a+(b+c)$ \\
Распределительный      & $a(b+c) = ab+ac$ & $a+bc = (a+b)(a+c)$ \\
Поглощение             & $a(a+b) = a$ & $a+ab = a$ \\
\textbf{Законы де Моргана} & $\overline{a\cdot b} = \overline{a}+\overline{b}$ & $\overline{a+b} = \overline{a}\cdot\overline{b}$ \\
\bottomrule
\end{tabular}
\end{table}

\begin{warnbox}[Внимание: необычный закон]
Второй распределительный закон $a+bc = (a+b)(a+c)$ не выполняется в обычной
арифметике, но верен в булевой алгебре. Проверьте его таблицей истинности!
\end{warnbox}

\subsection{Законы де Моргана}

Законы де Моргана --- самые полезные на практике. Они говорят: \emph{инверсию
над выражением можно «протащить» внутрь, заменив И на ИЛИ и наоборот}.
Графически это значит, что элемент И-НЕ можно нарисовать как ИЛИ с
инверсиями на входах:

\begin{figure}[H]
\centering
\begin{tikzpicture}
  \node[gNAND] (g1) at (0,0) {};
  \draw[wire] (g1.input 1) -- ++(-0.6,0) node[left] {$a$};
  \draw[wire] (g1.input 2) -- ++(-0.6,0) node[left] {$b$};
  \draw[wire] (g1.output) -- ++(0.6,0) node[right] {$\overline{ab}$};
  \node[font=\Large] at (2.6,0) {$=$};
  \begin{scope}[xshift=5.2cm]
    \node[gOR] (g2) at (0,0) {};
    \node[gNOT, minimum width=6mm, minimum height=4mm] (n1) at (-1.3,0.18) {};
    \node[gNOT, minimum width=6mm, minimum height=4mm] (n2) at (-1.3,-0.18) {};
    \draw[wire] (n1.output) -- (g2.input 1);
    \draw[wire] (n2.output) -- (g2.input 2);
    \draw[wire] (n1.input) -- ++(-0.4,0) node[left] {$a$};
    \draw[wire] (n2.input) -- ++(-0.4,0) node[left] {$b$};
    \draw[wire] (g2.output) -- ++(0.6,0) node[right] {$\overline{a}+\overline{b}$};
  \end{scope}
\end{tikzpicture}
\caption{Закон де Моргана в виде схемы}
\end{figure}

\begin{table}[H]
\centering
\caption{Проверка закона де Моргана таблицей истинности}
\begin{tabular}{cc|c|cc|c}
\toprule
$a$ & $b$ & $\overline{a\cdot b}$ & $\overline{a}$ & $\overline{b}$ & $\overline{a}+\overline{b}$ \\
\midrule
0 & 0 & \textbf{1} & 1 & 1 & \textbf{1} \\
0 & 1 & \textbf{1} & 1 & 0 & \textbf{1} \\
1 & 0 & \textbf{1} & 0 & 1 & \textbf{1} \\
1 & 1 & \textbf{0} & 0 & 0 & \textbf{0} \\
\bottomrule
\end{tabular}
\end{table}

\section{Универсальные элементы}

Удивительный факт: \textbf{из одних только элементов И-НЕ можно собрать любую
логическую схему}. То же верно для ИЛИ-НЕ. Такие элементы называют
\term{универсальными}. Достаточно показать, как из И-НЕ получить НЕ, И и ИЛИ ---
а из них уже строится всё остальное.

\begin{figure}[H]
\centering
\begin{tikzpicture}
  % НЕ
  \node[font=\sffamily\small\bfseries] at (0.3,1.3) {НЕ};
  \node[gNAND] (a1) at (0.3,0) {};
  \draw[wire] (a1.input 1) -- ++(-0.35,0) coordinate (t1);
  \draw[wire] (a1.input 2) -- ++(-0.35,0) coordinate (t2);
  \draw[wire] (t1) -- (t2); \draw[wire] ($(t1)!0.5!(t2)$) -- ++(-0.5,0) node[left] {$a$};
  \fill ($(t1)!0.5!(t2)$) circle (1.3pt);
  \draw[wire] (a1.output) -- ++(0.4,0) node[right] {$\overline{a}$};
  % И
  \begin{scope}[xshift=4.2cm]
  \node[font=\sffamily\small\bfseries] at (0.9,1.3) {И};
  \node[gNAND] (b1) at (0,0) {};
  \node[gNAND] (b2) at (2.0,0) {};
  \draw[wire] (b1.input 1) -- ++(-0.4,0) node[left] {$a$};
  \draw[wire] (b1.input 2) -- ++(-0.4,0) node[left] {$b$};
  \draw[wire] (b1.output) -- ++(0.2,0) coordinate (m);
  \draw[wire] (m) |- (b2.input 1); \draw[wire] (m) |- (b2.input 2); \fill (m) circle (1.3pt);
  \draw[wire] (b2.output) -- ++(0.4,0) node[right] {$ab$};
  \end{scope}
  % ИЛИ
  \begin{scope}[xshift=10cm]
  \node[font=\sffamily\small\bfseries] at (0.9,1.3) {ИЛИ};
  \node[gNAND, minimum width=9mm, minimum height=6mm] (c1) at (0,0.45) {};
  \node[gNAND, minimum width=9mm, minimum height=6mm] (c2) at (0,-0.45) {};
  \node[gNAND] (c3) at (2.0,0) {};
  \foreach \c/\n in {c1/a, c2/b} {
    \draw[wire] (\c.input 1) -- ++(-0.3,0) coordinate (u1);
    \draw[wire] (\c.input 2) -- ++(-0.3,0) coordinate (u2);
    \draw[wire] (u1) -- (u2); \draw[wire] ($(u1)!0.5!(u2)$) -- ++(-0.35,0) node[left] {$\n$};
  }
  \draw[wire] (c1.output) -- ++(0.3,0) |- (c3.input 1);
  \draw[wire] (c2.output) -- ++(0.3,0) |- (c3.input 2);
  \draw[wire] (c3.output) -- ++(0.4,0) node[right] {$a+b$};
  \end{scope}
\end{tikzpicture}
\caption{НЕ, И и ИЛИ, собранные только из элементов И-НЕ}
\end{figure}

Проверим схему ИЛИ по законам де Моргана:
$\overline{\overline{a}\cdot\overline{b}} = \overline{\overline{a}} + \overline{\overline{b}} = a + b$.

\begin{infobox}[Зачем это нужно]
В эпоху микросхем серии К155 (ТТЛ) разработчики часто собирали целые
устройства из микросхем К155ЛА3, в каждой из которых четыре элемента И-НЕ.
Так схема получалась дешевле и требовала меньше разных деталей. Внутри
ПЛИС вопрос решён иначе --- там используются таблицы истинности (глава~\ref{ch:inside}).
\end{infobox}

\section{Многовходовые элементы}

Элементы И, ИЛИ и Исключающее ИЛИ могут иметь больше двух входов. Их можно
собрать из двухвходовых, соединив «цепочкой» или «деревом»:

\begin{figure}[H]
\centering
\begin{tikzpicture}
  \node[font=\sffamily\small\bfseries] at (1.2,1.6) {Цепочка: 3 уровня};
  \node[gAND, minimum width=9mm, minimum height=6mm] (a1) at (0,0.6) {};
  \node[gAND, minimum width=9mm, minimum height=6mm] (a2) at (1.3,0.3) {};
  \node[gAND, minimum width=9mm, minimum height=6mm] (a3) at (2.6,0) {};
  \draw[wire] (a1.input 1) -- ++(-0.4,0) node[left] {$a$};
  \draw[wire] (a1.input 2) -- ++(-0.4,0) node[left] {$b$};
  \draw[wire] (a1.output) -- ++(0.1,0) |- (a2.input 1);
  \draw[wire] (a2.input 2) -- ++(-1.7,0) node[left] {$c$};
  \draw[wire] (a2.output) -- ++(0.1,0) |- (a3.input 1);
  \draw[wire] (a3.input 2) -- ++(-3.0,0) node[left] {$d$};
  \draw[wire] (a3.output) -- ++(0.4,0) node[right] {$abcd$};
  \begin{scope}[xshift=7.5cm]
  \node[font=\sffamily\small\bfseries] at (0.7,1.6) {Дерево: 2 уровня};
  \node[gAND, minimum width=9mm, minimum height=6mm] (b1) at (0,0.5) {};
  \node[gAND, minimum width=9mm, minimum height=6mm] (b2) at (0,-0.5) {};
  \node[gAND, minimum width=9mm, minimum height=6mm] (b3) at (1.5,0) {};
  \draw[wire] (b1.input 1) -- ++(-0.4,0) node[left] {$a$};
  \draw[wire] (b1.input 2) -- ++(-0.4,0) node[left] {$b$};
  \draw[wire] (b2.input 1) -- ++(-0.4,0) node[left] {$c$};
  \draw[wire] (b2.input 2) -- ++(-0.4,0) node[left] {$d$};
  \draw[wire] (b1.output) -- ++(0.1,0) |- (b3.input 1);
  \draw[wire] (b2.output) -- ++(0.1,0) |- (b3.input 2);
  \draw[wire] (b3.output) -- ++(0.4,0) node[right] {$abcd$};
  \end{scope}
\end{tikzpicture}
\caption{Четырёхвходовый элемент И из двухвходовых. Дерево работает быстрее:
сигнал проходит через два элемента, а не через три}
\end{figure}

\begin{ideabox}
Семь основных логических элементов (НЕ, И, ИЛИ, И-НЕ, ИЛИ-НЕ, Исключающее
ИЛИ и Исключающее ИЛИ-НЕ) задаются таблицами истинности. Выражения с ними упрощаются по
законам булевой алгебры, а законы де Моргана позволяют менять И на ИЛИ.
Элементы И-НЕ и ИЛИ-НЕ универсальны: из любого из них собирается любая схема.
\end{ideabox}

\chapter{Комбинационные схемы}
\label{ch:comb}

\section{Комбинационная и последовательностная логика}

Все цифровые схемы делятся на два больших класса.

\begin{description}
  \item[Комбинационные схемы] --- выходы зависят \emph{только от текущих
    значений входов}. У них нет памяти: подайте те же входы --- получите те же
    выходы, неважно, что было раньше. Примеры: логические элементы,
    дешифраторы, мультиплексоры, сумматоры.
  \item[Последовательностные схемы] --- выходы зависят не только от текущих
    входов, но и от \emph{предыстории}, то есть от того, что было раньше.
    У них есть память. Примеры: триггеры, регистры, счётчики. О них ---
    в главах~\ref{ch:ff}--\ref{ch:sync}.
\end{description}

\begin{figure}[H]
\centering
\begin{tikzpicture}
  \node[block, minimum width=30mm, minimum height=20mm] (c) {Комбинационная\\логика};
  \foreach \i in {1,2,3} \draw[arr] ($(c.west)+(-1.2,0.9-\i*0.45)$) -- ($(c.west)+(0,0.9-\i*0.45)$);
  \foreach \i in {1,2} \draw[arr] ($(c.east)+(0,0.68-\i*0.45)$) -- ($(c.east)+(1.2,0.68-\i*0.45)$);
  \node[lbl, left=12mm of c] {входы};
  \node[lbl, right=12mm of c] {выходы};
  \node[font=\sffamily\small\bfseries, above=2mm of c] {Комбинационная схема};
  \begin{scope}[xshift=8.3cm]
    \node[block, minimum width=30mm, minimum height=20mm] (s) {Комбинационная\\логика};
    \node[gblock, minimum width=30mm, below=6mm of s] (m) {Память\\(триггеры)};
    \foreach \i in {1,2,3} \draw[arr] ($(s.west)+(-1.2,0.9-\i*0.45)$) -- ($(s.west)+(0,0.9-\i*0.45)$);
    \foreach \i in {1,2} \draw[arr] ($(s.east)+(0,0.68-\i*0.45)$) -- ($(s.east)+(1.2,0.68-\i*0.45)$);
    \draw[arr] ($(s.east)+(0.4,-0.7)$) -- ++(0,-1.0) |- (m.east);
    \draw[arr] (m.west) -- ++(-0.4,0) |- ($(s.west)+(0,-0.7)$);
    \node[lbl, left=4mm of m, align=center] {состояние};
    \node[font=\sffamily\small\bfseries, above=2mm of s] {Последовательностная схема};
  \end{scope}
\end{tikzpicture}
\caption{Главное отличие: у последовательностной схемы есть память и обратная
связь через неё}
\label{fig:comb_vs_seq}
\end{figure}

\section{От словесного условия к схеме}

Разработка комбинационной схемы всегда идёт по одному плану:
\begin{enumerate}
  \item понять, какие входы и выходы нужны, и дать им имена;
  \item составить \textbf{таблицу истинности};
  \item записать логическое выражение;
  \item \textbf{упростить} его;
  \item нарисовать схему.
\end{enumerate}

\subsection{Пример: голосование трёх судей}

На соревнованиях работают три судьи: $a$, $b$, $c$. Каждый нажимает кнопку,
если засчитывает попытку. Лампа $y$ должна загораться, если попытку
засчитали \emph{хотя бы двое}. Такая функция называется
\term{мажоритарной} (от лат. \emph{major} --- большинство).

\textbf{Шаг 1--2. Таблица истинности.} Три входа --- $2^3 = 8$ строк:

\begin{table}[H]
\centering
\caption{Таблица истинности мажоритарной функции}
\label{tab:maj}
\begin{tabular}{c|ccc|c|l}
\toprule
\textbf{№} & $a$ & $b$ & $c$ & $y$ & \textbf{Слагаемое} \\
\midrule
0 & 0 & 0 & 0 & 0 & \\
1 & 0 & 0 & 1 & 0 & \\
2 & 0 & 1 & 0 & 0 & \\
3 & 0 & 1 & 1 & \textbf{1} & $\overline{a}bc$ \\
4 & 1 & 0 & 0 & 0 & \\
5 & 1 & 0 & 1 & \textbf{1} & $a\overline{b}c$ \\
6 & 1 & 1 & 0 & \textbf{1} & $ab\overline{c}$ \\
7 & 1 & 1 & 1 & \textbf{1} & $abc$ \\
\bottomrule
\end{tabular}
\end{table}

\textbf{Шаг 3. Выражение.} Для каждой строки, где $y=1$, запишем
произведение всех входов: вход берётся как есть, если он равен 1, и с
инверсией, если равен 0. Такое произведение равно единице \emph{только} в
этой строке. Сложив все такие произведения, получим функцию. Эта форма
записи называется \term{совершенной дизъюнктивной нормальной формой} (СДНФ):
\[
  y = \overline{a}bc + a\overline{b}c + ab\overline{c} + abc.
\]

\textbf{Шаг 4. Упрощение.} Слагаемое $abc$ можно повторить сколько угодно
раз ($x + x = x$). Добавим его ещё дважды и сгруппируем:
\[
  y = (\overline{a}bc + abc) + (a\overline{b}c + abc) + (ab\overline{c} + abc)
    = bc(\overline{a}+a) + ac(\overline{b}+b) + ab(\overline{c}+c)
    = bc + ac + ab.
\]
Вместо четырёх трёхвходовых элементов И и трёх инверторов нужно всего три
двухвходовых элемента И и один элемент ИЛИ.

\textbf{Шаг 5. Схема.}

\begin{figure}[H]
\centering
\begin{tikzpicture}
  \foreach \n/\x in {a/0, b/0.6, c/1.2} {
    \draw[very thick, FpgaBlue] (\x,0.6) node[above, text=FpgaDark] {$\n$} -- (\x,-3.4);
  }
  \node[gAND] (g1) at (3.4,0) {};
  \node[gAND] (g2) at (3.4,-1.4) {};
  \node[gAND] (g3) at (3.4,-2.8) {};
  \node[gOR, logic gate inputs=nnn, minimum height=14mm] (o) at (6,-1.4) {};
  \draw[wire] (g1.input 1) -- (0,0 |- g1.input 1) node[circle, fill, inner sep=1.2pt] {};
  \draw[wire] (g1.input 2) -- (0.6,0 |- g1.input 2) node[circle, fill, inner sep=1.2pt] {};
  \draw[wire] (g2.input 1) -- (0,0 |- g2.input 1) node[circle, fill, inner sep=1.2pt] {};
  \draw[wire] (g2.input 2) -- (1.2,0 |- g2.input 2) node[circle, fill, inner sep=1.2pt] {};
  \draw[wire] (g3.input 1) -- (0.6,0 |- g3.input 1) node[circle, fill, inner sep=1.2pt] {};
  \draw[wire] (g3.input 2) -- (1.2,0 |- g3.input 2) node[circle, fill, inner sep=1.2pt] {};
  \draw[wire] (g1.output) -- node[above, lbl] {$ab$} ++(0.6,0) |- (o.input 1);
  \draw[wire] (g2.output) -- node[above, lbl] {$ac$} (o.input 2);
  \draw[wire] (g3.output) -- node[below, lbl] {$bc$} ++(0.6,0) |- (o.input 3);
  \draw[wire] (o.output) -- ++(0.8,0) node[right] {$y$};
\end{tikzpicture}
\caption{Мажоритарный элемент «2 из 3»: $y = ab + ac + bc$}
\label{fig:maj}
\end{figure}

\section{Карты Карно}

Упрощать выражения алгебраически можно, но легко ошибиться. Для функций от
2--4 переменных есть наглядный способ --- \term{карта Карно}. Это та же
таблица истинности, но нарисованная в виде прямоугольника так, что
\emph{соседние клетки отличаются значением только одной переменной}.
Поэтому столбцы подписаны в порядке 00, 01, 11, 10 (код Грея), а не
00, 01, 10, 11.

Правило упрощения: единицы объединяют в прямоугольные группы по 2, 4, 8
клеток (степень двойки), стараясь сделать группы как можно больше. Каждой
группе соответствует одно произведение, в котором остаются только те
переменные, что \emph{не меняются} внутри группы.

\begin{figure}[H]
\centering
\begin{tikzpicture}[cell/.style={draw, minimum size=11mm, font=\sffamily\large}]
  \foreach \c/\lab in {0/00, 1/01, 2/11, 3/10} \node[font=\sffamily\small] at (\c*1.1,1.05) {\lab};
  \foreach \r/\lab in {0/0, 1/1} \node[font=\sffamily\small] at (-0.95,-\r*1.1) {\lab};
  \node[font=\sffamily\small] at (-1.1,1.05) {$a \backslash bc$};
  % значения: строка a=0: bc=00,01,11,10 -> 0,0,1,0; a=1: 0,1,1,1
  \foreach \r/\c/\v in {0/0/0, 0/1/0, 0/2/1, 0/3/0, 1/0/0, 1/1/1, 1/2/1, 1/3/1}
    {\ifnum\v=1 \def\cc{FpgaLight}\else\def\cc{white}\fi
     \node[cell, fill=\cc] at (\c*1.1,-\r*1.1) {\v};}
  \draw[very thick, FpgaRed, rounded corners=4pt] (1.7,-1.62) rectangle (2.7,0.52);
  \draw[very thick, FpgaGreen, rounded corners=4pt] (0.6,-1.56) rectangle (2.62,-0.64);
  \draw[very thick, FpgaOrange, rounded corners=4pt] (1.78,-1.48) rectangle (3.8,-0.72);
  \node[font=\sffamily\small, text=FpgaRed, anchor=west] at (4.6,0.2) {$bc$ --- столбец $bc=11$ ($a$ меняется)};
  \node[font=\sffamily\small, text=FpgaGreen, anchor=west] at (4.6,-0.5) {$ac$ --- клетки $abc=101, 111$};
  \node[font=\sffamily\small, text=FpgaOrange, anchor=west] at (4.6,-1.2) {$ab$ --- клетки $abc=111, 110$};
  \node[font=\sffamily\small\bfseries, anchor=west] at (4.6,-2.0) {$y = ab + ac + bc$};
\end{tikzpicture}
\caption{Карта Карно мажоритарной функции: три группы по две клетки дают
три слагаемых}
\label{fig:karnaugh}
\end{figure}

\begin{infobox}[Кто будет упрощать схемы в ПЛИС?]
Когда вы будете работать с ПЛИС, упрощение выполняет программа-синтезатор,
причём гораздо лучше человека. Но понимать, \emph{как} это делается, всё
равно полезно: это помогает оценить, сколько ресурсов займёт ваша схема и
насколько быстро она будет работать.
\end{infobox}

\section{Задержка распространения}

До сих пор мы считали, что элемент срабатывает мгновенно. На самом деле
транзисторам нужно время, чтобы переключиться, а проводам --- чтобы
зарядиться. Выход элемента меняется через время $t_\text{зд}$ после входа.
Это \term{задержка распространения}. У современных микросхем она составляет
от долей наносекунды до нескольких наносекунд.

\begin{figure}[H]
\centering
\begin{tikzpicture}[x=1cm, y=1cm]
  \draw[very thick, FpgaBlue] (0,1.6) -- (2,1.6) -- (2.1,2.4) -- (6,2.4) -- (6.1,1.6) -- (9,1.6);
  \draw[very thick, FpgaOrange] (0,0.8) -- (3,0.8) -- (3.1,0) -- (7,0) -- (7.1,0.8) -- (9,0.8);
  \node[font=\sffamily\small, anchor=east] at (-0.1,2.0) {$a$};
  \node[font=\sffamily\small, anchor=east] at (-0.1,0.4) {$\overline{a}$};
  \draw[dashed, FpgaGray] (2.05,2.6) -- (2.05,-0.4); \draw[dashed, FpgaGray] (3.05,2.6) -- (3.05,-0.4);
  \draw[<->, thick] (2.05,-0.3) -- (3.05,-0.3) node[midway, below, font=\sffamily\small] {$t_\text{зд}$};
  \draw[dashed, FpgaGray] (6.05,2.6) -- (6.05,-0.4); \draw[dashed, FpgaGray] (7.05,2.6) -- (7.05,-0.4);
  \draw[<->, thick] (6.05,-0.3) -- (7.05,-0.3) node[midway, below, font=\sffamily\small] {$t_\text{зд}$};
\end{tikzpicture}
\caption{Выход инвертора меняется с задержкой $t_\text{зд}$}
\end{figure}

Если сигнал проходит через несколько элементов подряд, задержки
складываются. Самый длинный путь от входа до выхода называется
\term{критическим путём}. Он определяет, насколько быстро может работать
схема.

\begin{figure}[H]
\centering
\begin{tikzpicture}
  \node[gAND, minimum width=10mm, minimum height=7mm] (g1) at (0,0) {};
  \node[gOR, minimum width=10mm, minimum height=7mm] (g2) at (1.8,-0.3) {};
  \node[gXOR, minimum width=10mm, minimum height=7mm] (g3) at (3.6,-0.6) {};
  \node[gNOT] (g4) at (5.3,-0.6) {};
  \draw[wire, FpgaRed, very thick] (g1.input 1) -- ++(-0.5,0) node[left, text=FpgaDark] {$a$};
  \draw[wire] (g1.input 2) -- ++(-0.5,0) node[left] {$b$};
  \draw[wire, FpgaRed, very thick] (g1.output) -- ++(0.2,0) |- (g2.input 1);
  \draw[wire] (g2.input 2) -- ++(-2.3,0) node[left] {$c$};
  \draw[wire, FpgaRed, very thick] (g2.output) -- ++(0.2,0) |- (g3.input 1);
  \draw[wire] (g3.input 2) -- ++(-4.1,0) node[left] {$d$};
  \draw[wire, FpgaRed, very thick] (g3.output) -- (g4.input);
  \draw[wire, FpgaRed, very thick] (g4.output) -- ++(0.6,0) node[right, text=FpgaDark] {$y$};
  \node[font=\sffamily\small, text=FpgaRed, anchor=west] at (7.2,-0.2) {критический путь:};
  \node[font=\sffamily\small, text=FpgaRed, anchor=west] at (7.2,-0.8) {$t = 4\,t_\text{зд}$};
\end{tikzpicture}
\caption{Критический путь проходит через четыре элемента}
\end{figure}

\subsection{Состязания и «иголки»}

Из-за задержек на выходе комбинационной схемы могут появляться короткие
ложные импульсы --- \term{иголки} (glitches). Классический пример:
$y = a\cdot\overline{a}$. По законам алгебры $y$ всегда равен 0. Но инвертор
срабатывает с задержкой, и в течение короткого времени после фронта $a$
\emph{оба} входа элемента И равны 1.

\begin{figure}[H]
\centering
\begin{tikzpicture}
  \node[gAND, minimum width=10mm, minimum height=7mm] (g) at (2.2,0) {};
  \node[gNOT] (n) at (0.9,-0.7) {};
  \draw[wire] (-0.6,0.12) node[left] {$a$} -- (g.input 1);
  \fill (0.0,0.12) circle (1.3pt);
  \draw[wire] (0.0,0.12) |- (n.input);
  \draw[wire] (n.output) -- ++(0.3,0) |- (g.input 2);
  \draw[wire] (g.output) -- ++(0.5,0) node[right] {$y$};
  \begin{scope}[xshift=5cm, yshift=0.9cm, x=1cm, y=1cm]
    \draw[very thick, FpgaBlue] (0,0) -- (1.5,0) -- (1.55,0.5) -- (5,0.5);
    \draw[very thick, FpgaOrange] (0,-0.4) -- (2.4,-0.4) -- (2.45,-0.9) -- (5,-0.9);
    \draw[very thick, FpgaRed] (0,-1.8) -- (1.6,-1.8) -- (1.65,-1.3) -- (2.5,-1.3) -- (2.55,-1.8) -- (5,-1.8);
    \node[font=\sffamily\small, anchor=east] at (-0.1,0.25) {$a$};
    \node[font=\sffamily\small, anchor=east] at (-0.1,-0.65) {$\overline{a}$};
    \node[font=\sffamily\small, anchor=east] at (-0.1,-1.55) {$y$};
    \node[font=\sffamily\small, text=FpgaRed, anchor=west] at (2.7,-1.4) {ложный импульс!};
  \end{scope}
\end{tikzpicture}
\caption{Возникновение «иголки» из-за задержки инвертора}
\end{figure}

\begin{warnbox}[Как с этим бороться]
Иголки --- главная причина, по которой в серьёзных схемах выходы
комбинационной логики «читают» не когда попало, а только по фронту
тактового сигнала, когда все переходные процессы уже закончились. Это
называется \term{синхронным проектированием} (глава~\ref{ch:sync}), и все
схемы для ПЛИС строятся именно так.
\end{warnbox}

\begin{ideabox}
Комбинационная схема не имеет памяти. Её проектируют по таблице истинности,
упрощают алгебраически или по карте Карно. Реальные элементы имеют
задержку, из-за которой возможны кратковременные ложные импульсы.
\end{ideabox}

\chapter{Дешифраторы и шифраторы}
\label{ch:coders}

Из логических элементов строят более крупные «кирпичи» --- типовые
комбинационные узлы. Они встречаются почти в любом цифровом устройстве, и их
полезно узнавать с первого взгляда. Начнём с узлов, которые преобразуют
один способ записи числа (код) в другой.

\begin{figure}[H]
\centering
\begin{tikzpicture}[node distance=6mm and 32mm]
  \node[block, minimum width=34mm] (bin) {Двоичный код\\\scriptsize \code{10} (число 2)};
  \node[oblock, minimum width=34mm, right=of bin] (oh) {Унитарный код\\\scriptsize \code{0100} (единица на 2-м месте)};
  \draw[arr] ([yshift=2mm]bin.east) -- node[above, font=\sffamily\small] {дешифратор} ([yshift=2mm]oh.west);
  \draw[arr] ([yshift=-2mm]oh.west) -- node[below, font=\sffamily\small] {шифратор} ([yshift=-2mm]bin.east);
\end{tikzpicture}
\caption{Дешифратор и шифратор выполняют взаимно обратные преобразования}
\end{figure}

\term{Унитарный код} (one-hot, «одна горячая») --- код, в котором из $N$
проводов единица стоит ровно на одном, а номер этого провода и есть
передаваемое значение. Он удобен, чтобы что-то \emph{выбрать}: включить одну из
нескольких ламп, одну из нескольких микросхем памяти.

\section{Дешифратор}

\term{Дешифратор} (decoder) имеет $n$ входов и $2^n$ выходов. Входы задают
двоичный номер, и на выходе с этим номером появляется 1, а на всех остальных
--- 0. На схемах дешифратор обозначают буквами DC.

\begin{figure}[H]
\centering
\begin{minipage}{0.4\textwidth}
\centering
\begin{tikzpicture}
  \draw[thick, fill=FpgaLight] (0,0) rectangle (2.6,3.2);
  \draw[thick] (0.6,0) -- (0.6,3.2); \draw[thick] (2.0,0) -- (2.0,3.2);
  \node[font=\sffamily\bfseries\large] at (1.3,1.6) {DC};
  \foreach \i/\n in {0/0, 1/1} {
    \draw[wire] (-0.6,2.1-\i*0.8) node[left] {$a_\i$} -- (0,2.1-\i*0.8);
    \node[font=\sffamily\small] at (0.3,2.1-\i*0.8) {\pgfmathparse{int(2^\i)}\pgfmathresult};
  }
  \draw[wire] (-0.6,0.4) node[left] {$E$} -- (0,0.4); \node[font=\sffamily\small] at (0.3,0.4) {E};
  \foreach \i in {0,...,3} {
    \draw[wire] (2.6,2.6-\i*0.7) -- (3.2,2.6-\i*0.7) node[right] {$y_\i$};
    \node[font=\sffamily\small] at (2.3,2.6-\i*0.7) {\i};
  }
\end{tikzpicture}
\end{minipage}\hfill
\begin{minipage}{0.56\textwidth}
\centering
\begin{tabular}{ccc|cccc}
\toprule
$E$ & $a_1$ & $a_0$ & $y_0$ & $y_1$ & $y_2$ & $y_3$ \\
\midrule
0 & $\times$ & $\times$ & 0 & 0 & 0 & 0 \\
1 & 0 & 0 & \textbf{1} & 0 & 0 & 0 \\
1 & 0 & 1 & 0 & \textbf{1} & 0 & 0 \\
1 & 1 & 0 & 0 & 0 & \textbf{1} & 0 \\
1 & 1 & 1 & 0 & 0 & 0 & \textbf{1} \\
\bottomrule
\end{tabular}

\smallskip
{\small $\times$ --- «неважно, 0 или 1»}
\end{minipage}
\caption{Дешифратор 2$\to$4 с входом разрешения $E$: обозначение и таблица истинности}
\label{fig:dc24}
\end{figure}

Вход $E$ (Enable, разрешение) позволяет «выключить» дешифратор: при $E=0$
все выходы равны нулю, какой бы номер ни был подан.

Как построить дешифратор из элементов? Каждый выход --- это произведение,
равное единице ровно для одной комбинации входов:
\[
  y_0 = E\,\overline{a_1}\,\overline{a_0}, \quad
  y_1 = E\,\overline{a_1}\,a_0, \quad
  y_2 = E\,a_1\,\overline{a_0}, \quad
  y_3 = E\,a_1\,a_0.
\]

\begin{figure}[H]
\centering
\begin{tikzpicture}
  % входы
  \draw[wire] (-0.4,1.6) node[left] {$a_1$} -- (1.0,1.6);
  \draw[wire] (-0.4,0.8) node[left] {$a_0$} -- (2.4,0.8);
  \draw[wire] (-0.4,0.0) node[left] {$E$} -- (3.8,0.0);
  % инверторы
  \node[gNOT, minimum width=7mm, minimum height=5mm] (n1) at (1.0,2.4) {};
  \draw[wire] (0.3,1.6) |- (n1.input); \fill (0.3,1.6) circle (1.3pt);
  \draw[very thick, FpgaBlue] (n1.output) -| (1.6,-4.8);
  \node[gNOT, minimum width=7mm, minimum height=5mm] (n0) at (2.4,1.5) {};
  \draw[wire] (1.95,0.8) |- (n0.input); \fill (1.95,0.8) circle (1.3pt);
  \draw[very thick, FpgaBlue] (n0.output) -| (3.0,-4.8);
  % шины
  \draw[very thick, FpgaBlue] (1.0,1.6) -- (1.0,-4.8);
  \draw[very thick, FpgaBlue] (2.4,0.8) -- (2.4,-4.8);
  \draw[very thick, FpgaBlue] (3.8,0.0) -- (3.8,-4.8);
  \foreach \x/\l in {1.0/a_1, 1.6/\overline{a_1}, 2.4/a_0, 3.0/\overline{a_0}, 3.8/E}
    \node[below, font=\small] at (\x,-4.8) {$\l$};
  % элементы И: y0 = ~a1 ~a0, y1 = ~a1 a0, y2 = a1 ~a0, y3 = a1 a0
  \foreach \i/\x/\z in {0/1.6/3.0, 1/1.6/2.4, 2/1.0/3.0, 3/1.0/2.4} {
    \pgfmathsetmacro{\y}{-0.9-\i*1.1}
    \node[gAND, logic gate inputs=nnn, minimum height=9mm] (g\i) at (5.6,\y) {};
    \draw[wire] (g\i.input 1) -- (\x,0 |- g\i.input 1) node[circle, fill, inner sep=1.1pt] {};
    \draw[wire] (g\i.input 2) -- (\z,0 |- g\i.input 2) node[circle, fill, inner sep=1.1pt] {};
    \draw[wire] (g\i.input 3) -- (3.8,0 |- g\i.input 3) node[circle, fill, inner sep=1.1pt] {};
    \draw[wire] (g\i.output) -- ++(0.6,0) node[right] {$y_\i$};
  }
\end{tikzpicture}
\caption{Дешифратор 2$\to$4 на элементах И и НЕ}
\label{fig:dc24_gates}
\end{figure}

\begin{figure}[H]
\centering
\begin{tikztimingtable}[timing/slope=0.05, xscale=2.2, yscale=1.4, timing/font=\sffamily\small, timing/rowdist=3.2ex]
  $a_1a_0$ & 2D{00} 2D{01} 2D{10} 2D{11} 2D{00} \\
  $y_0$ & 2H 6L 2H \\
  $y_1$ & 2L 2H 6L \\
  $y_2$ & 4L 2H 4L \\
  $y_3$ & 6L 2H 2L \\
\end{tikztimingtable}
\caption{При переборе номеров единица «переходит» с одного выхода на другой}
\end{figure}

\subsection{Где применяются дешифраторы}

\begin{itemize}
  \item \textbf{Выбор устройства по адресу.} В компьютере к одной шине
        подключено много микросхем памяти и периферии. Дешифратор по
        старшим битам адреса выбирает, какая из них сейчас должна ответить.
  \item \textbf{Управление индикаторами}: включение одного из нескольких
        светодиодов, разрядов индикатора, строк матрицы.
  \item \textbf{Выбор команды в процессоре}: код операции дешифрируется в
        управляющие сигналы «сложить», «прочитать из памяти» и т.\,д.
\end{itemize}

\begin{figure}[H]
\centering
\begin{tikzpicture}[node distance=5mm]
  \node[hblock, minimum height=30mm, minimum width=18mm] (cpu) {Процессор};
  \node[block, right=22mm of cpu, minimum height=26mm] (dc) {DC\\2$\to$4};
  \foreach \i/\t in {0/ПЗУ, 1/ОЗУ, 2/Порт ввода, 3/Порт вывода} {
    \node[oblock, minimum width=24mm, minimum height=7mm] (d\i) at (9.0,1.65-\i*1.1) {\t};
  }
  \draw[arr] (cpu.east) -- node[above, lbl] {A[15:14]} (dc.west);
  \foreach \i in {0,...,3} \draw[arr] ($(dc.north east)!{(\i+0.5)/4}!(dc.south east)$) -- ++(0.6,0) |- (d\i.west);
  \foreach \i in {0,...,3} \node[lbl, above=0mm of d\i.north west, anchor=south west] {CS$_\i$};
\end{tikzpicture}
\caption{Дешифратор адреса: два старших бита адреса выбирают одну из четырёх микросхем}
\end{figure}

\subsection{Дешифратор для семисегментного индикатора}

Семисегментный индикатор состоит из семи светодиодных полосок, обозначаемых
буквами от $a$ до $g$. Зажигая разные сегменты, можно показать любую цифру.
Узел, который переводит двоичный код цифры в сигналы для сегментов, тоже
называют дешифратором (хотя он выдаёт не унитарный код).

\newcommand{\sevenseg}[8]{%
  \begin{tikzpicture}[scale=0.32, seg/.style={line width=2.6pt, line cap=round}]
    \def\segc##1{\ifnum##1=1 \def\cc{FpgaRed}\else\def\cc{gray!20}\fi}
    \segc{#1}\draw[seg, color=\cc] (0.2,4) -- (1.8,4);   % a
    \segc{#2}\draw[seg, color=\cc] (2,3.8) -- (2,2.2);   % b
    \segc{#3}\draw[seg, color=\cc] (2,1.8) -- (2,0.2);   % c
    \segc{#4}\draw[seg, color=\cc] (0.2,0) -- (1.8,0);   % d
    \segc{#5}\draw[seg, color=\cc] (0,1.8) -- (0,0.2);   % e
    \segc{#6}\draw[seg, color=\cc] (0,3.8) -- (0,2.2);   % f
    \segc{#7}\draw[seg, color=\cc] (0.2,2) -- (1.8,2);   % g
    \node[font=\sffamily\scriptsize] at (1,-1.1) {#8};
  \end{tikzpicture}}

\begin{figure}[H]
\centering
\begin{tikzpicture}[scale=0.8, seg/.style={line width=4pt, line cap=round, FpgaRed}]
  \draw[seg] (0.25,4) -- (1.75,4) node[midway, above=1mm, text=FpgaDark, font=\sffamily\small] {a};
  \draw[seg] (2,3.75) -- (2,2.25) node[midway, right=1mm, text=FpgaDark, font=\sffamily\small] {b};
  \draw[seg] (2,1.75) -- (2,0.25) node[midway, right=1mm, text=FpgaDark, font=\sffamily\small] {c};
  \draw[seg] (0.25,0) -- (1.75,0) node[midway, below=1mm, text=FpgaDark, font=\sffamily\small] {d};
  \draw[seg] (0,1.75) -- (0,0.25) node[midway, left=1mm, text=FpgaDark, font=\sffamily\small] {e};
  \draw[seg] (0,3.75) -- (0,2.25) node[midway, left=1mm, text=FpgaDark, font=\sffamily\small] {f};
  \draw[seg] (0.25,2) -- (1.75,2) node[midway, above=0.5mm, text=FpgaDark, font=\sffamily\small] {g};
\end{tikzpicture}
\hspace{12mm}
\sevenseg1111110{0}\,\sevenseg0110000{1}\,\sevenseg1101101{2}\,\sevenseg1111001{3}\,\sevenseg0110011{4}\,%
\sevenseg1011011{5}\,\sevenseg1011111{6}\,\sevenseg1110000{7}\,\sevenseg1111111{8}\,\sevenseg1111011{9}
\caption{Семисегментный индикатор: обозначение сегментов и изображение цифр}
\end{figure}

\begin{table}[H]
\centering
\caption{Таблица истинности дешифратора для семисегментного индикатора (1 --- сегмент горит)}
\label{tab:7seg}
\small
\begin{tabular}{c|cccc|ccccccc}
\toprule
Цифра & $x_3$ & $x_2$ & $x_1$ & $x_0$ & a & b & c & d & e & f & g \\
\midrule
0 & 0&0&0&0 & 1&1&1&1&1&1&0 \\
1 & 0&0&0&1 & 0&1&1&0&0&0&0 \\
2 & 0&0&1&0 & 1&1&0&1&1&0&1 \\
3 & 0&0&1&1 & 1&1&1&1&0&0&1 \\
4 & 0&1&0&0 & 0&1&1&0&0&1&1 \\
5 & 0&1&0&1 & 1&0&1&1&0&1&1 \\
6 & 0&1&1&0 & 1&0&1&1&1&1&1 \\
7 & 0&1&1&1 & 1&1&1&0&0&0&0 \\
8 & 1&0&0&0 & 1&1&1&1&1&1&1 \\
9 & 1&0&0&1 & 1&1&1&1&0&1&1 \\
\bottomrule
\end{tabular}
\end{table}

Каждый столбец a--g --- это отдельная логическая функция от четырёх
переменных. Её можно упростить по карте Карно и собрать из элементов, как в
главе~\ref{ch:comb}.

\section{Шифратор}

\term{Шифратор} (encoder, на схемах обозначается CD) выполняет обратное
преобразование: из унитарного кода на $2^n$ входах делает $n$-разрядный
двоичный номер той линии, на которой стоит единица.

\begin{figure}[H]
\centering
\begin{minipage}{0.38\textwidth}
\centering
\begin{tikzpicture}
  \draw[thick, fill=FpgaLight] (0,0) rectangle (2.6,3.0);
  \draw[thick] (0.6,0) -- (0.6,3.0); \draw[thick] (2.0,0) -- (2.0,3.0);
  \node[font=\sffamily\bfseries\large] at (1.3,1.5) {CD};
  \foreach \i in {0,...,3} {
    \draw[wire] (-0.6,2.4-\i*0.65) node[left] {$x_\i$} -- (0,2.4-\i*0.65);
    \node[font=\sffamily\small] at (0.3,2.4-\i*0.65) {\i};
  }
  \foreach \i/\w in {0/1, 1/2} {
    \draw[wire] (2.6,1.9-\i*0.9) -- (3.2,1.9-\i*0.9) node[right] {$y_\i$};
    \node[font=\sffamily\small] at (2.3,1.9-\i*0.9) {\w};
  }
\end{tikzpicture}
\end{minipage}\hfill
\begin{minipage}{0.58\textwidth}
\centering
\begin{tabular}{cccc|cc}
\toprule
$x_0$ & $x_1$ & $x_2$ & $x_3$ & $y_1$ & $y_0$ \\
\midrule
1 & 0 & 0 & 0 & 0 & 0 \\
0 & 1 & 0 & 0 & 0 & 1 \\
0 & 0 & 1 & 0 & 1 & 0 \\
0 & 0 & 0 & 1 & 1 & 1 \\
\bottomrule
\end{tabular}
\end{minipage}
\caption{Шифратор 4$\to$2: обозначение и таблица истинности}
\end{figure}

По таблице видно: $y_0$ равен 1, когда активна линия 1 или 3; $y_1$ --- когда
активна линия 2 или 3. Значит, шифратор собирается из одних элементов ИЛИ:
\[
  y_0 = x_1 + x_3, \qquad y_1 = x_2 + x_3.
\]

\begin{figure}[H]
\centering
\begin{tikzpicture}
  \foreach \i in {0,...,3} \draw[very thick, FpgaBlue] (\i*0.7,0.6) node[above, text=FpgaDark] {$x_\i$} -- (\i*0.7,-2.4);
  \node[gOR] (o0) at (4.2,-0.4) {};
  \node[gOR] (o1) at (4.2,-1.7) {};
  \draw[wire] (o0.input 1) -- (0.7,0 |- o0.input 1) node[circle, fill, inner sep=1.1pt] {};
  \draw[wire] (o0.input 2) -- (2.1,0 |- o0.input 2) node[circle, fill, inner sep=1.1pt] {};
  \draw[wire] (o1.input 1) -- (1.4,0 |- o1.input 1) node[circle, fill, inner sep=1.1pt] {};
  \draw[wire] (o1.input 2) -- (2.1,0 |- o1.input 2) node[circle, fill, inner sep=1.1pt] {};
  \draw[wire] (o0.output) -- ++(0.6,0) node[right] {$y_0$};
  \draw[wire] (o1.output) -- ++(0.6,0) node[right] {$y_1$};
  \node[lbl, text=FpgaGray] at (0,-2.7) {не используется};
\end{tikzpicture}
\caption{Шифратор 4$\to$2 из двух элементов ИЛИ}
\end{figure}

\subsection{Приоритетный шифратор}

Что будет, если единица окажется сразу на двух входах, например, если
одновременно нажаты кнопки 1 и 2? Простой шифратор выдаст $y_1y_0 = 11$, то есть
номер 3, --- явную ошибку. Поэтому на практике применяют
\term{приоритетный шифратор}: он выдаёт номер \emph{старшего} из активных
входов, а остальные игнорирует. Дополнительный выход $V$ (valid) показывает,
что хотя бы один вход активен.

\begin{table}[H]
\centering
\caption{Приоритетный шифратор 4$\to$2 ($\times$ --- значение не важно)}
\begin{tabular}{cccc|ccc}
\toprule
$x_3$ & $x_2$ & $x_1$ & $x_0$ & $y_1$ & $y_0$ & $V$ \\
\midrule
0 & 0 & 0 & 0 & 0 & 0 & 0 \\
0 & 0 & 0 & 1 & 0 & 0 & 1 \\
0 & 0 & 1 & $\times$ & 0 & 1 & 1 \\
0 & 1 & $\times$ & $\times$ & 1 & 0 & 1 \\
1 & $\times$ & $\times$ & $\times$ & 1 & 1 & 1 \\
\bottomrule
\end{tabular}
\end{table}

Выражения: $y_1 = x_3 + x_2$, \ $y_0 = x_3 + \overline{x_2}\,x_1$, \
$V = x_3 + x_2 + x_1 + x_0$.

\begin{infobox}[Где применяются шифраторы]
\begin{itemize}
  \item \textbf{Клавиатуры}: номер нажатой клавиши превращается в код.
  \item \textbf{Контроллеры прерываний}: из нескольких одновременных запросов
        процессору передаётся номер самого важного.
  \item \textbf{Датчики положения}: из показаний линейки фотодатчиков
        получают число --- координату.
\end{itemize}
\end{infobox}

\section{Преобразователи кодов}

Дешифраторы и шифраторы --- частные случаи \term{преобразователей кодов}.
Ещё один полезный пример --- перевод двоичного кода в \term{код Грея}, в
котором соседние числа отличаются ровно одним битом. Это важно для датчиков
угла поворота: при переходе от 3 (\code{011}) к 4 (\code{100}) в двоичном коде
меняются сразу три бита, и в промежуточный момент датчик может выдать
что угодно. В коде Грея меняется только один бит.

\begin{figure}[H]
\centering
\begin{minipage}{0.45\textwidth}
\centering
\begin{tabular}{c c c}
\toprule
\textbf{Число} & \textbf{Двоичный} & \textbf{Грея} \\
\midrule
0 & \code{000} & \code{000} \\
1 & \code{001} & \code{001} \\
2 & \code{010} & \code{011} \\
3 & \code{011} & \code{010} \\
4 & \code{100} & \code{110} \\
5 & \code{101} & \code{111} \\
6 & \code{110} & \code{101} \\
7 & \code{111} & \code{100} \\
\bottomrule
\end{tabular}
\end{minipage}\hfill
\begin{minipage}{0.5\textwidth}
\centering
\begin{tikzpicture}
  \draw[very thick, FpgaBlue] (0,0) node[left, text=FpgaDark] {$b_2$} -- (3.2,0) node[right, text=FpgaDark] {$g_2$};
  \node[gXOR, minimum width=10mm, minimum height=7mm] (x1) at (2.2,-1.0) {};
  \node[gXOR, minimum width=10mm, minimum height=7mm] (x0) at (2.2,-2.2) {};
  \draw[wire] (x1.input 1) -- (0.6,0 |- x1.input 1) -- (0.6,0); \fill (0.6,0) circle (1.3pt);
  \draw[wire] (x1.input 2) -- (0,0 |- x1.input 2) node[left] {$b_1$};
  \fill (1.0,0 |- x1.input 2) circle (1.3pt);
  \draw[wire] (x0.input 1) -- (1.0,0 |- x0.input 1) -- (1.0,0 |- x1.input 2);
  \draw[wire] (x0.input 2) -- (0,0 |- x0.input 2) node[left] {$b_0$};
  \draw[wire] (x1.output) -- ++(0.5,0) node[right] {$g_1$};
  \draw[wire] (x0.output) -- ++(0.5,0) node[right] {$g_0$};
\end{tikzpicture}

\small $g_i = b_{i+1} \oplus b_i$
\end{minipage}
\caption{Двоичный код и код Грея; преобразователь на элементах «Исключающее ИЛИ»}
\end{figure}

\begin{ideabox}
\textbf{Дешифратор}: номер $\to$ единица на одном из $2^n$ выходов.
\textbf{Шифратор}: единица на одном из входов $\to$ её номер.
\textbf{Приоритетный шифратор} при нескольких активных входах выбирает старший.
\end{ideabox}

\chapter{Мультиплексоры и демультиплексоры}
\label{ch:mux}

\section{Цифровой переключатель}

Представьте старый галетный переключатель: ручка поворачивается и соединяет
общий контакт с одним из нескольких. \term{Мультиплексор} (MUX) делает то же
самое с цифровыми сигналами: у него несколько информационных входов, один
выход и \term{адресные} (управляющие) входы, которые выбирают, какой из
информационных входов соединить с выходом.

\begin{figure}[H]
\centering
\begin{tikzpicture}
  % механический переключатель
  \foreach \i in {0,...,3} {
    \fill (0,1.5-\i) circle (2pt);
    \draw[wire] (-1,1.5-\i) node[left] {$x_\i$} -- (0,1.5-\i);
  }
  \fill (2,0) circle (2.5pt);
  \draw[very thick, FpgaOrange] (2,0) -- (0.08,0.5);
  \draw[wire] (2,0) -- (3,0) node[right] {$y$};
  \node[lbl] at (1.2,-2.3) {механический переключатель};
  % мультиплексор
  \begin{scope}[xshift=7.5cm]
    \draw[thick, fill=orange!15] (0,2) -- (1.6,1.2) -- (1.6,-1.2) -- (0,-2) -- cycle;
    \node[font=\sffamily\bfseries] at (0.8,0) {MUX};
    \foreach \i in {0,...,3} {
      \draw[wire] (-1,1.5-\i) node[left] {$x_\i$} -- (0,1.5-\i);
      \node[font=\sffamily\scriptsize, anchor=west] at (0.05,1.5-\i) {\i};
    }
    \draw[wire] (1.6,0) -- (2.6,0) node[right] {$y$};
    \draw[wire, FpgaBlue] (0.6,-1.7) -- (0.6,-2.6) node[below] {$s_1$};
    \draw[wire, FpgaBlue] (1.1,-1.45) -- (1.1,-2.6) node[below] {$s_0$};
    \node[lbl] at (0.8,-3.3) {цифровой мультиплексор 4:1};
  \end{scope}
  \draw[-{Stealth[length=4mm]}, line width=2pt, FpgaOrange] (4.2,0) -- (5.6,0);
\end{tikzpicture}
\caption{Мультиплексор --- это управляемый переключатель: при $s_1s_0 = 01$
выход соединён со входом $x_1$}
\end{figure}

Мультиплексор с $n$ адресными входами может выбирать из $2^n$
информационных входов. Его обозначают дробью: \textbf{2:1}, \textbf{4:1},
\textbf{8:1} (сколько входов : сколько выходов).

\section{Мультиплексор 2:1}

Самый простой мультиплексор имеет два входа $a$, $b$ и один адресный вход $s$:
если $s = 0$, на выход проходит $a$; если $s = 1$ --- проходит $b$.

\begin{figure}[H]
\centering
\begin{minipage}{0.28\textwidth}
\centering
\begin{tikzpicture}
  \draw[thick, fill=orange!15] (0,1.1) -- (1.1,0.6) -- (1.1,-0.6) -- (0,-1.1) -- cycle;
  \node[font=\sffamily\small] at (0.55,0) {MUX};
  \draw[wire] (-0.7,0.6) node[left] {$a$} -- (0,0.6); \node[font=\sffamily\scriptsize, anchor=west] at (0.02,0.6) {0};
  \draw[wire] (-0.7,-0.6) node[left] {$b$} -- (0,-0.6); \node[font=\sffamily\scriptsize, anchor=west] at (0.02,-0.6) {1};
  \draw[wire] (1.1,0) -- (1.7,0) node[right] {$y$};
  \draw[wire, FpgaBlue] (0.55,-0.85) -- (0.55,-1.5) node[below] {$s$};
\end{tikzpicture}
\end{minipage}
\begin{minipage}{0.27\textwidth}
\centering
\begin{tabular}{c|c}
\toprule
$s$ & $y$ \\
\midrule
0 & $a$ \\
1 & $b$ \\
\bottomrule
\end{tabular}

\medskip
$y = \overline{s}\,a + s\,b$
\end{minipage}
\begin{minipage}{0.43\textwidth}
\centering
\begin{tikzpicture}
  \node[gAND, minimum width=10mm, minimum height=7mm] (g1) at (2,0.6) {};
  \node[gAND, minimum width=10mm, minimum height=7mm] (g2) at (2,-0.6) {};
  \node[gOR, minimum width=10mm, minimum height=7mm] (o) at (3.7,0) {};
  \node[gNOT, minimum width=7mm, minimum height=5mm] (n) at (0.9,0.48) {};
  \draw[wire] (g1.input 1) -- ++(-1.8,0) node[left] {$a$};
  \draw[wire] (g2.input 2) -- ++(-1.8,0) node[left] {$b$};
  \draw[wire, FpgaBlue] (0.2,-1.4) node[below] {$s$} -- (0.2,0.48) -- (n.input);
  \draw[wire, FpgaBlue] (0.2,-0.48) -- (g2.input 1); \fill[FpgaBlue] (0.2,-0.48) circle (1.3pt);
  \draw[wire] (n.output) -- (g1.input 2);
  \draw[wire] (g1.output) -- ++(0.15,0) |- (o.input 1);
  \draw[wire] (g2.output) -- ++(0.15,0) |- (o.input 2);
  \draw[wire] (o.output) -- ++(0.4,0) node[right] {$y$};
\end{tikzpicture}
\end{minipage}
\caption{Мультиплексор 2:1: обозначение, таблица, формула и схема на элементах}
\label{fig:mux21}
\end{figure}

Как работает схема: при $s = 0$ верхний элемент И «пропускает» $a$, а нижний
«закрыт» (на его входе 0). При $s = 1$ --- наоборот. Элемент ИЛИ собирает
результат.

\section{Мультиплексор 4:1}

У мультиплексора 4:1 два адресных входа $s_1 s_0$ задают номер входа от 0 до 3.
\[
  y = \overline{s_1}\,\overline{s_0}\,x_0 + \overline{s_1}\,s_0\,x_1 + s_1\,\overline{s_0}\,x_2 + s_1\,s_0\,x_3.
\]
Присмотритесь к формуле: множители при $x_i$ --- это ровно выходы
\emph{дешифратора} 2$\to$4 из главы~\ref{ch:coders}! Значит, мультиплексор
состоит из дешифратора адреса, четырёх элементов И и одного ИЛИ.

\begin{figure}[H]
\centering
\begin{tikzpicture}
  \draw[thick, fill=FpgaLight] (0,-2.6) rectangle (1.6,1.0);
  \node[font=\sffamily\bfseries] at (0.8,-0.8) {DC};
  \draw[wire, FpgaBlue] (-0.8,0.2) node[left, text=FpgaDark] {$s_1$} -- (0,0.2);
  \draw[wire, FpgaBlue] (-0.8,-1.8) node[left, text=FpgaDark] {$s_0$} -- (0,-1.8);
  \foreach \i in {0,...,3} {
    \pgfmathsetmacro{\y}{0.6-\i*1.0}
    \node[gAND, minimum width=10mm, minimum height=7mm] (g\i) at (4.6,\y-0.12) {};
    \draw[wire] (1.6,\y) -- ++(0.8,0) |- (g\i.input 2);
    \node[font=\sffamily\scriptsize, anchor=east] at (1.55,\y) {\i};
    \draw[wire] (g\i.input 1) -- ++(-0.8,0) node[left, font=\small] {$x_\i$};
  }
  \node[gOR, logic gate inputs=nnnn, minimum height=18mm] (o) at (7.2,-0.9) {};
  \foreach \i in {0,...,3} {
    \pgfmathtruncatemacro{\j}{\i+1}
    \draw[wire] (g\i.output) -- ++(0.35,0) |- (o.input \j);
  }
  \draw[wire] (o.output) -- ++(0.7,0) node[right] {$y$};
\end{tikzpicture}
\caption{Мультиплексор 4:1 = дешифратор адреса + элементы И + элемент ИЛИ}
\label{fig:mux41}
\end{figure}

\begin{figure}[H]
\centering
\begin{tikztimingtable}[timing/slope=0.05, xscale=0.9, yscale=1.4, timing/font=\sffamily\small, timing/rowdist=3.2ex]
  $x_0$     & 16{1L 1H} \\
  $x_1$     & 8{2L 2H} \\
  $x_2$     & 32H \\
  $x_3$     & 32L \\
  $s_1s_0$  & 8D{00} 8D{01} 8D{10} 8D{11} \\
  $y$       & 4{1L 1H} 2{2L 2H} 8H 8L \\
\end{tikztimingtable}
\caption{Мультиплексор 4:1 по очереди передаёт на выход разные входные
сигналы}
\end{figure}

\section{Дерево мультиплексоров}

Большой мультиплексор можно собрать из маленьких. Мультиплексор 4:1 --- это
три мультиплексора 2:1: первый уровень выбирает по младшему биту адреса,
второй --- по старшему. Мультиплексор 8:1 --- дерево из семи мультиплексоров
2:1 в три уровня, и так далее.

\newcommand{\muxsym}[3]{% имя, x, y
  \draw[thick, fill=orange!15] (#2,#3+0.55) -- (#2+0.7,#3+0.3) -- (#2+0.7,#3-0.3) -- (#2,#3-0.55) -- cycle;
  \coordinate (#1-i0) at (#2,#3+0.3); \coordinate (#1-i1) at (#2,#3-0.3);
  \coordinate (#1-o) at (#2+0.7,#3); \coordinate (#1-s) at (#2+0.35,#3-0.42);}

\begin{figure}[H]
\centering
\begin{tikzpicture}
  \muxsym{m0}{0}{1}
  \muxsym{m1}{0}{-1}
  \muxsym{m2}{2}{0}
  \foreach \m/\a/\b in {m0/x_0/x_1, m1/x_2/x_3} {
    \draw[wire] (\m-i0) -- ++(-0.6,0) node[left] {$\a$};
    \draw[wire] (\m-i1) -- ++(-0.6,0) node[left] {$\b$};
  }
  \draw[wire] (m0-o) -- ++(0.4,0) |- (m2-i0);
  \draw[wire] (m1-o) -- ++(0.4,0) |- (m2-i1);
  \draw[wire] (m2-o) -- ++(0.6,0) node[right] {$y$};
  \draw[wire, FpgaBlue] (m0-s) -- (0.35,0.25) -- (-0.9,0.25) node[left, text=FpgaDark] {$s_0$};
  \draw[wire, FpgaBlue] (m1-s) -- (0.35,-1.8) -- (-0.4,-1.8) -- (-0.4,0.25); \fill[FpgaBlue] (-0.4,0.25) circle (1.3pt);
  \draw[wire, FpgaBlue] (m2-s) -- (2.35,-1.2) node[below, text=FpgaDark] {$s_1$};
  \node[lbl] at (0.35,-2.4) {уровень 1};
  \node[lbl] at (2.35,-2.4) {уровень 2};
\end{tikzpicture}
\caption{Мультиплексор 4:1 из трёх мультиплексоров 2:1}
\label{fig:muxtree}
\end{figure}

\section{Мультиплексор как универсальный логический элемент}

Вот удивительное свойство мультиплексора. Подадим на адресные входы
\emph{переменные} $a$ и $b$, а на информационные входы --- \emph{константы}
0 и 1, взятые из правого столбца таблицы истинности какой-либо функции. Тогда
мультиплексор будет выбирать нужную строку таблицы и выдавать значение
функции! Любую функцию от $n$ переменных можно реализовать одним
мультиплексором $2^n$:1.

\begin{figure}[H]
\centering
\begin{tikzpicture}
  \draw[thick, fill=orange!15] (0,2) -- (1.6,1.2) -- (1.6,-1.2) -- (0,-2) -- cycle;
  \node[font=\sffamily\bfseries] at (0.8,0) {MUX};
  \foreach \i/\v/\ab in {0/0/00, 1/1/01, 2/1/10, 3/0/11} {
    \draw[wire] (-0.8,1.5-\i) -- (0,1.5-\i);
    \node[draw=FpgaGreen, fill=green!8, thick, minimum size=6mm, font=\ttfamily] at (-1.1,1.5-\i) {\v};
    \node[font=\sffamily\scriptsize, anchor=west] at (0.05,1.5-\i) {\i};
    \node[lbl, anchor=east] at (-1.5,1.5-\i) {$ab=\ab$};
  }
  \draw[wire] (1.6,0) -- (2.6,0) node[right] {$y = a\oplus b$};
  \draw[wire, FpgaBlue] (0.6,-1.7) -- (0.6,-2.6) node[below, text=FpgaDark] {$a$};
  \draw[wire, FpgaBlue] (1.1,-1.45) -- (1.1,-2.6) node[below, text=FpgaDark] {$b$};
  \node[lbl, align=center, text=FpgaGreen] at (-1.1,2.3) {таблица\\истинности};
  \node[font=\sffamily\small, align=left, anchor=west] at (5.0,0.8) {Подставьте \code{0001} --- получится И,};
  \node[font=\sffamily\small, align=left, anchor=west] at (5.0,0.2) {\code{0111} --- ИЛИ, \code{1110} --- И-НЕ\ldots};
  \node[font=\sffamily\small, align=left, anchor=west] at (5.0,-0.4) {Схема та же --- меняются только};
  \node[font=\sffamily\small, align=left, anchor=west] at (5.0,-1.0) {четыре бита слева.};
\end{tikzpicture}
\caption{Мультиплексор 4:1 с константами на входах вычисляет любую функцию
двух переменных (здесь --- Исключающее ИЛИ)}
\label{fig:mux_lut}
\end{figure}

\begin{ideabox}[Запомните этот рисунок]
Мультиплексор, на входы которого поданы биты из маленькой памяти, --- это и
есть \textbf{LUT}, основной строительный блок ПЛИС. Во второй части книги
(глава~\ref{ch:inside}) мы увидим, что ПЛИС состоит из тысяч таких
«программируемых мультиплексоров».
\end{ideabox}

\section{Демультиплексор}

\term{Демультиплексор} (DEMUX) выполняет обратное действие: у него один
информационный вход $x$ и $2^n$ выходов. Адрес $s$ определяет, на какой из
выходов будет передан входной сигнал; на остальных выходах --- 0.

\begin{figure}[H]
\centering
\begin{minipage}{0.35\textwidth}
\centering
\begin{tikzpicture}
  \draw[thick, fill=orange!15] (0,1.2) -- (1.6,2) -- (1.6,-2) -- (0,-1.2) -- cycle;
  \node[font=\sffamily\small\bfseries] at (0.8,0) {DEMUX};
  \draw[wire] (-0.8,0) node[left] {$x$} -- (0,0);
  \foreach \i in {0,...,3} {
    \draw[wire] (1.6,1.5-\i) -- (2.4,1.5-\i) node[right] {$y_\i$};
    \node[font=\sffamily\scriptsize, anchor=east] at (1.55,1.5-\i) {\i};
  }
  \draw[wire, FpgaBlue] (0.5,-1.45) -- (0.5,-2.6) node[below, text=FpgaDark] {$s_1$};
  \draw[wire, FpgaBlue] (1.0,-1.7) -- (1.0,-2.6) node[below, text=FpgaDark] {$s_0$};
\end{tikzpicture}
\end{minipage}\hfill
\begin{minipage}{0.6\textwidth}
\centering
\begin{tabular}{cc|cccc}
\toprule
$s_1$ & $s_0$ & $y_0$ & $y_1$ & $y_2$ & $y_3$ \\
\midrule
0 & 0 & $x$ & 0 & 0 & 0 \\
0 & 1 & 0 & $x$ & 0 & 0 \\
1 & 0 & 0 & 0 & $x$ & 0 \\
1 & 1 & 0 & 0 & 0 & $x$ \\
\bottomrule
\end{tabular}

\medskip
$y_0 = x\,\overline{s_1}\,\overline{s_0}, \ \ y_1 = x\,\overline{s_1}\,s_0, \ \ \ldots$
\end{minipage}
\caption{Демультиплексор 1:4}
\end{figure}

Сравните формулы с дешифратором со входом разрешения
(рис.~\ref{fig:dc24}): они совпадают, если вход $E$ назвать $x$. Поэтому
\textbf{демультиплексор --- это дешифратор, у которого на вход разрешения
подаётся информационный сигнал}. В справочниках эти две микросхемы часто
объединяют под одним названием.

\section{Передача по одной линии}

Вместе мультиплексор и демультиплексор позволяют передать несколько сигналов
по \emph{одному} проводу. Адрес быстро перебирает значения $0, 1, 2, 3, 0,
1\ldots$, и в каждый момент по линии идёт сигнал одного из источников. На
приёмной стороне демультиплексор с таким же адресом раскладывает их по своим
выходам. Это \term{временное разделение каналов}: так работают телефонные
линии, шины компьютеров, а в дронах --- приёмник пульта, передающий по одному
проводу сигналы всех каналов (протоколы PPM, SBUS).

\begin{figure}[H]
\centering
\begin{tikzpicture}
  \draw[thick, fill=orange!15] (0,1.3) -- (1.2,0.7) -- (1.2,-0.7) -- (0,-1.3) -- cycle;
  \node[font=\sffamily\scriptsize\bfseries] at (0.6,0) {MUX};
  \foreach \i/\t in {0/газ, 1/крен, 2/тангаж, 3/рыскание}
    \draw[wire] (-0.6,0.9-\i*0.6) node[left, font=\sffamily\small] {\t} -- (0,0.9-\i*0.6);
  \draw[very thick, FpgaOrange] (1.2,0) -- node[above, font=\sffamily\small, text=FpgaDark] {одна линия} (5.0,0);
  \begin{scope}[xshift=5cm]
  \draw[thick, fill=orange!15] (0,0.7) -- (1.2,1.3) -- (1.2,-1.3) -- (0,-0.7) -- cycle;
  \node[font=\sffamily\tiny\bfseries] at (0.6,0) {DEMUX};
  \foreach \i/\t in {0/газ, 1/крен, 2/тангаж, 3/рыскание}
    \draw[wire] (1.2,0.9-\i*0.6) -- (1.8,0.9-\i*0.6) node[right, font=\sffamily\small] {\t};
  \end{scope}
  \node[block, minimum width=20mm, minimum height=7mm] (ctr) at (3.1,-2.1) {Счётчик\\\scriptsize адреса};
  \draw[arr, FpgaBlue] (ctr.west) -| node[pos=0.25, below, lbl] {$s$} (0.6,-1.05);
  \draw[arr, FpgaBlue] (ctr.east) -| node[pos=0.25, below, lbl] {$s$} (5.6,-1.05);
\end{tikzpicture}
\caption{Четыре канала пульта по одному проводу}
\end{figure}

\section{Многоразрядные мультиплексоры}

Часто нужно выбирать не между битами, а между \emph{числами}. Для этого
несколько мультиплексоров ставят параллельно с общим адресом. Например,
8-разрядный мультиплексор 2:1 --- это восемь мультиплексоров 2:1, у
которых соединены адресные входы. На схемах такой узел рисуют одним символом
с шинами.

\begin{figure}[H]
\centering
\begin{tikzpicture}
  \draw[thick, fill=orange!15] (0,1.1) -- (1.1,0.6) -- (1.1,-0.6) -- (0,-1.1) -- cycle;
  \node[font=\sffamily\small] at (0.55,0) {MUX};
  \draw[bus] (-1.4,0.6) node[left, text=FpgaDark] {$A$} -- (0,0.6);
  \draw[bus] (-1.4,-0.6) node[left, text=FpgaDark] {$B$} -- (0,-0.6);
  \draw[bus] (1.1,0) -- (2.4,0) node[right, text=FpgaDark] {$Y$};
  \foreach \x/\y in {-0.8/0.6, -0.8/-0.6, 1.7/0} {
    \draw[thick] (\x-0.1,\y-0.15) -- (\x+0.1,\y+0.15);
    \node[font=\sffamily\scriptsize, above] at (\x,\y+0.1) {8};
  }
  \draw[wire, FpgaBlue] (0.55,-0.85) -- (0.55,-1.5) node[below, text=FpgaDark] {$s$};
  \node[font=\sffamily\small, anchor=west] at (4.0,0.3) {$s=0$: \ $Y = A$};
  \node[font=\sffamily\small, anchor=west] at (4.0,-0.3) {$s=1$: \ $Y = B$};
\end{tikzpicture}
\caption{Мультиплексор 8-разрядных чисел. Косая черта с цифрой показывает число проводов в шине}
\end{figure}

\begin{ideabox}
\textbf{Мультиплексор} выбирает один из $2^n$ входов по адресу. Он состоит из
дешифратора адреса и элементов И-ИЛИ и может реализовать любую функцию, если
подать на входы константы. \textbf{Демультиплексор} передаёт один вход на
выбранный выход --- это дешифратор с информационным входом разрешения.
\end{ideabox}

\chapter{Арифметические узлы}
\label{ch:arith}

Компьютер --- по-английски «вычислитель», и главное, что он умеет, ---
считать. В этой главе мы увидим, что сложение, вычитание и сравнение чисел
выполняются всё теми же логическими элементами.

\section{Сложение двоичных чисел}

Двоичные числа складываются «столбиком», как десятичные, только перенос в
следующий разряд возникает, когда сумма достигает двух:
\[
  0+0 = 0, \quad 0+1 = 1, \quad 1+0 = 1, \quad 1+1 = 10_2 \ (\text{ноль и перенос единицы}).
\]

\begin{figure}[H]
\centering
\begin{tikzpicture}[x=9mm, y=7mm, font=\ttfamily\large]
  \node[text=FpgaRed, font=\ttfamily\small] at (1,3) {1}; \node[text=FpgaRed, font=\ttfamily\small] at (2,3) {1};
  \node[text=FpgaRed, font=\ttfamily\small] at (3,3) {1};
  \foreach \d [count=\i] in {0,1,1,1} \node at (\i,2) {\d};
  \foreach \d [count=\i] in {0,0,1,1} \node at (\i,1) {\d};
  \foreach \d [count=\i from 0] in {1,0,1,0} \node at (\i+1,-0.2) {\d};
  \node at (-0.2,1.5) {$+$};
  \draw[thick] (0.4,0.5) -- (4.5,0.5);
  \node[font=\sffamily\small, anchor=west] at (5.2,2) {$0111_2 = 7$};
  \node[font=\sffamily\small, anchor=west] at (5.2,1) {$0011_2 = 3$};
  \node[font=\sffamily\small, anchor=west] at (5.2,-0.2) {$1010_2 = 10$};
  \node[font=\sffamily\small, anchor=west, text=FpgaRed] at (5.2,3) {переносы};
\end{tikzpicture}
\caption{Сложение столбиком в двоичной системе}
\end{figure}

\section{Полусумматор}

Сложение двух одноразрядных чисел $a$ и $b$ даёт двухразрядный результат:
бит суммы $s$ и бит переноса $c$. Таблица истинности:

\begin{figure}[H]
\centering
\begin{minipage}{0.35\textwidth}
\centering
\begin{tabular}{cc|cc}
\toprule
$a$ & $b$ & $c$ & $s$ \\
\midrule
0 & 0 & 0 & 0 \\
0 & 1 & 0 & 1 \\
1 & 0 & 0 & 1 \\
1 & 1 & 1 & 0 \\
\bottomrule
\end{tabular}

\medskip
$s = a \oplus b, \qquad c = a\cdot b$
\end{minipage}
\begin{minipage}{0.55\textwidth}
\centering
\begin{tikzpicture}
  \node[gXOR, minimum height=12mm] (x) at (2,0.8) {};
  \node[gAND, minimum height=12mm] (g) at (2,-0.8) {};
  \draw[wire] (x.input 1) -- ++(-2.0,0) node[left] {$a$};
  \draw[wire] (x.input 2) -- ++(-2.0,0) node[left] {$b$};
  \coordinate (A) at ($(x.input 1)+(-1.4,0)$); \coordinate (B) at ($(x.input 2)+(-0.9,0)$);
  \draw[wire] (A) |- (g.input 1); \fill (A) circle (1.3pt);
  \draw[wire] (B) |- (g.input 2); \fill (B) circle (1.3pt);
  \draw[wire] (x.output) -- ++(0.6,0) node[right] {$s$ (сумма)};
  \draw[wire] (g.output) -- ++(0.6,0) node[right] {$c$ (перенос)};
\end{tikzpicture}
\end{minipage}
\caption{Полусумматор (half adder)}
\label{fig:ha}
\end{figure}

\section{Полный сумматор}

Полусумматор годится только для младшего разряда. В остальных разрядах
складываются \emph{три} бита: $a_i$, $b_i$ и перенос $c_\text{in}$ из
предыдущего разряда. Такой узел называется \term{полным сумматором}.

\begin{figure}[H]
\centering
\begin{minipage}{0.38\textwidth}
\centering
\small
\begin{tabular}{ccc|cc}
\toprule
$a$ & $b$ & $c_\text{in}$ & $c_\text{out}$ & $s$ \\
\midrule
0&0&0 & 0&0 \\
0&0&1 & 0&1 \\
0&1&0 & 0&1 \\
0&1&1 & 1&0 \\
1&0&0 & 0&1 \\
1&0&1 & 1&0 \\
1&1&0 & 1&0 \\
1&1&1 & 1&1 \\
\bottomrule
\end{tabular}
\end{minipage}
\begin{minipage}{0.58\textwidth}
\[
  s = a \oplus b \oplus c_\text{in}
\]
\[
  c_\text{out} = ab + ac_\text{in} + bc_\text{in}
\]
\small Заметьте: $c_\text{out}$ --- это мажоритарная функция «2 из 3» из
главы~\ref{ch:comb}! Перенос возникает, когда единиц хотя бы две.
\end{minipage}
\caption{Полный сумматор (full adder): таблица истинности и формулы}
\end{figure}

\begin{figure}[H]
\centering
\begin{tikzpicture}[ha/.style={draw=FpgaBlue, fill=FpgaLight, thick, minimum width=14mm, minimum height=14mm, font=\sffamily\small, align=center}]
  \node[ha] (h1) at (0,0) {HA};
  \node[ha] (h2) at (2.8,-0.6) {HA};
  \node[gOR, minimum width=10mm, minimum height=7mm] (o) at (5.4,-1.8) {};
  \draw[wire] ([yshift=3.5mm]h1.west) -- ++(-0.8,0) node[left] {$a$};
  \draw[wire] ([yshift=-3.5mm]h1.west) -- ++(-0.8,0) node[left] {$b$};
  \draw[wire] ([yshift=3.5mm]h1.east) -- node[above, lbl] {$s$} ++(0.4,0) |- ([yshift=3.5mm]h2.west);
  \draw[wire] (-0.8,-1.4) node[left] {$c_\text{in}$} -- (1.7,-1.4) |- ([yshift=-3.5mm]h2.west);
  \draw[wire] ([yshift=3.5mm]h2.east) -- ++(3.4,0) node[right] {$s$};
  \draw[wire] ([yshift=-3.5mm]h1.east) -- node[below, lbl, pos=0.3] {$c$} ++(0.25,0) |- (0.95,-2.2) -- (4.7,-2.2) |- (o.input 2);
  \draw[wire] ([yshift=-3.5mm]h2.east) -- node[above, lbl] {$c$} ++(0.4,0) |- (o.input 1);
  \draw[wire] (o.output) -- ++(0.6,0) node[right] {$c_\text{out}$};
\end{tikzpicture}
\caption{Полный сумматор из двух полусумматоров и элемента ИЛИ}
\label{fig:fa}
\end{figure}

\section{Многоразрядный сумматор}

Чтобы сложить $n$-разрядные числа, $n$ полных сумматоров соединяют в цепочку:
выход переноса каждого разряда идёт на вход переноса следующего. Такой
сумматор называется \term{сумматором с последовательным переносом} (ripple
carry adder).

\begin{figure}[H]
\centering
\begin{tikzpicture}[fa/.style={draw=FpgaBlue, fill=FpgaLight, thick, minimum width=15mm, minimum height=13mm, font=\sffamily\small}]
  \foreach \i [evaluate=\i as \x using (3-\i)*2.9] in {0,...,3} {
    \node[fa] (f\i) at (\x,0) {FA$_\i$};
    \draw[wire] ([xshift=-3mm]f\i.north) -- ++(0,0.6) node[above] {$a_\i$};
    \draw[wire] ([xshift=3mm]f\i.north) -- ++(0,0.6) node[above] {$b_\i$};
    \draw[wire] (f\i.south) -- ++(0,-0.6) node[below] {$s_\i$};
  }
  \foreach \i [evaluate=\i as \j using int(\i+1)] in {0,1,2}
    \draw[arr, FpgaRed] (f\i.west) -- node[above, lbl] {$c_\j$} (f\j.east);
  \draw[arr] ($(f0.east)+(0.8,0)$) node[right] {$c_0 = 0$} -- (f0.east);
  \draw[arr, FpgaRed] (f3.west) -- ++(-0.8,0) node[left, text=FpgaDark] {$c_4$ (перенос)};
\end{tikzpicture}
\caption{Четырёхразрядный сумматор с последовательным переносом. Красным ---
цепочка переноса}
\label{fig:rca}
\end{figure}

Недостаток такого сумматора: старший разряд не может посчитать результат,
пока не получит перенос от младшего, а тот --- от ещё более младшего.
Задержка растёт пропорционально числу разрядов. Поэтому в быстрых
процессорах используют сумматоры с \term{ускоренным переносом}, в которых
перенос вычисляется параллельно специальной схемой. А в ПЛИС для этой цели
рядом с каждым LUT проложена отдельная быстрая \emph{цепочка переноса}.

\begin{figure}[H]
\centering
\begin{tikzpicture}
\begin{axis}[
  width=0.8\textwidth, height=5.4cm, xlabel={Разрядность, бит}, ylabel={Задержка, $t_\text{зд}$},
  xmin=0, xmax=64, ymin=0, ymax=140, xtick={8,16,32,48,64},
  legend pos=north west, legend style={font=\sffamily\small, draw=none, fill=white, fill opacity=0.8},
  label style={font=\sffamily\small}, tick label style={font=\sffamily\small},
  grid=major, grid style={FpgaGray!20}, axis lines*=left,
]
  \addplot[very thick, FpgaRed, domain=1:64, samples=2] {2*x};
  \addlegendentry{последовательный перенос ($\sim n$)}
  \addplot[very thick, FpgaGreen, domain=1:64, samples=64] {4+4*ln(x)/ln(2)};
  \addlegendentry{ускоренный перенос ($\sim \log_2 n$)}
\end{axis}
\end{tikzpicture}
\caption{Задержка сумматора в зависимости от разрядности (оценка)}
\end{figure}

\section{Отрицательные числа и вычитание}

Как записать отрицательное число, если у нас есть только нули и единицы?
Используется \term{дополнительный код}: старший бит считается имеющим
\emph{отрицательный} вес. Для 4 битов веса разрядов равны $-8, 4, 2, 1$.

\begin{table}[H]
\centering
\caption{Четырёхбитные числа в дополнительном коде}
\small
\begin{tabular}{c c @{\hspace{10mm}} c c}
\toprule
\textbf{Код} & \textbf{Число} & \textbf{Код} & \textbf{Число} \\
\midrule
\code{0000} & 0 & \code{1000} & $-8$ \\
\code{0001} & 1 & \code{1001} & $-7$ \\
\code{0010} & 2 & \code{1010} & $-6$ \\
\code{0011} & 3 & \code{1011} & $-5$ \\
\code{0100} & 4 & \code{1100} & $-4$ \\
\code{0101} & 5 & \code{1101} & $-3$ \\
\code{0110} & 6 & \code{1110} & $-2$ \\
\code{0111} & 7 & \code{1111} & $-1$ \\
\bottomrule
\end{tabular}
\end{table}

\begin{figure}[H]
\centering
\begin{tikzpicture}
  \foreach \i in {0,...,15} {
    \pgfmathsetmacro{\ang}{90-\i*22.5}
    \pgfmathtruncatemacro{\v}{\i<8 ? \i : \i-16}
    \fill[\ifnum\i<8 FpgaBlue\else FpgaOrange\fi] (\ang:2.2) circle (0.09);
    \node[font=\sffamily\small] at (\ang:2.65) {\v};
  }
  \draw[thick, FpgaGray] (0,0) circle (2.2);
  \node[font=\ttfamily\scriptsize] at (90:1.75) {0000};
  \node[font=\ttfamily\scriptsize] at (90-7*22.5:1.75) {0111};
  \node[font=\ttfamily\scriptsize] at (90-8*22.5:1.75) {1000};
  \node[font=\ttfamily\scriptsize] at (90-15*22.5:1.75) {1111};
  \draw[-{Stealth}, thick, FpgaGreen] (60:1.2) arc (60:-30:1.2);
  \node[font=\sffamily\small, text=FpgaGreen] at (0,0) {$+1$};
  \node[font=\sffamily\small, align=left, anchor=west] at (3.6,0.8) {\textcolor{FpgaBlue}{$\bullet$} положительные: старший бит 0};
  \node[font=\sffamily\small, align=left, anchor=west] at (3.6,0.1) {\textcolor{FpgaOrange}{$\bullet$} отрицательные: старший бит 1};
  \node[font=\sffamily\small, align=left, anchor=west] at (3.6,-0.8) {$7 + 1 = -8$: \textbf{переполнение}!};
\end{tikzpicture}
\caption{Числа в дополнительном коде удобно представлять на круге: прибавление
единицы --- шаг по часовой стрелке}
\end{figure}

Главное достоинство дополнительного кода: \textbf{складывать и вычитать можно
одним и тем же сумматором}. Чтобы получить $-b$, нужно инвертировать все
биты $b$ и прибавить единицу:
\[
  a - b = a + \overline{b} + 1.
\]
Инверсию сделают элементы «Исключающее ИЛИ» (при управляющем сигнале 1
элемент $x\oplus 1 = \overline{x}$ инвертирует), а единицу --- вход переноса
младшего разряда.

\begin{figure}[H]
\centering
\begin{tikzpicture}
  \node[hblock, minimum width=40mm, minimum height=14mm] (add) at (0,0) {4-разрядный\\сумматор};
  \node[block, minimum width=22mm, minimum height=10mm] (xr) at (1.0,2.2) {4 элемента\\\scriptsize Исключающее ИЛИ};
  \draw[bus, -{Stealth}] (1.0,3.6) node[above, text=FpgaDark] {$B$} -- (xr.north);
  \draw[bus, -{Stealth}] (xr.south) -- node[right, lbl] {$B$ или $\overline{B}$} (1.0,0.7);
  \draw[bus, -{Stealth}] (-1.0,3.6) node[above, text=FpgaDark] {$A$} -- (-1.0,0.7);
  \draw[arr, FpgaRed] (-3.2,2.2) node[left, text=FpgaDark] {SUB} -- (xr.west);
  \draw[arr, FpgaRed] (-2.6,2.2) |- node[pos=0.75, above, lbl, text=FpgaRed] {$c_0$} (add.west);
  \fill[FpgaRed] (-2.6,2.2) circle (1.5pt);
  \draw[bus, -{Stealth}] (add.south) -- ++(0,-0.8) node[below, text=FpgaDark] {$S = A \pm B$};
\end{tikzpicture}
\caption{Сумматор-вычитатель: при SUB $= 0$ считает $A+B$, при SUB $= 1$ ---
$A-B$}
\end{figure}

\section{Компараторы}

\term{Компаратор} сравнивает два числа. Проще всего проверить
\emph{равенство}: числа равны, если равны все их разряды. Равенство двух
битов проверяет элемент «Исключающее ИЛИ-НЕ» ($\overline{a\oplus b} = 1$,
когда $a = b$), а элемент И объединяет результаты всех разрядов.

\begin{figure}[H]
\centering
\begin{tikzpicture}
  \foreach \i in {0,...,3} {
    \node[gXNOR, minimum width=10mm, minimum height=7mm] (x\i) at (1.5,1.8-\i*1.2) {};
    \draw[wire] (x\i.input 1) -- ++(-0.6,0) node[left] {$a_\i$};
    \draw[wire] (x\i.input 2) -- ++(-0.6,0) node[left] {$b_\i$};
  }
  \node[gAND, logic gate inputs=nnnn, minimum height=18mm] (g) at (4.2,0) {};
  \foreach \i [evaluate=\i as \j using int(\i+1)] in {0,...,3}
    \draw[wire] (x\i.output) -- ++(0.3,0) |- (g.input \j);
  \draw[wire] (g.output) -- ++(0.6,0) node[right] {$A = B$};
\end{tikzpicture}
\caption{Компаратор равенства 4-разрядных чисел}
\label{fig:cmp}
\end{figure}

Сравнение «больше/меньше» устроено сложнее: числа сравниваются, начиная со
старшего разряда. Первый разряд, в котором они различаются, определяет
результат. Для одного разряда:
\[
  (a > b) = a\,\overline{b}, \qquad (a < b) = \overline{a}\,b, \qquad (a = b) = \overline{a\oplus b}.
\]
Кстати, $A < B$ можно проверить и вычитателем: если $A - B$ отрицательно
(старший бит результата равен 1, при отсутствии переполнения), то $A < B$.

\section{Умножение и АЛУ}

Двоичное умножение ещё проще десятичного: каждая цифра множителя --- это 0
или 1, поэтому частичные произведения --- это либо ноль, либо сдвинутый
множимый. Их остаётся сложить.

\begin{figure}[H]
\centering
\begin{tikzpicture}[x=8mm, y=6.5mm, font=\ttfamily]
  \foreach \d [count=\i] in {1,0,1,1} \node at (\i+3,4) {\d};
  \foreach \d [count=\i] in {0,1,1,0} \node at (\i+3,3) {\d};
  \node at (2.8,3.5) {$\times$};
  \draw (3.3,2.5) -- (7.6,2.5);
  \foreach \d [count=\i] in {0,0,0,0} \node[text=FpgaGray] at (\i+3,2) {\d};
  \foreach \d [count=\i] in {1,0,1,1} \node[text=FpgaBlue] at (\i+2,1) {\d};
  \foreach \d [count=\i] in {1,0,1,1} \node[text=FpgaBlue] at (\i+1,0) {\d};
  \foreach \d [count=\i] in {0,0,0,0} \node[text=FpgaGray] at (\i,-1) {\d};
  \draw (0.5,-1.5) -- (7.6,-1.5);
  \foreach \d [count=\i] in {1,0,0,0,0,1,0} \node at (\i+0.0,-2) {\d};
  \node[font=\sffamily\small, anchor=west] at (8.5,4) {$1011_2 = 11$};
  \node[font=\sffamily\small, anchor=west] at (8.5,3) {$0110_2 = 6$};
  \node[font=\sffamily\small, anchor=west] at (8.5,-2) {$1000010_2 = 66$};
  \node[font=\sffamily\small, anchor=west, text=FpgaBlue] at (8.5,0.5) {сдвинутые копии};
\end{tikzpicture}
\caption{Умножение столбиком: сдвиги и сложения}
\end{figure}

Умножитель $n\times n$ бит требует примерно $n^2$ полных сумматоров, поэтому
он получается большим. В ПЛИС для этого есть готовые аппаратные умножители
(блоки DSP).

Сумматор, вычитатель, логические операции и сдвиги, собранные вместе и
выбираемые мультиплексором по коду операции, образуют \term{арифметико-логическое
устройство} (АЛУ) --- сердце любого процессора.

\begin{figure}[H]
\centering
\begin{tikzpicture}
  \draw[thick, fill=FpgaLight] (0,1.8) -- (0.9,1.8) -- (1.35,1.1) -- (1.8,1.8) -- (2.7,1.8) -- (2.0,-0.4) -- (0.7,-0.4) -- cycle;
  \node[font=\sffamily\bfseries] at (1.35,0.5) {АЛУ};
  \draw[bus, -{Stealth}] (0.45,2.6) node[above, text=FpgaDark] {$A$} -- (0.45,1.8);
  \draw[bus, -{Stealth}] (2.25,2.6) node[above, text=FpgaDark] {$B$} -- (2.25,1.8);
  \draw[bus, -{Stealth}] (1.35,-0.4) -- (1.35,-1.2) node[below, text=FpgaDark] {результат};
  \draw[arr, FpgaBlue] (-1.3,0.7) node[left, text=FpgaDark, align=right] {код\\операции} -- (0.33,0.7);
  \draw[arr] (2.37,0.7) -- (3.4,0.7) node[right, font=\sffamily\small, align=left] {флаги: ноль,\\перенос, знак,\\переполнение};
  \node[font=\sffamily\small, align=left, anchor=west] at (7.2,0.7) {\code{000}: $A+B$\\\code{001}: $A-B$\\\code{010}: $A \,\&\, B$\\\code{011}: $A \mid B$\\\code{100}: $A \oplus B$\ldots};
\end{tikzpicture}
\caption{Арифметико-логическое устройство}
\end{figure}

\begin{ideabox}
Сложение выполняется полными сумматорами, соединёнными цепочкой переноса.
Отрицательные числа записываются в дополнительном коде, и тогда вычитание
сводится к сложению. Компаратор равенства строится из элементов
«Исключающее ИЛИ-НЕ» и И.
\end{ideabox}

\chapter{Триггеры}
\label{ch:ff}

\section{Зачем схеме память}

Представьте кнопку «включить/выключить» на пульте: нажали --- свет
загорелся, отпустили --- продолжает гореть, нажали ещё раз --- погас. В момент,
когда кнопка отпущена, входы схемы одинаковы, а выход может быть и 0, и 1 ---
в зависимости от того, \emph{что было раньше}. Комбинационная схема так не
умеет. Нужна \term{память}.

Простейший элемент памяти хранит один бит и называется \term{триггером}
(от англ. \emph{trigger} --- «спусковой крючок»; в английской литературе ---
latch и flip-flop).

\section{Бистабильная ячейка}

Соединим два инвертора в кольцо. Если на выходе первого 1, то на выходе
второго 0, и этот 0 поддерживает 1 на выходе первого. Состояние
\emph{само себя поддерживает}. Возможно и противоположное устойчивое
состояние. У схемы два устойчивых состояния --- она \term{бистабильна} и
может хранить один бит.

\begin{figure}[H]
\centering
\begin{tikzpicture}
  \node[gNOT, minimum width=12mm, minimum height=8mm] (n1) at (0,0.8) {};
  \node[gNOT, minimum width=12mm, minimum height=8mm, rotate=180] (n2) at (0,-0.8) {};
  \draw[wire] (n1.output) -- ++(0.6,0) |- (n2.input);
  \draw[wire] (n2.output) -- ++(-0.6,0) |- (n1.input);
  \node[font=\sffamily\small, text=FpgaRed] at (1.4,1.15) {1};
  \node[font=\sffamily\small, text=FpgaBlue] at (-1.4,-0.45) {0};
  \fill ($(n1.output)+(0.6,-0.8)$) circle (1.3pt);
  \draw[wire] ($(n1.output)+(0.6,-0.8)$) -- ++(0.6,0) node[right] {$Q$};
  \node[font=\sffamily\small, align=left, anchor=west] at (3.4,0.3) {Состояние поддерживает само себя,};
  \node[font=\sffamily\small, align=left, anchor=west] at (3.4,-0.3) {но изменить его пока нечем.};
\end{tikzpicture}
\caption{Два инвертора в кольце --- бистабильная ячейка}
\end{figure}

\section{RS-триггер}

Чтобы записывать в ячейку новое значение, заменим инверторы на элементы
ИЛИ-НЕ и используем их вторые входы. Получится \term{RS-триггер}
(асинхронная RS-защёлка):
\begin{itemize}
  \item вход $S$ (Set --- «установить») записывает 1;
  \item вход $R$ (Reset --- «сбросить») записывает 0;
  \item если $S = R = 0$, триггер \textbf{хранит} прежнее значение.
\end{itemize}

\begin{figure}[H]
\centering
\begin{minipage}{0.48\textwidth}
\centering
\begin{tikzpicture}
  \node[gNOR, minimum height=10mm] (g1) at (0,1) {};
  \node[gNOR, minimum height=10mm] (g2) at (0,-1) {};
  \draw[wire] (g1.input 1) -- ++(-1.2,0) node[left] {$R$};
  \draw[wire] (g2.input 2) -- ++(-1.2,0) node[left] {$S$};
  \draw[wire] (g1.output) -- ++(1.2,0) node[right] {$Q$};
  \draw[wire] (g2.output) -- ++(1.2,0) node[right] {$\overline{Q}$};
  \draw[wire] ($(g1.output)+(0.4,0)$) -- ++(0,-0.6) -- ($(g2.input 1)+(-0.4,0.8)$) -- ($(g2.input 1)+(-0.4,0)$) -- (g2.input 1);
  \draw[wire] ($(g2.output)+(0.4,0)$) -- ++(0,0.6) -- ($(g1.input 2)+(-0.4,-0.8)$) -- ($(g1.input 2)+(-0.4,0)$) -- (g1.input 2);
  \fill ($(g1.output)+(0.4,0)$) circle (1.3pt); \fill ($(g2.output)+(0.4,0)$) circle (1.3pt);
\end{tikzpicture}
\end{minipage}\hfill
\begin{minipage}{0.48\textwidth}
\centering
\begin{tabular}{cc|c|l}
\toprule
$S$ & $R$ & $Q$ & \textbf{Режим} \\
\midrule
0 & 0 & $Q$ & хранение \\
1 & 0 & 1 & установка \\
0 & 1 & 0 & сброс \\
1 & 1 & \textcolor{FpgaRed}{?} & \textcolor{FpgaRed}{запрещено} \\
\bottomrule
\end{tabular}
\end{minipage}
\caption{RS-триггер на двух элементах ИЛИ-НЕ и его таблица переходов}
\label{fig:rs}
\end{figure}

Разберём, почему он работает. Помним: выход ИЛИ-НЕ равен 0, если хотя бы один
вход равен 1.
\begin{itemize}
  \item $S=1$: нижний элемент выдаёт $\overline{Q}=0$. У верхнего оба входа
        равны нулю ($R=0$, $\overline{Q}=0$), поэтому $Q=1$.
  \item Теперь $S$ вернули в 0. На нижнем элементе $Q=1$, поэтому
        $\overline{Q}$ по-прежнему 0, а $Q$ по-прежнему 1: триггер
        \emph{запомнил} единицу.
  \item $S=R=1$ --- оба выхода становятся нулями (хотя должны быть
        противоположны), а при одновременном снятии сигналов неизвестно, в
        какое состояние триггер «упадёт». Поэтому это сочетание запрещено.
\end{itemize}

\begin{figure}[H]
\centering
\begin{tikztimingtable}[timing/slope=0.1, xscale=3.0, yscale=1.4, timing/font=\sffamily\small, timing/rowdist=3.2ex]
  $S$ & 1L 1H 10L \\
  $R$ & 5L 1H 6L \\
  $Q$ & 1L 4H 7L \\
\extracode
  \node[lbl, text=FpgaOrange] at (1.5,1.1) {уст.};
  \node[lbl, text=FpgaOrange] at (3.5,1.1) {хранение 1};
  \node[lbl, text=FpgaOrange] at (5.5,1.1) {сброс};
  \node[lbl, text=FpgaOrange] at (8.0,1.1) {хранение 0};
\end{tikztimingtable}
\caption{Работа RS-триггера: короткий импульс на $S$ или $R$ переключает
триггер, а без импульсов он хранит значение}
\end{figure}

\begin{infobox}[RS-триггер на элементах И-НЕ]
Точно так же триггер собирают из двух элементов И-НЕ. Но тогда входы
срабатывают от \emph{нуля}: $\overline{S}=0$ устанавливает, $\overline{R}=0$
сбрасывает, а запрещено сочетание $\overline{S}=\overline{R}=0$. Такой
триггер --- классическое средство подавления дребезга механического
переключателя: при первом же касании контакта он переключается и дальнейших
«прыжков» не замечает.
\end{infobox}

\section{D-защёлка}

У RS-триггера два неудобства: есть запрещённое сочетание входов, и он
реагирует на входы в любой момент. Добавим перед ним два элемента И, вход
данных $D$ и вход разрешения $C$ (clock, «синхронизация»). Получится
\term{D-защёлка} (D latch):
\begin{itemize}
  \item при $C=1$ защёлка \textbf{прозрачна}: выход повторяет вход, $Q = D$;
  \item при $C=0$ защёлка \textbf{хранит} значение, которое было на $D$ в момент
        перехода $C$ из 1 в 0.
\end{itemize}

\begin{figure}[H]
\centering
\begin{tikzpicture}
  \node[gAND, minimum height=9mm] (a1) at (1.6,1) {};
  \node[gAND, minimum height=9mm] (a2) at (1.6,-1) {};
  \node[draw=FpgaGreen, fill=green!8, thick, minimum width=18mm, minimum height=30mm, font=\sffamily\small, align=center] (rs) at (4.4,0) {RS-\\триггер};
  \node[font=\sffamily\scriptsize] at ($(rs.west)+(0.25,1)$) {R};
  \node[font=\sffamily\scriptsize] at ($(rs.west)+(0.25,-1)$) {S};
  \node[gNOT, minimum width=7mm, minimum height=5mm] (n) at (-0.2,1.12) {};
  \draw[wire] (-1.6,-0.88) node[left] {$D$} -- (a2.input 1);
  \draw[wire] (-0.9,-0.88) |- (n.input); \fill (-0.9,-0.88) circle (1.3pt);
  \draw[wire] (n.output) -- (a1.input 1);
  \draw[wire, FpgaBlue] (-1.6,0) node[left, text=FpgaDark] {$C$} -- (0.4,0);
  \draw[wire, FpgaBlue] (0.4,0) |- (a1.input 2); \draw[wire, FpgaBlue] (0.4,0) |- (a2.input 2);
  \fill[FpgaBlue] (0.4,0) circle (1.3pt);
  \draw[wire] (a1.output) -- ($(rs.west)+(0,1)$);
  \draw[wire] (a2.output) -- ($(rs.west)+(0,-1)$);
  \draw[wire] ($(rs.east)+(0,1)$) -- ++(0.8,0) node[right] {$Q$};
  \draw[wire] ($(rs.east)+(0,-1)$) -- ++(0.8,0) node[right] {$\overline{Q}$};
\end{tikzpicture}
\caption{D-защёлка: запрещённое сочетание $S=R=1$ теперь невозможно, так как
на $S$ и $R$ подаются $D$ и $\overline{D}$}
\end{figure}

\begin{figure}[H]
\centering
\begin{tikztimingtable}[timing/slope=0.08, xscale=1.6, yscale=1.4, timing/font=\sffamily\small, timing/rowdist=3.2ex]
  $C$ & 2L 4H 4L 4H 2L \\
  $D$ & 1L 2H 1L 4H 3L 1H 4L \\
  $Q$ & 2L 1H 1L 6H 1L 1H 4L \\
\extracode
  \begin{pgfonlayer}{background}
    \fill[green!10] (2,0.6) rectangle (6,-5.2);
    \fill[green!10] (10,0.6) rectangle (14,-5.2);
  \end{pgfonlayer}
  \node[lbl, text=FpgaGreen] at (4,1.0) {прозрачна};
  \node[lbl] at (8,1.0) {хранит};
  \node[lbl, text=FpgaGreen] at (12,1.0) {прозрачна};
\end{tikztimingtable}
\caption{D-защёлка при $C=1$ передаёт все изменения $D$ на выход, при $C=0$
«замораживает» выход}
\label{fig:dlatch_t}
\end{figure}

\section{D-триггер}

Прозрачность защёлки --- это и достоинство, и недостаток. Если выход защёлки
через комбинационную логику идёт обратно на её же вход (как в счётчике), то
пока $C=1$, сигнал будет бегать по кругу, и результат станет непредсказуемым.
Нужен элемент, который запоминает значение в \emph{один момент времени} ---
по \emph{фронту} тактового сигнала. Это \term{D-триггер} (D flip-flop),
главный элемент памяти всей современной цифровой техники и ПЛИС.

Его можно собрать из двух D-защёлок, включённых последовательно, с
противоположными сигналами разрешения. Такая схема называется
\term{«ведущий -- ведомый»} (master--slave):

\begin{figure}[H]
\centering
\begin{tikzpicture}[lt/.style={draw=FpgaGreen, fill=green!8, thick, minimum width=24mm, minimum height=16mm, font=\sffamily\small, align=center}]
  \node[lt] (m) at (0,0) {\quad ведущая\\\quad защёлка};
  \node[lt] (s) at (4.2,0) {\quad ведомая\\\quad защёлка};
  \node[font=\sffamily\scriptsize] at ($(m.west)+(0.2,0.45)$) {D};
  \node[font=\sffamily\scriptsize] at ($(m.west)+(0.2,-0.45)$) {C};
  \node[font=\sffamily\scriptsize] at ($(s.west)+(0.2,0.45)$) {D};
  \node[font=\sffamily\scriptsize] at ($(s.west)+(0.2,-0.45)$) {C};
  \draw[wire] ($(m.west)+(-1.2,0.45)$) node[left] {$D$} -- ($(m.west)+(0,0.45)$);
  \draw[wire] ($(m.east)+(0,0.45)$) -- node[above, lbl] {$Q_m$} ($(s.west)+(0,0.45)$);
  \draw[wire] ($(s.east)+(0,0.45)$) -- ++(1,0) node[right] {$Q$};
  \node[gNOT, minimum width=7mm, minimum height=5mm] (n) at (-1.2,-1.3) {};
  \draw[wire, FpgaBlue] (-2.4,-1.3) node[left, text=FpgaDark] {CLK} -- (n.input);
  \draw[wire, FpgaBlue] (n.output) -- ++(0.2,0) |- ($(m.west)+(0,-0.45)$);
  \draw[wire, FpgaBlue] (-1.9,-1.3) |- (2.5,-1.9) |- ($(s.west)+(0,-0.45)$);
  \fill[FpgaBlue] (-1.9,-1.3) circle (1.3pt);
  \node[font=\sffamily\small, align=left, anchor=north west] at (-2.4,-2.2) {CLK $=0$: ведущая прозрачна, ведомая хранит;\\
  CLK $=1$: ведущая хранит, ведомая прозрачна;\\ $\Rightarrow$ $Q$ меняется только в момент перехода CLK из 0 в 1.};
\end{tikzpicture}
\caption{D-триггер, построенный по схеме «ведущий -- ведомый»}
\label{fig:ms}
\end{figure}

На схемах D-триггер рисуют прямоугольником, где вход тактового сигнала
отмечен треугольником (знак «срабатывает по фронту»):

\begin{figure}[H]
\centering
\begin{tikzpicture}
  \draw[thick, fill=green!8] (0,0) rectangle (1.8,2.2);
  \node[font=\sffamily\small] at (0.3,1.7) {D}; \node[font=\sffamily\small] at (1.5,1.7) {Q};
  \node[font=\sffamily\small] at (1.5,0.5) {$\overline{\text{Q}}$};
  \draw[thick] (0,0.3) -- (0.25,0.5) -- (0,0.7);
  \draw[wire] (-0.7,1.7) node[left] {$D$} -- (0,1.7);
  \draw[wire] (-0.7,0.5) node[left] {CLK} -- (0,0.5);
  \draw[wire] (1.8,1.7) -- (2.5,1.7) node[right] {$Q$};
  \draw[wire] (1.8,0.5) -- (2.1,0.5); \draw[thick] (2.18,0.5) circle (0.08); \draw[wire] (2.26,0.5) -- (2.5,0.5) node[right] {$\overline{Q}$};
  \draw[wire] (0.9,2.2) -- (0.9,2.7) node[above] {$\overline{\text{SET}}$}; \draw[thick, fill=white] (0.9,2.28) circle (0.08);
  \draw[wire] (0.9,0) -- (0.9,-0.5) node[below] {$\overline{\text{RST}}$}; \draw[thick, fill=white] (0.9,-0.08) circle (0.08);
  \node[font=\sffamily\small, align=left, anchor=west] at (4.2,1.6) {треугольник --- срабатывает по фронту};
  \node[font=\sffamily\small, align=left, anchor=west] at (4.2,1.0) {кружок --- активный уровень 0};
  \node[font=\sffamily\small, align=left, anchor=west] at (4.2,0.4) {SET, RST --- асинхронные установка и сброс:};
  \node[font=\sffamily\small, align=left, anchor=west] at (4.2,-0.1) {срабатывают сразу, не дожидаясь фронта};
\end{tikzpicture}
\caption{Условное обозначение D-триггера}
\end{figure}

\begin{figure}[H]
\centering
\begin{tikztimingtable}[timing/slope=0.08, xscale=1.6, yscale=1.4, timing/font=\sffamily\small, timing/rowdist=3.2ex]
  CLK            & 2L 4H 4L 4H 2L \\
  $D$            & 1L 2H 1L 4H 3L 1H 4L \\
  $Q$ защёлки    & 2L 1H 1L 6H 1L 1H 4L \\
  $Q$ триггера   & 2L 8H 6L \\
\extracode
  \begin{pgfonlayer}{background}
    \draw[FpgaOrange, dashed] (2,0.6) -- (2,-7.2);
    \draw[FpgaOrange, dashed] (10,0.6) -- (10,-7.2);
  \end{pgfonlayer}
  \node[lbl, text=FpgaOrange] at (2,1.0) {фронт}; \node[lbl, text=FpgaOrange] at (10,1.0) {фронт};
\end{tikztimingtable}
\caption{Сравнение: защёлка (рис.~\ref{fig:dlatch_t}) пропускает все изменения
$D$, пока CLK $=1$, а триггер берёт значение $D$ \emph{только в момент фронта}}
\label{fig:latch_vs_ff}
\end{figure}

\section{Временные параметры триггера}

Чтобы триггер надёжно запомнил значение, вход $D$ должен быть неизменным в
течение некоторого времени до фронта и после него:
\begin{itemize}
  \item $t_\text{su}$ (setup, \term{время предустановки}) --- сколько $D$ должен
        быть стабилен \emph{до} фронта;
  \item $t_\text{h}$ (hold, \term{время удержания}) --- сколько \emph{после} фронта;
  \item $t_\text{co}$ (clock-to-output) --- через сколько после фронта меняется $Q$.
\end{itemize}

\begin{figure}[H]
\centering
\begin{tikzpicture}[x=1cm, y=1cm]
  \draw[very thick, FpgaBlue] (0,2.2) -- (4,2.2) -- (4.1,3) -- (8,3);
  \node[font=\sffamily\small, anchor=east] at (-0.1,2.6) {CLK};
  \draw[very thick] (0,1.1) -- (1.5,1.1) -- (1.7,1.5) -- (6,1.5) -- (6.2,1.1) -- (8,1.1);
  \draw[very thick] (0,1.5) -- (1.5,1.5) -- (1.7,1.1) -- (6,1.1) -- (6.2,1.5) -- (8,1.5);
  \node[font=\sffamily\small, anchor=east] at (-0.1,1.3) {$D$};
  \node[font=\sffamily\scriptsize] at (3.9,1.3) {стабильно};
  \draw[very thick, FpgaGreen] (0,0) -- (4.8,0) -- (4.9,0.4) -- (8,0.4);
  \node[font=\sffamily\small, anchor=east] at (-0.1,0.2) {$Q$};
  \fill[FpgaRed, opacity=0.15] (2.6,0.9) rectangle (4.05,1.7);
  \fill[FpgaOrange, opacity=0.2] (4.05,0.9) rectangle (4.8,1.7);
  \draw[dashed, FpgaGray] (4.05,3.3) -- (4.05,-0.5);
  \draw[<->] (2.6,-0.35) -- (4.05,-0.35) node[midway, below, font=\sffamily\small] {$t_\text{su}$};
  \draw[<->] (4.05,-0.35) -- (4.8,-0.35) node[midway, below, font=\sffamily\small] {$t_\text{h}$};
  \draw[<->] (4.05,-1.1) -- (4.85,-1.1) node[midway, below, font=\sffamily\small] {$t_\text{co}$};
\end{tikzpicture}
\caption{Время предустановки, удержания и задержка выхода триггера}
\label{fig:setup_hold}
\end{figure}

\begin{warnbox}[Метастабильность]
Если $D$ меняется внутри окна $t_\text{su}$--$t_\text{h}$, триггер может на
некоторое время «зависнуть» в промежуточном состоянии --- ни 0, ни 1 --- и
лишь потом случайно упасть в одно из них. Это \term{метастабильность}. Она
неизбежно возникает, когда на триггер приходит сигнал, не связанный с
тактом (например, от кнопки). Защита --- \term{синхронизатор} из двух
триггеров подряд: к моменту следующего фронта первый триггер почти наверняка
успокоится.
\end{warnbox}

\section{Вход разрешения}

Часто нужно, чтобы триггер обновлялся не на каждом фронте, а только по
команде. Для этого добавляют вход разрешения CE (clock enable). Устроен он
просто: перед входом $D$ ставится мультиплексор, который при CE $=0$ подаёт
на вход триггера его же выход --- триггер переписывает сам себя и не меняется.

\begin{figure}[H]
\centering
\begin{tikzpicture}
  \draw[thick, fill=orange!15] (0,0.9) -- (0.9,0.5) -- (0.9,-0.5) -- (0,-0.9) -- cycle;
  \node[font=\sffamily\scriptsize] at (0.45,0) {MUX};
  \node[font=\tiny, anchor=west] at (0.02,0.5) {0}; \node[font=\tiny, anchor=west] at (0.02,-0.5) {1};
  \draw[thick, fill=green!8] (1.8,-0.8) rectangle (3.2,0.8);
  \node[font=\sffamily\small] at (2.05,0) {D}; \node[font=\sffamily\small] at (2.95,0) {Q};
  \draw[thick] (2.3,-0.8) -- (2.5,-0.6) -- (2.7,-0.8);
  \draw[wire] (0.9,0) -- (1.8,0);
  \draw[wire] (3.2,0) -- (4.4,0) node[right] {$Q$};
  \draw[wire] (3.8,0) -- (3.8,1.5) -| (-0.6,0.5) -- (0,0.5); \fill (3.8,0) circle (1.3pt);
  \draw[wire] (-1.2,-0.5) node[left] {$D$} -- (0,-0.5);
  \draw[wire, FpgaBlue] (0.45,-0.7) -- (0.45,-1.4) node[below, text=FpgaDark] {CE};
  \draw[wire, FpgaBlue] (2.5,-0.8) -- (2.5,-1.4) node[below, text=FpgaDark] {CLK};
\end{tikzpicture}
\caption{D-триггер с входом разрешения: мультиплексор + обратная связь}
\label{fig:ce}
\end{figure}

\section{T-триггер и JK-триггер}

\term{T-триггер} (toggle, «переключатель») меняет своё состояние на
противоположное по каждому фронту, если на входе $T=1$, и хранит его при
$T=0$. Его легко сделать из D-триггера и элемента «Исключающее ИЛИ»:
$D = T \oplus Q$.

\begin{figure}[H]
\centering
\begin{minipage}{0.5\textwidth}
\centering
\begin{tikzpicture}
  \node[gXOR, minimum width=10mm, minimum height=8mm] (x) at (0,0) {};
  \draw[thick, fill=green!8] (1.2,-0.8) rectangle (2.6,0.8);
  \node[font=\sffamily\small] at (1.45,0) {D}; \node[font=\sffamily\small] at (2.35,0) {Q};
  \draw[thick] (1.7,-0.8) -- (1.9,-0.6) -- (2.1,-0.8);
  \draw[wire] (x.output) -- (1.2,0);
  \draw[wire] (2.6,0) -- (3.6,0) node[right] {$Q$};
  \draw[wire] (3.1,0) -- (3.1,1.3) -| (-1.0,0 |- x.input 1) -- (x.input 1); \fill (3.1,0) circle (1.3pt);
  \draw[wire] (x.input 2) -- ++(-0.9,0) node[left] {$T$};
  \draw[wire, FpgaBlue] (1.9,-0.8) -- (1.9,-1.3) node[below, text=FpgaDark] {CLK};
\end{tikzpicture}
\end{minipage}\hfill
\begin{minipage}{0.45\textwidth}
\centering
\begin{tabular}{c|c|l}
\toprule
$T$ & $Q$ после фронта & \\
\midrule
0 & $Q$ & хранение \\
1 & $\overline{Q}$ & переключение \\
\bottomrule
\end{tabular}
\end{minipage}
\caption{T-триггер из D-триггера и элемента «Исключающее ИЛИ»}
\end{figure}

Если на T-триггер постоянно подавать $T=1$, он переключается на каждом
фронте, и частота на его выходе будет \emph{вдвое меньше} тактовой. Это
простейший \term{делитель частоты на 2}.

\begin{figure}[H]
\centering
\begin{tikztimingtable}[timing/slope=0.08, xscale=1.5, yscale=1.4, timing/font=\sffamily\small, timing/rowdist=3.2ex]
  CLK            & 8{1L 1H} \\
  $Q$ ($T=1$)    & 1L 2H 2L 2H 2L 2H 2L 2H 1L \\
\end{tikztimingtable}
\caption{T-триггер с $T=1$ делит частоту на 2}
\end{figure}

\term{JK-триггер} объединяет возможности RS- и T-триггеров: $J$ ---
установка, $K$ --- сброс, $J = K = 1$ --- переключение (запрещённого сочетания
нет). В старых схемах на микросхемах серии 155 JK-триггеры были очень
популярны. В ПЛИС же используются только D-триггеры, а остальные виды
при необходимости получаются из них добавлением логики на входе.

\begin{table}[H]
\centering
\caption{Сводка: виды триггеров}
\small
\begin{tabularx}{\textwidth}{l l X}
\toprule
\textbf{Тип} & \textbf{Когда срабатывает} & \textbf{Что делает} \\
\midrule
RS-защёлка & сразу & $S$ --- установить 1, $R$ --- сбросить в 0, $S=R=1$ запрещено \\
D-защёлка  & пока $C = 1$ & повторяет $D$, при $C = 0$ хранит \\
\textbf{D-триггер} & \textbf{по фронту} CLK & запоминает $D$ в момент фронта \\
T-триггер  & по фронту CLK & при $T=1$ меняет состояние на противоположное \\
JK-триггер & по фронту CLK & $J$ --- установка, $K$ --- сброс, $J=K=1$ --- переключение \\
\bottomrule
\end{tabularx}
\end{table}

\begin{ideabox}
Триггер хранит один бит. Главный элемент памяти в цифровой технике ---
\textbf{D-триггер}: он запоминает значение входа $D$ в момент \emph{переднего
фронта} тактового сигнала и хранит его до следующего фронта. Чтобы
триггер работал надёжно, вход не должен меняться в окне $t_\text{su}$--$t_\text{h}$
вокруг фронта.
\end{ideabox}

\chapter{Регистры и счётчики}
\label{ch:regs}

Один триггер хранит один бит. Соединив несколько триггеров, получают узлы,
которые хранят числа, сдвигают их и считают.

\newcommand{\dff}[3]{% имя, x, y  -- D-триггер 12x14 мм
  \draw[thick, fill=green!8] (#2,#3) rectangle ++(1.2,1.4);
  \node[font=\sffamily\scriptsize] at (#2+0.22,#3+1.05) {D};
  \node[font=\sffamily\scriptsize] at (#2+0.98,#3+1.05) {Q};
  \draw[thick] (#2,#3+0.2) -- (#2+0.18,#3+0.35) -- (#2,#3+0.5);
  \coordinate (#1-d) at (#2,#3+1.05); \coordinate (#1-q) at (#2+1.2,#3+1.05);
  \coordinate (#1-c) at (#2,#3+0.35);}

\section{Регистр хранения}

\term{Регистр} --- это группа D-триггеров с общим тактовым сигналом. По
фронту все триггеры одновременно запоминают свои входы, и регистр хранит
$n$-разрядное число до следующего фронта. На схемах регистр обозначают
буквами RG.

\begin{figure}[H]
\centering
\begin{tikzpicture}
  \foreach \i in {0,...,3} {
    \pgfmathsetmacro{\x}{(3-\i)*2.3}
    \dff{r\i}{\x}{0}
    \draw[wire] (r\i-d) -- ++(-0.4,0) |- ($(r\i-d)+(-0.4,1.0)$) node[above] {$d_\i$};
    \draw[wire] (r\i-q) -- ++(0.35,0) -- ++(0,-1.5) node[below] {$q_\i$};
    \draw[wire, FpgaBlue] (r\i-c) -- ++(-0.3,0) -- ++(0,-0.75);
  }
  \draw[wire, FpgaBlue] (-1.4,-0.4) node[left, text=FpgaDark] {CLK} -- (6.6,-0.4);
  \foreach \i in {0,...,3} { \pgfmathsetmacro{\x}{(3-\i)*2.3-0.3} \fill[FpgaBlue] (\x,-0.4) circle (1.3pt); }
\end{tikzpicture}
\caption{Четырёхразрядный регистр: четыре D-триггера с общим тактом}
\label{fig:reg}
\end{figure}

\begin{figure}[H]
\centering
\begin{tikztimingtable}[timing/slope=0.08, xscale=1.5, yscale=1.4, timing/font=\sffamily\small, timing/rowdist=3.2ex]
  CLK      & 8{1L 1H} \\
  $D[3:0]$ & 2D{5} 1.5D{9} 1D{7} 1.5D{3} 4D{C} 6D{0} \\
  $Q[3:0]$ & 1D{?} 2D{5} 2D{9} 2D{3} 4D{C} 5D{0} \\
\end{tikztimingtable}
\caption{Регистр «фотографирует» значение шины $D$ в моменты передних фронтов
CLK. Значение 7 продержалось на входе меньше такта и пришлось между
фронтами --- регистр его просто не заметил}
\end{figure}

Регистры --- это «переменные» аппаратуры. В процессоре в регистрах хранятся
операнды и результаты вычислений; в любом устройстве на ПЛИС --- текущее
состояние, счётчики, данные между этапами обработки. Чтобы регистр менял
значение не на каждом фронте, а только по команде, используют триггеры с
входом разрешения CE (рис.~\ref{fig:ce}).

\section{Сдвиговый регистр}

Соединим триггеры цепочкой: выход каждого подаётся на вход следующего. По
каждому фронту содержимое \term{сдвигается} на один разряд, а в первый
триггер записывается новый бит с входа $D_\text{in}$.

\begin{figure}[H]
\centering
\begin{tikzpicture}
  \foreach \i in {0,...,3} {
    \pgfmathsetmacro{\x}{\i*2.3}
    \dff{s\i}{\x}{0}
    \draw[wire, FpgaBlue] (s\i-c) -- ++(-0.3,0) -- ++(0,-0.75);
  }
  \draw[wire] (-1.0,1.05) node[left] {$D_\text{in}$} -- (s0-d);
  \foreach \i [evaluate=\i as \j using int(\i+1)] in {0,1,2} {
    \draw[wire] (s\i-q) -- (s\j-d);
    \fill ($(s\i-q)+(0.3,0)$) circle (1.3pt);
    \draw[wire] ($(s\i-q)+(0.3,0)$) -- ++(0,1.0) node[above] {$q_\i$};
  }
  \draw[wire] (s3-q) -- ++(0.7,0) node[right] {$q_3 = D_\text{out}$};
  \draw[wire, FpgaBlue] (-1.0,-0.4) node[left, text=FpgaDark] {CLK} -- (6.6,-0.4);
\end{tikzpicture}
\caption{Четырёхразрядный сдвиговый регистр}
\label{fig:shift}
\end{figure}

\begin{figure}[H]
\centering
\begin{tikztimingtable}[timing/slope=0.08, xscale=1.5, yscale=1.4, timing/font=\sffamily\small, timing/rowdist=3.2ex]
  CLK            & 8{1L 1H} \\
  $D_\text{in}$  & 2H 2L 12L \\
  $q_0$          & 1L 2H 13L \\
  $q_1$          & 3L 2H 11L \\
  $q_2$          & 5L 2H 9L \\
  $q_3$          & 7L 2H 7L \\
\end{tikztimingtable}
\caption{Единица, поданная на вход, продвигается по регистру на один разряд
за такт}
\end{figure}

\subsection{Зачем сдвигать}

\begin{description}
  \item[Последовательный $\to$ параллельный код.] Биты приходят по одному
    проводу друг за другом (так работают интерфейсы UART, SPI, I$^2$C). После
    $n$ тактов все $n$ битов оказываются в регистре одновременно и могут быть
    прочитаны как число.
  \item[Параллельный $\to$ последовательный код.] Обратная задача при передаче:
    число загружается в регистр целиком, а затем «выталкивается» по одному
    биту за такт. Для загрузки перед каждым триггером ставят мультиплексор
    «сдвиг/загрузка».
  \item[Задержка сигнала] на заданное число тактов.
  \item[Умножение и деление на 2]: сдвиг двоичного числа влево на один разряд
    умножает его на 2, вправо --- делит.
\end{description}

\begin{figure}[H]
\centering
\begin{tikzpicture}[bit/.style={draw=FpgaBlue, fill=FpgaLight, minimum size=8mm, font=\ttfamily}]
  \node[font=\sffamily\small, anchor=east] at (-0.6,0) {передатчик};
  \foreach \b [count=\i from 0] in {1,0,1,1} \node[bit] at (\i*0.8,0) {\b};
  \draw[very thick, FpgaOrange, -{Stealth}] (3.0,0) -- node[above, font=\sffamily\small, text=FpgaDark] {1 провод: \code{1}, \code{1}, \code{0}, \code{1}\ldots} (7.4,0);
  \foreach \b [count=\i from 0] in {1,0,1,1} \node[bit] at (7.8+\i*0.8,0) {\b};
  \node[font=\sffamily\small, anchor=west] at (10.8,0) {приёмник};
  \node[lbl] at (1.2,-0.8) {параллельный $\to$ послед.};
  \node[lbl] at (9.0,-0.8) {послед. $\to$ параллельный};
\end{tikzpicture}
\caption{Передача числа по одному проводу с помощью двух сдвиговых регистров}
\end{figure}

\subsection{Кольцевой регистр и генератор псевдослучайных чисел}

Если выход последнего триггера соединить со входом первого, получится
\term{кольцевой регистр}: записанная в него единица будет бесконечно бегать по
кругу (эффект «бегущего огня» в главе~\ref{ch:examples}).

А если на вход подать «Исключающее ИЛИ» от нескольких разрядов, получится
\term{регистр сдвига с линейной обратной связью} (LFSR). При правильном выборе
отводов он перебирает все $2^n-1$ ненулевых состояний в «перемешанном»
порядке. Такие регистры используются как простые генераторы
псевдослучайных чисел, в помехоустойчивом кодировании и для шифрования.

\begin{figure}[H]
\centering
\begin{minipage}{0.55\textwidth}
\centering
\begin{tikzpicture}
  \foreach \i in {0,...,2} {
    \pgfmathsetmacro{\x}{\i*2.0}
    \dff{l\i}{\x}{0}
  }
  \draw[wire] (l0-q) -- (l1-d); \draw[wire] (l1-q) -- (l2-d);
  \node[gXOR, minimum width=9mm, minimum height=7mm, rotate=180] (x) at (2.6,2.3) {};
  \draw[wire] (l2-q) -- ++(0.4,0) |- (x.input 1);
  \fill ($(l1-q)+(0.3,0)$) circle (1.3pt);
  \draw[wire] ($(l1-q)+(0.3,0)$) |- (x.input 2);
  \draw[wire] (x.output) -| ($(l0-d)+(-0.4,0)$) -- (l0-d);
  \foreach \i in {0,...,2} \node[lbl, below=1mm] at ($(l\i-c)+(0.6,-0.4)$) {$q_\i$};
\end{tikzpicture}
\end{minipage}\hfill
\begin{minipage}{0.4\textwidth}
\centering
\small
\begin{tabular}{c|ccc}
\toprule
такт & $q_0$ & $q_1$ & $q_2$ \\
\midrule
0 & 1 & 0 & 0 \\
1 & 0 & 1 & 0 \\
2 & 1 & 0 & 1 \\
3 & 1 & 1 & 0 \\
4 & 1 & 1 & 1 \\
5 & 0 & 1 & 1 \\
6 & 0 & 0 & 1 \\
7 & 1 & 0 & 0 \\
\bottomrule
\end{tabular}
\end{minipage}
\caption{Трёхразрядный LFSR ($q_0 \leftarrow q_1\oplus q_2$) перебирает все семь
ненулевых состояний и повторяется}
\end{figure}

\section{Счётчики}

\term{Счётчик} --- регистр, значение которого по каждому такту
увеличивается (или уменьшается) на единицу. Счётчики измеряют время, делят
частоту, перебирают адреса памяти.

\subsection{Асинхронный счётчик}

Посмотрите, как меняются разряды при счёте $0, 1, 2, \ldots, 7$: младший бит
переключается каждый шаг, следующий --- в два раза реже, следующий --- ещё в
два раза реже. Каждый разряд переключается, когда предыдущий переходит из 1
в 0. Значит, можно соединить T-триггеры цепочкой: тактовый вход каждого
следующего триггера подключить к \emph{инверсному} выходу предыдущего.

\begin{figure}[H]
\centering
\begin{tikzpicture}
  \foreach \i in {0,...,2} {
    \pgfmathsetmacro{\x}{\i*2.8}
    \draw[thick, fill=green!8] (\x,0) rectangle ++(1.2,1.4);
    \node[font=\sffamily\scriptsize] at (\x+0.22,1.05) {T};
    \node[font=\sffamily\scriptsize] at (\x+0.98,1.05) {Q};
    \node[font=\sffamily\scriptsize] at (\x+0.98,0.35) {$\overline{\text{Q}}$};
    \draw[thick] (\x,0.2) -- (\x+0.18,0.35) -- (\x,0.5);
    \draw[wire] (\x,1.05) -- ++(-0.3,0) node[left, font=\scriptsize] {1};
    \draw[wire] (\x+1.2,1.05) -- ++(0.4,0) -- ++(0,0.8) node[above] {$q_\i$};
  }
  \draw[wire, FpgaBlue] (-1.0,0.35) node[left, text=FpgaDark] {CLK} -- (0,0.35);
  \draw[wire, FpgaRed] (1.2,0.35) -- (2.8,0.35);
  \draw[wire, FpgaRed] (4.0,0.35) -- (5.6,0.35);
\end{tikzpicture}
\caption{Асинхронный (последовательный) трёхразрядный счётчик на T-триггерах}
\label{fig:ripple}
\end{figure}

\begin{figure}[H]
\centering
\begin{tikztimingtable}[timing/slope=0.05, xscale=1.2, yscale=1.4, timing/font=\sffamily\small, timing/rowdist=3.2ex]
  CLK     & 16{1L 1H} \\
  $q_0$   & 1L 7{2H 2L} 2H 1L \\
  $q_1$   & 3L 3{4H 4L} 4H 1L \\
  $q_2$   & 7L 8H 8L 8H 1L \\
  число   & 1D{0} 2D{1} 2D{2} 2D{3} 2D{4} 2D{5} 2D{6} 2D{7} 2D{0} 2D{1} 2D{2} 2D{3} 2D{4} 2D{5} 2D{6} 2D{7} 1D{0} \\
\end{tikztimingtable}
\caption{Каждый следующий разряд меняется вдвое реже: счётчик перебирает числа
$0\ldots7$, а разряд $q_k$ имеет частоту $f_\text{CLK}/2^{k+1}$}
\label{fig:count_t}
\end{figure}

Недостаток асинхронного счётчика --- триггеры переключаются не одновременно,
а «волной», каждый с задержкой относительно предыдущего. В момент перехода
$011 \to 100$ на выходе на короткое время появляются ложные значения
($011\to010\to000\to100$). В длинном счётчике эта «волна» идёт долго.

\subsection{Синхронный счётчик}

В синхронном счётчике все триггеры тактируются \emph{одним} сигналом, а новое
значение вычисляет комбинационная схема --- сумматор, прибавляющий единицу.
Все разряды меняются одновременно по фронту. Именно так строятся все
счётчики в ПЛИС.

\begin{figure}[H]
\centering
\begin{tikzpicture}
  \node[draw=FpgaGreen, fill=green!8, thick, minimum width=20mm, minimum height=18mm, font=\sffamily\small, align=center] (reg) at (4.4,0) {регистр\\\scriptsize $n$ бит};
  \node[font=\sffamily\scriptsize, anchor=west] at (reg.west) {D};
  \node[font=\sffamily\scriptsize, anchor=east] at (reg.east) {Q};
  \draw[thick] ([xshift=-1.2mm]reg.south) -- ([yshift=1.8mm]reg.south) -- ([xshift=1.2mm]reg.south);
  \draw[wire, FpgaBlue] (reg.south) -- ++(0,-0.8) node[below, text=FpgaDark] {CLK};
  \node[draw=FpgaBlue, fill=FpgaLight, thick, circle, minimum size=12mm, font=\sffamily\bfseries] (add) at (1.4,0) {$+1$};
  \draw[bus, -{Stealth}] (add.east) -- (reg.west);
  \draw[bus, -{Stealth}] (reg.east) -- ++(1.8,0) node[right, text=FpgaDark] {$Q$};
  \draw[bus, -{Stealth}] ($(reg.east)+(0.9,0)$) -- ++(0,1.6) -| (add.north);
  \fill[FpgaBlue] ($(reg.east)+(0.9,0)$) circle (1.8pt);
  \node[lbl] at (4.4,1.9) {обратная связь};
\end{tikzpicture}
\caption{Синхронный счётчик: регистр + сумматор. Каждый такт регистр
записывает $Q+1$}
\label{fig:sync_counter}
\end{figure}

\subsection{Делитель частоты и счётчик по модулю}

Как видно на рис.~\ref{fig:count_t}, каждый разряд счётчика --- это
делитель частоты на степень двойки. А как поделить на число, не являющееся
степенью двойки, например на 10? Нужен \term{счётчик по модулю $N$}: он считает
$0, 1, \ldots, N-1$ и затем сбрасывается в 0. Для этого к схеме добавляют
компаратор, который сравнивает значение с $N-1$, и мультиплексор, который
вместо $Q+1$ подаёт на регистр 0.

\begin{figure}[H]
\centering
\begin{tikzpicture}
  \node[draw=FpgaGreen, fill=green!8, thick, minimum width=18mm, minimum height=16mm, font=\sffamily\small] (reg) at (5.6,0) {регистр};
  \draw[wire, FpgaBlue] (reg.south) -- ++(0,-0.6) node[right, text=FpgaDark, font=\sffamily\small] {CLK};
  \draw[thick, fill=orange!15] (3.0,0.9) -- (3.8,0.5) -- (3.8,-0.5) -- (3.0,-0.9) -- cycle;
  \node[font=\sffamily\scriptsize] at (3.4,0) {MUX};
  \node[font=\tiny, anchor=west] at (3.02,0.5) {0}; \node[font=\tiny, anchor=west] at (3.02,-0.5) {1};
  \node[draw=FpgaBlue, fill=FpgaLight, thick, circle, minimum size=10mm, font=\sffamily\bfseries\small] (add) at (1.2,0.5) {$+1$};
  \draw[bus, -{Stealth}] (add.east) -- (3.0,0.5);
  \draw[bus, -{Stealth}] (1.8,-0.5) node[left, text=FpgaDark] {0} -- (3.0,-0.5);
  \draw[bus, -{Stealth}] (3.8,0) -- (reg.west);
  \draw[bus, -{Stealth}] (reg.east) -- ++(1.6,0) node[right, text=FpgaDark] {$Q$};
  \draw[bus] ($(reg.east)+(0.8,0)$) -- ++(0,1.6) -| (add.north);
  \fill[FpgaBlue] ($(reg.east)+(0.8,0)$) circle (1.8pt);
  \node[block, minimum width=20mm] (cmp) at (5.6,-2.9) {$Q = N-1$?};
  \draw[bus, -{Stealth}] ($(reg.east)+(0.8,0)$) |- (cmp.east);
  \draw[arr] (cmp.west) -| node[pos=0.3, below, lbl] {сброс} (3.4,-0.72);
\end{tikzpicture}
\caption{Счётчик по модулю $N$}
\end{figure}

\begin{figure}[H]
\centering
\begin{tikzpicture}
\begin{axis}[
  width=0.82\textwidth, height=5.6cm,
  xlabel={Разрядность счётчика $n$}, ylabel={Частота, Гц},
  ymode=log, ymin=0.1, ymax=3e7, xmin=0, xmax=26, xtick={1,4,8,12,16,20,25},
  label style={font=\sffamily\small}, tick label style={font=\sffamily\small},
  axis lines*=left, grid=major, grid style={FpgaGray!20},
]
  \addplot[very thick, FpgaBlue, mark=*, mark size=1.5pt, samples at={1,2,...,25}] {27e6/2^x};
  \addplot[FpgaRed, thick, dashed, domain=0:26, samples=2] {1};
  \node[font=\sffamily\small, text=FpgaRed, anchor=south west] at (axis cs:1,1.3) {1 Гц --- мигание, заметное глазу};
  \node[font=\sffamily\small, text=FpgaBlue, anchor=south west] at (axis cs:25,0.8) {$n=25$};
\end{axis}
\end{tikzpicture}
\caption{Частота на выходе старшего разряда $n$-битного счётчика при такте
27 МГц: $f = 27\,\text{МГц}/2^n$. Чтобы получить мигание около 1 Гц, нужен
счётчик примерно на 25 бит}
\label{fig:div_plot}
\end{figure}

\begin{ideabox}
\textbf{Регистр} --- группа триггеров с общим тактом, хранит число.
\textbf{Сдвиговый регистр} перемещает биты по цепочке и преобразует
последовательный код в параллельный и обратно. \textbf{Синхронный счётчик}
--- регистр с сумматором в обратной связи; его разряды делят частоту такта
на степени двойки.
\end{ideabox}

\chapter{Синхронные схемы, автоматы и память}
\label{ch:sync}

\section{Синхронное проектирование}

Мы уже знаем два «материала»: комбинационную логику (она вычисляет) и
триггеры (они хранят). Как собрать из них большое надёжное устройство?
Почти все современные цифровые схемы строятся по одному принципу ---
\term{синхронному}:

\begin{itemize}
  \item все триггеры тактируются \textbf{одним} тактовым сигналом;
  \item между регистрами находится только комбинационная логика;
  \item результат вычислений записывается в регистры по фронту, когда все
        переходные процессы и «иголки» уже закончились.
\end{itemize}

\begin{figure}[H]
\centering
\begin{tikzpicture}[rg/.style={draw=FpgaGreen, fill=green!8, thick, minimum width=10mm, minimum height=18mm, font=\sffamily\small},
  cl/.style={draw=FpgaBlue, fill=FpgaLight, thick, cloud, cloud puffs=11, aspect=1.8, minimum width=26mm, minimum height=16mm, font=\sffamily\small, align=center}]
  \node[rg] (r1) at (0,0) {RG};
  \node[cl] (c1) at (2.6,0) {логика};
  \node[rg] (r2) at (5.2,0) {RG};
  \node[cl] (c2) at (7.8,0) {логика};
  \node[rg] (r3) at (10.4,0) {RG};
  \draw[bus, -{Stealth}] (-1.2,0) node[left, text=FpgaDark] {входы} -- (r1);
  \draw[bus, -{Stealth}] (r1) -- (c1); \draw[bus, -{Stealth}] (c1) -- (r2);
  \draw[bus, -{Stealth}] (r2) -- (c2); \draw[bus, -{Stealth}] (c2) -- (r3);
  \draw[bus, -{Stealth}] (r3) -- ++(1.2,0) node[right, text=FpgaDark] {выходы};
  \draw[wire, FpgaBlue] (-0.8,-1.6) node[left, text=FpgaDark] {CLK} -- (10.4,-1.6);
  \foreach \r in {r1,r2,r3} { \draw[wire, FpgaBlue] (\r.south) -- (\r.south |- 0,-1.6); \fill[FpgaBlue] (\r.south |- 0,-1.6) circle (1.4pt); }
\end{tikzpicture}
\caption{Синхронная схема: регистры, разделённые комбинационной логикой, и
один общий такт}
\label{fig:sync}
\end{figure}

\subsection{Максимальная тактовая частота}

За один такт сигнал должен пройти от выхода одного регистра через логику
до входа следующего и успеть «установиться» до фронта. Отсюда главное
уравнение синхронной схемы:
\[
  T_\text{CLK} \ \ge\ t_\text{co} + t_\text{логики} + t_\text{su},
  \qquad
  f_\text{max} = \frac{1}{t_\text{co} + t_\text{логики} + t_\text{su}}.
\]

\begin{figure}[H]
\centering
\begin{tikzpicture}[x=1cm, y=1cm]
  \draw[very thick, FpgaBlue] (0,2.2) -- (1,2.2) -- (1.05,2.8) -- (5,2.8) -- (5.05,2.2) -- (9,2.2) -- (9.05,2.8) -- (11,2.8);
  \node[font=\sffamily\small, anchor=east] at (-0.1,2.5) {CLK};
  \fill[FpgaOrange!40] (1.05,1.2) rectangle (2.0,1.8); \node[font=\sffamily\scriptsize] at (1.52,1.5) {$t_\text{co}$};
  \fill[FpgaBlue!25] (2.0,1.2) rectangle (7.6,1.8); \node[font=\sffamily\scriptsize] at (4.8,1.5) {$t_\text{логики}$ (критический путь)};
  \fill[FpgaRed!30] (7.6,1.2) rectangle (8.6,1.8); \node[font=\sffamily\scriptsize] at (8.1,1.5) {$t_\text{su}$};
  \fill[green!25] (8.6,1.2) rectangle (9.05,1.8);
  \node[font=\sffamily\scriptsize, anchor=west] at (9.2,1.5) {запас};
  \draw[dashed, FpgaGray] (1.05,3.0) -- (1.05,0.6); \draw[dashed, FpgaGray] (9.05,3.0) -- (9.05,0.6);
  \draw[<->, thick] (1.05,0.8) -- (9.05,0.8) node[midway, below, font=\sffamily\small] {период $T_\text{CLK}$};
\end{tikzpicture}
\caption{Бюджет времени одного такта. Если сумма задержек больше периода,
в регистр запишется неверное значение}
\label{fig:budget}
\end{figure}

Пример: $t_\text{co} = 0{,}5$~нс, $t_\text{su} = 0{,}5$~нс, а сигнал проходит
через 10 элементов по 0,5~нс. Тогда
$T \ge 0{,}5 + 5 + 0{,}5 = 6$~нс и $f_\text{max} \approx 166$~МГц. Чем
длиннее цепочка логики между регистрами, тем ниже допустимая частота.

\begin{figure}[H]
\centering
\begin{tikzpicture}
\begin{axis}[
  width=0.78\textwidth, height=5.2cm, xlabel={Число уровней логики на критическом пути}, ylabel={$f_\text{max}$, МГц},
  xmin=1, xmax=20, ymin=0, ymax=520,
  label style={font=\sffamily\small}, tick label style={font=\sffamily\small},
  grid=major, grid style={FpgaGray!20}, axis lines*=left,
]
  \addplot[very thick, FpgaBlue, domain=1:20, samples=40] {1000/(1+0.5*x)};
  \addplot[only marks, mark=*, FpgaRed] coordinates {(10,166.7)};
  \node[font=\sffamily\small, text=FpgaRed, anchor=south west] at (axis cs:10,170) {пример: 166 МГц};
\end{axis}
\end{tikzpicture}
\caption{Максимальная частота падает с ростом числа элементов между регистрами
($t_\text{co}+t_\text{su} = 1$~нс, задержка элемента 0,5~нс)}
\end{figure}

\subsection{Конвейер}

Что делать, если логика слишком длинная? Её можно разрезать на части и
поставить между ними дополнительные регистры. Это \term{конвейер}
(pipeline): каждая «станция» делает свою часть работы за один короткий такт,
и хотя результат для одних данных появляется через несколько тактов, новые
данные можно подавать \emph{каждый} такт. Так работает сборочный конвейер на
заводе --- и все современные процессоры.

\begin{figure}[H]
\centering
\begin{tikzpicture}[c/.style={draw, minimum width=11mm, minimum height=6mm, font=\sffamily\scriptsize}]
  \foreach \t in {1,...,6} \node[font=\sffamily\scriptsize] at (\t*1.15+0.9,0.55) {такт \t};
  \foreach \s/\n [count=\r from 0] in {этап 1/A, этап 2/B, этап 3/C} {
    \node[font=\sffamily\small, anchor=east] at (1.3,-\r*0.75) {\s};
  }
  \foreach \d/\col [count=\k from 0] in {1/FpgaBlue!30, 2/FpgaOrange!40, 3/FpgaGreen!35, 4/FpgaPurple!30} {
    \foreach \r in {0,1,2} {
      \pgfmathsetmacro{\x}{(\k+\r+1)*1.15+0.9}
      \node[c, fill=\col] at (\x,-\r*0.75) {данные \d};
    }
  }
\end{tikzpicture}
\caption{Трёхэтапный конвейер: на такте 3 одновременно обрабатываются три
порции данных, каждая на своём этапе}
\end{figure}

\section{Конечные автоматы}

Многие устройства работают по сценарию: «ждать нажатия, потом сделать то, потом
это, при ошибке вернуться в начало». Такое поведение описывают
\term{конечным автоматом} (FSM, finite state machine). Автомат находится в
одном из конечного числа \term{состояний}; по каждому такту, в зависимости от
текущего состояния и входов, он \term{переходит} в следующее состояние.

Схема любого автомата одинакова: регистр хранит номер текущего состояния,
комбинационная логика вычисляет следующее состояние и выходы.

\begin{figure}[H]
\centering
\begin{tikzpicture}[node distance=12mm]
  \node[block, minimum width=28mm, minimum height=14mm] (next) {Логика\\переходов};
  \node[gblock, right=of next, minimum height=14mm] (st) {Регистр\\состояния};
  \node[block, right=of st, minimum width=28mm, minimum height=14mm] (out) {Логика\\выходов};
  \draw[bus, -{Stealth}] (next) -- node[above, lbl] {след.} (st);
  \draw[bus, -{Stealth}] (st) -- node[above, lbl] {текущее} (out);
  \draw[arr] (out.east) -- ++(1,0) node[right, font=\sffamily\small] {выходы};
  \draw[bus, -{Stealth}] ($(st.east)+(0.6,0)$) -- ++(0,-1.3) -| (next.south);
  \draw[arr] ($(next.west)+(-1.2,0)$) node[left, font=\sffamily\small] {входы} -- (next.west);
  \draw[arr, dashed, FpgaPurple] ($(next.west)+(-0.6,0)$) -- ++(0,1.3) -| node[pos=0.25, above, lbl, text=FpgaPurple] {только у автомата Мили} (out.north);
  \draw[wire, FpgaBlue] (st.south) ++(0,0) -- ++(0,-0.5) node[below, font=\sffamily\scriptsize, text=FpgaDark] {CLK};
\end{tikzpicture}
\caption{Структура конечного автомата. У автомата \textbf{Мура} выходы зависят
только от состояния, у автомата \textbf{Мили} --- ещё и от входов (пунктир)}
\label{fig:fsm_struct}
\end{figure}

\subsection{Пример: детектор последовательности}

Построим автомат, который следит за входом $x$ (по одному биту за такт) и
выдаёт $y=1$, когда последние три принятых бита равны \code{101}. Такие
детекторы ищут стартовые метки в потоке данных.

\begin{figure}[H]
\centering
\begin{tikzpicture}[
  st/.style={circle, draw=FpgaDark, very thick, minimum size=17mm, font=\sffamily\small, align=center, fill=FpgaLight},
  tr/.style={-{Stealth[length=3mm]}, thick, FpgaDark},
  lb/.style={font=\sffamily\small, fill=white, inner sep=1pt}]
  \node[st] (s0) at (0,0) {S0\\\scriptsize $y=0$};
  \node[st] (s1) at (3.6,0) {S1\\\scriptsize «1»\\\scriptsize $y=0$};
  \node[st] (s2) at (7.2,0) {S2\\\scriptsize «10»\\\scriptsize $y=0$};
  \node[st, fill=green!20] (s3) at (10.8,0) {S3\\\scriptsize «101»\\\scriptsize $y=1$};
  \draw[tr] (s0) -- node[lb, above] {1} (s1);
  \draw[tr] (s1) -- node[lb, above] {0} (s2);
  \draw[tr] (s2) -- node[lb, above] {1} (s3);
  \draw[tr] (s0) edge[loop above, looseness=6] node[lb, above] {0} (s0);
  \draw[tr] (s1) edge[loop above, looseness=6] node[lb, above] {1} (s1);
  \draw[tr] (s2) to[bend left=35] node[lb, below] {0} (s0);
  \draw[tr] (s3) to[bend left=32] node[lb, below] {0} (s2);
  \draw[tr] (s3) to[bend right=40] node[lb, above] {1} (s1);
  \draw[tr, FpgaRed] (-1.6,1.4) node[left, font=\sffamily\small, text=FpgaRed] {сброс} -- (s0);
\end{tikzpicture}
\caption{Граф переходов детектора \code{101} (автомат Мура). На стрелках ---
значение входа $x$}
\label{fig:fsm101}
\end{figure}

\begin{table}[H]
\centering
\caption{Таблица переходов детектора}
\begin{tabular}{l|cc|c}
\toprule
\textbf{Текущее состояние} & \multicolumn{2}{c|}{\textbf{Следующее состояние}} & \textbf{Выход} $y$ \\
 & $x=0$ & $x=1$ & \\
\midrule
S0 (ничего не найдено) & S0 & S1 & 0 \\
S1 (принято «1»)       & S2 & S1 & 0 \\
S2 (принято «10»)      & S0 & S3 & 0 \\
S3 (принято «101»)     & S2 & S1 & 1 \\
\bottomrule
\end{tabular}
\end{table}

\begin{figure}[H]
\centering
\begin{tikztimingtable}[timing/slope=0.05, xscale=1.3, yscale=1.4, timing/font=\sffamily\small, timing/rowdist=3.2ex]
  CLK     & 10{1L 1H} \\
  $x$     & 2H 2L 2H 2L 2H 2H 2L 2L 2H 2L \\
  state   & 1D{S0} 2D{S1} 2D{S2} 2D{S3} 2D{S2} 2D{S3} 2D{S1} 2D{S2} 2D{S0} 2D{S1} 1D{S2} \\
  $y$     & 5L 2H 2L 2H 9L \\
\end{tikztimingtable}
\caption{Работа детектора: вход \code{1010110010}. Выход равен 1 после
каждой найденной тройки \code{101}, в том числе перекрывающейся}
\end{figure}

Проверим по таблице: S0 --(1)$\to$ S1 --(0)$\to$ S2 --(1)$\to$ S3 ($y=1$)
--(0)$\to$ S2 --(1)$\to$ S3 ($y=1$) --(1)$\to$ S1 --(0)$\to$ S2 --(0)$\to$ S0
--(1)$\to$ S1.

\subsection{Кодирование состояний}

Состояния хранятся в регистре в виде двоичных кодов. Четыре состояния можно
закодировать двумя битами (\code{00}, \code{01}, \code{10}, \code{11}) ---
это \term{двоичное кодирование}. Можно взять четыре триггера и держать единицу
только в одном из них (\code{0001}, \code{0010}, \ldots) --- это
\term{унитарное кодирование} (one-hot). Оно требует больше триггеров, зато
логика переходов проще и быстрее. В ПЛИС, где триггеров много, синтезатор
часто выбирает именно его.

\section{Память}

Регистр из $n$ триггеров хранит одно число. Чтобы хранить много чисел,
используют \term{память}: массив ячеек, каждая из которых имеет свой номер ---
\term{адрес}. Память характеризуется \term{глубиной} (числом ячеек) и
\term{шириной} (числом битов в ячейке), например $1024 \times 8$ ---
1024 ячейки по 8 бит.

\subsection{ПЗУ: память только для чтения}

\term{ПЗУ} (постоянное запоминающее устройство, ROM) хранит данные, записанные
заранее. По сути это комбинационная схема: адрес --- вход, данные --- выход,
а содержимое --- большая таблица истинности. ПЗУ можно построить из
дешифратора адреса и элементов ИЛИ: каждая ячейка --- это выход дешифратора,
подключённый к тем элементам ИЛИ, где в данных стоит 1.

\begin{figure}[H]
\centering
\begin{tikzpicture}
  \draw[thick, fill=FpgaLight] (0,-2.8) rectangle (1.4,0.6);
  \node[font=\sffamily\bfseries] at (0.7,-1.1) {DC};
  \draw[wire] (-0.7,-0.4) node[left] {$a_1$} -- (0,-0.4);
  \draw[wire] (-0.7,-1.8) node[left] {$a_0$} -- (0,-1.8);
  \foreach \i in {0,...,3} {
    \pgfmathsetmacro{\y}{0.2-\i*0.9}
    \draw[wire] (1.4,\y) -- (6.4,\y);
    \node[font=\sffamily\scriptsize, anchor=east] at (1.35,\y) {\i};
  }
  \foreach \j in {0,...,3} \draw[very thick, FpgaBlue] (2.4+\j*1.1,0.5) -- (2.4+\j*1.1,-3.2);
  % данные: ячейка0=1011, 1=0110, 2=1100, 3=0001 (d3..d0)
  \foreach \i/\row in {0/1011, 1/0110, 2/1100, 3/0001} {
    \pgfmathsetmacro{\y}{0.2-\i*0.9}
    \foreach \j in {0,...,3} {
      \StrChar{\row}{\numexpr\j+1\relax}[\bitv]
      \ifnum\bitv=1 \fill[FpgaRed] (2.4+\j*1.1,\y) circle (2.4pt);\fi
    }
    \node[font=\ttfamily\small, anchor=west] at (6.6,\y) {\row};
  }
  \foreach \j/\n in {0/3,1/2,2/1,3/0} {
    \node[gOR, rotate=-90, minimum width=9mm, minimum height=6mm] (o\j) at (2.4+\j*1.1,-3.8) {};
    \draw[wire] (o\j.east) -- ++(0,-0.4) node[below] {$d_\n$};
  }
  \node[lbl, anchor=west] at (6.6,0.75) {содержимое};
  \node[lbl, text=FpgaRed, align=left, anchor=west] at (6.6,-3.6) {точка --- соединение\\(в этом бите 1)};
\end{tikzpicture}
\caption{ПЗУ $4\times4$: дешифратор выбирает строку, элементы ИЛИ собирают
биты. Тот же принцип --- «память как таблица истинности» --- лежит в основе LUT}
\label{fig:rom}
\end{figure}

\subsection{ОЗУ: оперативная память}

\term{ОЗУ} (оперативное запоминающее устройство, RAM) позволяет и читать, и
записывать. Концептуально это набор регистров, дешифратор адреса, который при
записи выбирает, в какой регистр записать данные, и мультиплексор, который
при чтении выбирает, какой регистр подать на выход.

\begin{figure}[H]
\centering
\begin{tikzpicture}
  \draw[thick, fill=FpgaLight] (0,-2.9) rectangle (1.3,0.7);
  \node[font=\sffamily\bfseries] at (0.65,-1.1) {DC};
  \foreach \i in {0,...,3} {
    \pgfmathsetmacro{\y}{0.25-\i*0.95}
    \node[draw=FpgaGreen, fill=green!8, thick, minimum width=22mm, minimum height=7mm, font=\sffamily\small] (r\i) at (3.6,\y) {ячейка \i};
    \draw[arr] (1.3,\y) -- node[above, lbl] {CE} (r\i.west);
    \draw[bus, -{Stealth}] (r\i.east) -- (5.9,\y);
  }
  \draw[thick, fill=orange!15] (5.9,0.7) -- (7.0,0.2) -- (7.0,-2.4) -- (5.9,-2.9) -- cycle;
  \node[font=\sffamily\small] at (6.45,-1.1) {MUX};
  \draw[bus, -{Stealth}] (7.0,-1.1) -- (8.2,-1.1) node[right, text=FpgaDark] {данные чтения};
  \draw[bus] (3.6,1.6) node[above, text=FpgaDark] {данные записи} -- (3.6,0.6);
  \draw[arr] (-0.8,0.3) node[left, font=\sffamily\small] {WE (запись)} -- (0,0.3);
  \draw[bus, -{Stealth}] (-0.8,-1.1) node[left, text=FpgaDark] {адрес} -- (0,-1.1);
  \draw[bus] (-0.4,-1.1) |- (6.45,-3.6) -- (6.45,-2.65); \fill[FpgaBlue] (-0.4,-1.1) circle (1.8pt);
\end{tikzpicture}
\caption{Устройство ОЗУ «на пальцах»: дешифратор адреса при записи,
мультиплексор при чтении}
\label{fig:ram}
\end{figure}

На практике ячейки памяти делают не из D-триггеров, а из более компактных
ячеек:
\begin{itemize}
  \item \textbf{статическая память SRAM} --- 6 транзисторов на бит, быстрая,
        не требует обслуживания; из неё сделаны блоки памяти внутри ПЛИС;
  \item \textbf{динамическая память DRAM} --- 1 транзистор и 1 конденсатор на
        бит, очень плотная, но требует периодической регенерации; такая
        память --- SDRAM --- встроена в корпус ПЛИС на Tang Nano 20K
        (глава~\ref{ch:sdram}).
\end{itemize}

\section{Что дальше}

Мы прошли путь от сигнала до памяти и конечных автоматов. Этого достаточно,
чтобы понять, как устроено практически любое цифровое устройство: процессор,
контроллер дрона, видеокарта. Все эти узлы можно собрать и внутри ПЛИС ---
причём не паяльником, а описав их текстом.

\begin{table}[H]
\centering
\caption{Узлы из первой части и то, как они реализуются в ПЛИС}
\small
\begin{tabularx}{\textwidth}{X X}
\toprule
\textbf{Узел} & \textbf{Ресурс ПЛИС (часть II)} \\
\midrule
Логические элементы, дешифраторы, шифраторы, мультиплексоры & таблицы LUT (рис.~\ref{fig:mux_lut}) \\
Сумматоры, счётчики & LUT + быстрые цепочки переноса (ALU) \\
Умножители & аппаратные блоки DSP \\
Триггеры, регистры, сдвиговые регистры & триггеры в каждой логической ячейке \\
Конечные автоматы & LUT + триггеры \\
ПЗУ и ОЗУ & блочная память BSRAM, распределённая память SSRAM \\
Большая память & встроенная SDRAM 64 Мбит \\
Тактовый сигнал & генератор 27 МГц, PLL, глобальная тактовая сеть \\
\bottomrule
\end{tabularx}
\end{table}

\begin{ideabox}
Синхронная схема --- это регистры и комбинационная логика между ними с одним
общим тактом. Её быстродействие определяется критическим путём. Конечный
автомат --- регистр состояния плюс логика переходов. Память --- массив ячеек
с дешифратором адреса. Во второй части мы соберём всё это в ПЛИС.
\end{ideabox}


% ============================================================
% ЧАСТЬ II. ПЛИС НА ПРАКТИКЕ: TANG NANO 20K
% ============================================================
\part{ПЛИС на практике: Tang Nano 20K}
\chapter{Что такое ПЛИС}

\section{Схема, которую можно перепаять без паяльника}

Представьте, что вам нужно собрать устройство, которое включает лампу,
когда нажаты \emph{обе} кнопки. Вы берёте микросхему с элементами И,
паяете провода, и всё работает. А завтра заказчик просит: «Пусть лампа
горит, когда нажата \emph{ровно одна} кнопка». Придётся выпаивать
элемент И, ставить элемент «Исключающее ИЛИ» и перекладывать провода.

А теперь представьте микросхему, внутри которой лежат тысячи
«заготовок» логических элементов и триггеров и огромная сеть
проводов с управляемыми переключателями на каждом перекрёстке. Какие
заготовки чем станут и какие переключатели будут замкнуты, задаётся
специальным файлом конфигурации, который загружается в микросхему по кабелю
за доли секунды. Это и есть \term{ПЛИС}.

\begin{ideabox}
\textbf{ПЛИС} (программируемая логическая интегральная схема, англ.
\textbf{FPGA}, Field-Programmable Gate Array: «вентильная матрица,
программируемая в полевых условиях») --- это микросхема, внутреннюю
цифровую схему которой можно многократно изменять после её производства.
\end{ideabox}

Важно сразу понять разницу: ПЛИС \emph{не выполняет программу}, как
процессор. Она \emph{становится схемой}. После загрузки конфигурации внутри
кристалла существуют настоящие логические элементы и настоящие провода,
по которым сигналы бегут одновременно, как в схеме на макетной плате.

\begin{figure}[H]
\centering
\begin{tikzpicture}
  % Вариант 1
  \begin{scope}
    \node[font=\sffamily\bfseries\small] at (1.4,1.9) {Конфигурация №1};
    \node[gAND] (g1) at (1.4,0) {};
    \draw[wire] (g1.input 1) -- ++(-0.8,0) node[left] {$a$};
    \draw[wire] (g1.input 2) -- ++(-0.8,0) node[left] {$b$};
    \draw[wire] (g1.output) -- ++(0.8,0) node[right] {$y=a\cdot b$};
  \end{scope}
  % Стрелка
  \draw[-{Stealth[length=4mm]}, line width=2pt, FpgaOrange] (4.8,0) -- (6.4,0)
    node[midway, above=2mm, font=\sffamily\scriptsize, text=FpgaDark, align=center] {новый файл\\конфигурации};
  % Вариант 2
  \begin{scope}[xshift=8.6cm]
    \node[font=\sffamily\bfseries\small] at (0.3,1.9) {Конфигурация №2};
    \node[gXOR] (g2) at (0.3,0) {};
    \draw[wire] (g2.input 1) -- ++(-0.8,0) node[left] {$a$};
    \draw[wire] (g2.input 2) -- ++(-0.8,0) node[left] {$b$};
    \draw[wire] (g2.output) -- ++(0.8,0) node[right] {$y=a\oplus b$};
  \end{scope}
  \draw[dashed, FpgaGray, rounded corners] (-0.3,-1.1) rectangle (11.8,2.4);
  \node[lbl, anchor=north] at (5.75,-1.2) {одна и та же микросхема ПЛИС, одни и те же ножки};
\end{tikzpicture}
\caption{Одна и та же ПЛИС после загрузки разных конфигураций}
\label{fig:reconfig}
\end{figure}

\section{Немного истории}

Идея программируемой логики появилась в 1970-х. Сначала это были
небольшие матрицы И--ИЛИ (PAL, PLA), в которых пережигались перемычки.
Затем появились CPLD (сложные программируемые логические устройства), а в 1985 году компания Xilinx
выпустила первую FPGA XC2064, в которой было всего 64 логических блока.
Современные ПЛИС содержат миллионы логических ячеек.

\begin{figure}[H]
\centering
\begin{tikzpicture}
  \draw[line width=2pt, FpgaBlue, -{Stealth[length=4mm]}] (0,0) -- (15,0);
  \foreach \x/\y/\t/\d in {
      0.8/1/1975/PLA, 3.2/-1/1978/PAL,
      5.6/1/1984/CPLD (Altera), 8/-1/1985/{Xilinx XC2064,\\ первая FPGA},
      10.4/1/2000-е/{FPGA с DSP\\и памятью}, 13/-1/2020-е/{Tang Nano 20K:\\ 20\,736 LUT за копейки}} {
    \fill[FpgaOrange] (\x,0) circle (1.1mm);
    \draw[FpgaGray] (\x,0) -- (\x,\y*0.55);
    \node[font=\sffamily\bfseries\small, text=FpgaDark] at (\x,\y*0.85) {\t};
    \node[font=\sffamily\scriptsize, align=center, text=FpgaDark, anchor={\ifnum\y>0 south\else north\fi}] at (\x,\y*1.1) {\d};
  }
\end{tikzpicture}
\caption{Эволюция программируемой логики}
\end{figure}

\section{ПЛИС, микроконтроллер и заказная микросхема}

Чтобы понять, зачем нужны ПЛИС, сравним три способа сделать «умное»
цифровое устройство.

\begin{description}
  \item[Микроконтроллер (МК)] --- готовый процессор, который
    \emph{последовательно}, команда за командой, выполняет программу. Например,
    на МК STM32 работает автопилот квадрокоптера «Геоскан Пионер».
  \item[Заказная микросхема (ASIC)] --- схема, «выжженная» в кремнии на
    заводе. Она работает быстрее и экономичнее всех, но её разработка стоит
    миллионы долларов, и изменить её уже нельзя.
  \item[ПЛИС (FPGA)] --- золотая середина: аппаратная схема, как у ASIC,
    которую можно менять, как программу МК.
\end{description}

\begin{figure}[H]
\centering
\begin{tikzpicture}
  % МК
  \begin{scope}
    \node[font=\sffamily\bfseries] at (2.6,3.6) {Микроконтроллер};
    \node[hblock, minimum width=16mm] (cpu) at (2.6,2.4) {ЦП};
    \foreach \i/\t in {0/{$a\cdot b$}, 1/{$c+d$}, 2/{$e\oplus f$}, 3/{$g\cdot h$}} {
      \node[block, minimum width=14mm, minimum height=6mm, font=\sffamily\scriptsize]
        (t\i) at (2.6, 1.2-\i*0.8) {\t};
    }
    \draw[arr] (cpu) -- (t0);
    \draw[arr] (t0) -- (t1); \draw[arr] (t1) -- (t2); \draw[arr] (t2) -- (t3);
    \draw[arr, FpgaOrange] (0.9,1.3) -- (0.9,-1.3) node[midway, left, font=\sffamily\scriptsize, align=right] {время\\4 шага};
  \end{scope}
  % ПЛИС
  \begin{scope}[xshift=7.2cm]
    \node[font=\sffamily\bfseries] at (2.8,3.6) {ПЛИС};
    \foreach \i/\t in {0/{$a\cdot b$}, 1/{$c+d$}, 2/{$e\oplus f$}, 3/{$g\cdot h$}} {
      \node[gblock, minimum width=12mm, minimum height=6mm, font=\sffamily\scriptsize]
        (p\i) at (\i*1.5+0.5, 1.2) {\t};
      \draw[arr] (\i*1.5+0.5, 2.5) -- (p\i);
      \draw[arr] (p\i) -- (\i*1.5+0.5, 0);
    }
    \draw[arr, FpgaOrange] (-0.5,1.5) -- (-0.5,0.9) node[midway, left, font=\sffamily\scriptsize, align=right] {1 шаг};
    \node[lbl, align=center] at (2.75,-0.6) {все четыре операции --\\ отдельные схемы, работают \emph{одновременно}};
  \end{scope}
\end{tikzpicture}
\caption{Процессор выполняет действия по очереди, ПЛИС --- параллельно}
\label{fig:par}
\end{figure}

На рис.~\ref{fig:par} видна главная суперсила ПЛИС --- \term{параллелизм}.
Если нужно выполнить 1000 независимых операций, процессору нужно
1000 шагов, а в ПЛИС можно построить 1000 схем, которые справятся за один
такт.

\begin{figure}[H]
\centering
\begin{tikzpicture}
\begin{axis}[
  width=0.82\textwidth, height=6.2cm,
  xlabel={Число независимых операций}, ylabel={Время, мкс},
  xmin=0, xmax=1000, ymin=0, ymax=12,
  grid=major, grid style={FpgaGray!25},
  legend pos=north west, legend style={font=\sffamily\small, draw=none, fill=white, fill opacity=0.8},
  tick label style={font=\small}, label style={font=\sffamily\small},
]
  \addplot[very thick, FpgaRed, domain=0:1000, samples=2] {x*0.0104};
  \addlegendentry{Микроконтроллер 96 МГц (1 операция за такт)}
  \addplot[very thick, FpgaGreen, domain=0:1000, samples=2] {0.037};
  \addlegendentry{ПЛИС 27 МГц (все операции за 1 такт)}
\end{axis}
\end{tikzpicture}
\caption{Время обработки при росте числа операций (упрощённая оценка)}
\end{figure}

\begin{table}[H]
\centering
\caption{Сравнение трёх подходов}
\small
\begin{tabularx}{\textwidth}{>{\bfseries}X c c c}
\toprule
 & \textbf{Микроконтроллер} & \textbf{ПЛИС} & \textbf{ASIC} \\
\midrule
Как работает          & программа, по шагам & схема, параллельно & схема, параллельно \\
Скорость реакции      & микросекунды        & наносекунды        & наносекунды \\
Можно изменить?       & да                  & да                 & нет \\
Цена одного образца   & низкая              & средняя            & огромная \\
Сложность разработки  & низкая              & средняя            & очень высокая \\
Энергопотребление     & низкое              & среднее            & самое низкое \\
\bottomrule
\end{tabularx}
\end{table}

\section{Где применяются ПЛИС}

ПЛИС выбирают там, где нужно одно из трёх: \emph{очень быстро},
\emph{очень много параллельно} или \emph{точно по времени}.

\begin{figure}[H]
\centering
\begin{tikzpicture}[
  app/.style={block, minimum width=34mm, minimum height=13mm, font=\sffamily\scriptsize},
]
  \node[hblock, circle, minimum size=24mm, font=\sffamily\bfseries] (c) at (0,0) {ПЛИС};
  \foreach \a/\t/\s in {
     90/{Обработка видео\\и машинное зрение}/app,
     30/{Связь: 5G, радары,\\программное радио}/app,
    -30/{Управление моторами,\\ШИМ, датчики дронов}/app,
    -90/{Прототипирование\\процессоров (ASIC)}/app,
   -150/{Сопряжение\\интерфейсов}/app,
    150/{Ускорение ИИ\\и вычислений}/app} {
    \node[\s] (n\a) at (\a:4.1) {\t};
    \draw[thick, FpgaBlue!60] (c) -- (n\a);
  }
\end{tikzpicture}
\caption{Основные области применения ПЛИС}
\end{figure}

\begin{itemize}
  \item \textbf{Цифровая обработка сигналов}: фильтры, преобразование
        Фурье, обработка звука и радиосигналов в реальном времени.
  \item \textbf{Видео}: генерация изображения для HDMI, обработка кадров с камер,
        поиск объектов прямо «на лету».
  \item \textbf{Связь и интерфейсы}: ПЛИС может «говорить» на любом
        протоколе (UART, SPI, I$^2$C, Ethernet, USB), даже на придуманном вами.
  \item \textbf{Робототехника и БПЛА}: точная генерация десятков ШИМ-сигналов,
        одновременный опрос многих датчиков, стабилизация видео.
  \item \textbf{Прототипирование}: прежде чем заказывать ASIC за
        миллионы долларов, инженеры проверяют его на ПЛИС.
  \item \textbf{Ретро-компьютеры}: на ПЛИС воссоздают старые игровые
        приставки и компьютеры (проект MiSTer), причём с точностью до такта.
\end{itemize}

\begin{ideabox}
ПЛИС --- это «пластилин» для цифровой схемотехники. Вы описываете схему,
программа-синтезатор раскладывает её по ресурсам кристалла, и через
минуту ваша схема уже работает в железе.
\end{ideabox}

\chapter{Что внутри ПЛИС}
\label{ch:inside}

Если снять крышку с ПЛИС и посмотреть на кристалл под микроскопом, мы
увидим аккуратную «шахматную доску»: одинаковые логические блоки, окружённые
сетью проводников, а по краям --- блоки ввода-вывода. Разберём все эти
части по порядку.

\begin{figure}[H]
\centering
\begin{tikzpicture}[scale=0.92]
  % Кристалл
  \fill[ChipBlack!8] (-0.6,-0.6) rectangle (12.6,8.6);
  \draw[very thick, FpgaDark] (-0.6,-0.6) rectangle (12.6,8.6);
  % Трассировочные каналы
  \foreach \x in {0,...,12} \draw[FpgaOrange!55, line width=0.7pt] (\x,0) -- (\x,8);
  \foreach \y in {0,...,8} \draw[FpgaOrange!55, line width=0.7pt] (0,\y) -- (12,\y);
  % Логические блоки / память / DSP
  \foreach \x in {0,...,11} \foreach \y in {0,...,7} {
    \pgfmathtruncatemacro{\col}{\x}
    \ifnum\col=4
      \fill[FpgaPurple!65] (\x+0.15,\y+0.15) rectangle ++(0.7,0.7);
    \else\ifnum\col=8
      \fill[FpgaRed!65] (\x+0.15,\y+0.15) rectangle ++(0.7,0.7);
    \else
      \fill[FpgaBlue!75] (\x+0.15,\y+0.15) rectangle ++(0.7,0.7);
    \fi\fi
  }
  % Коммутаторы на пересечениях
  \foreach \x in {0,...,12} \foreach \y in {0,...,8}
    \fill[FpgaOrange] (\x,\y) circle (0.06);
  % Блоки ввода-вывода
  \foreach \x in {0,...,11} {
    \fill[FpgaGreen] (\x+0.25,-0.5) rectangle ++(0.5,0.3);
    \fill[FpgaGreen] (\x+0.25,8.2) rectangle ++(0.5,0.3);
  }
  \foreach \y in {0,...,7} {
    \fill[FpgaGreen] (-0.5,\y+0.25) rectangle ++(0.3,0.5);
    \fill[FpgaGreen] (12.2,\y+0.25) rectangle ++(0.3,0.5);
  }
  % PLL
  \fill[yellow!80!orange] (-0.55,-0.55) rectangle (-0.05,-0.05);
  \fill[yellow!80!orange] (12.05,8.05) rectangle (12.55,8.55);
  % Легенда
  \begin{scope}[xshift=13.3cm, yshift=7.4cm, every node/.style={font=\sffamily\small, anchor=west}]
    \fill[FpgaBlue!75] (0,0) rectangle (0.4,0.4); \node at (0.5,0.2) {CFU (логика)};
    \fill[FpgaPurple!65] (0,-0.8) rectangle (0.4,-0.4); \node at (0.5,-0.6) {BSRAM (память)};
    \fill[FpgaRed!65] (0,-1.6) rectangle (0.4,-1.2); \node at (0.5,-1.4) {DSP (умножители)};
    \fill[FpgaGreen] (0,-2.4) rectangle (0.4,-2.0); \node at (0.5,-2.2) {IOB (ввод-вывод)};
    \fill[yellow!80!orange] (0,-3.2) rectangle (0.4,-2.8); \node at (0.5,-3.0) {PLL (такты)};
    \draw[FpgaOrange!70, line width=1pt] (0,-3.8) -- (0.4,-3.8); \node at (0.5,-3.8) {трассировка};
    \fill[FpgaOrange] (0.2,-4.5) circle (0.08); \node at (0.5,-4.5) {коммутатор};
  \end{scope}
\end{tikzpicture}
\caption{Упрощённое устройство кристалла ПЛИС}
\label{fig:fpga_grid}
\end{figure}

\section{Главная идея: таблица истинности как память}

Как сделать универсальный логический элемент, который по желанию
становится И, ИЛИ, «Исключающим ИЛИ» или чем угодно ещё? Инженеры
придумали изящный трюк.

Любая логическая функция от $n$ входов полностью задаётся своей
таблицей истинности, в которой $2^n$ строк. Значит, достаточно
\emph{запомнить} правый столбец таблицы в маленькой памяти, а входы
использовать как \emph{адрес} ячейки. Такое устройство называется
\term{LUT} (Look-Up Table, таблица поиска).

\begin{table}[H]
\centering
\caption{Одна «заготовка» LUT2 --- четыре разных функции}
\label{tab:lut2}
\begin{tabular}{cc|c|cccc}
\toprule
$a$ & $b$ & адрес & И & ИЛИ & И-НЕ & Искл. ИЛИ \\
\midrule
0 & 0 & 0 & 0 & 0 & 1 & 0 \\
0 & 1 & 1 & 0 & 1 & 1 & 1 \\
1 & 0 & 2 & 0 & 1 & 1 & 1 \\
1 & 1 & 3 & 1 & 1 & 0 & 0 \\
\midrule
\multicolumn{3}{r|}{\textbf{содержимое памяти:}} & \code{1000} & \code{1110} & \code{0111} & \code{0110} \\
\bottomrule
\end{tabular}
\end{table}

Чтобы LUT «стал» элементом И, в его память записывают числа
$0,0,0,1$; чтобы он стал ИЛИ --- $0,1,1,1$. Никакой перепайки:
меняются лишь четыре бита конфигурации.

\begin{figure}[H]
\centering
\begin{tikzpicture}[
  cell/.style={draw=FpgaDark, fill=FpgaLight, minimum width=9mm, minimum height=6.5mm, font=\ttfamily\small},
  mux/.style={draw=FpgaDark, fill=orange!15, thick, trapezium, trapezium angle=70,
              shape border rotate=270, minimum width=13mm, minimum height=6mm, font=\sffamily\tiny},
]
  % Ячейки памяти
  \foreach \i/\v in {0/0,1/0,2/0,3/1} {
    \node[cell] (m\i) at (0,-\i*1.0) {\v};
    \node[lbl, left=1mm of m\i] {ячейка \i};
  }
  % Уровень 1 (управляется b)
  \node[mux] (x0) at (2.6,-0.5) {MUX};
  \node[mux] (x1) at (2.6,-2.5) {MUX};
  \draw[wire] (m0.east) -- ++(0.5,0) |- ([yshift=1.3mm]x0.west);
  \draw[wire] (m1.east) -- ++(0.5,0) |- ([yshift=-1.3mm]x0.west);
  \draw[wire] (m2.east) -- ++(0.5,0) |- ([yshift=1.3mm]x1.west);
  \draw[wire] (m3.east) -- ++(0.5,0) |- ([yshift=-1.3mm]x1.west);
  % Уровень 2 (управляется a)
  \node[mux] (y) at (5.0,-1.5) {MUX};
  \draw[wire] (x0.east) -- ++(0.4,0) |- ([yshift=1.3mm]y.west);
  \draw[wire] (x1.east) -- ++(0.4,0) |- ([yshift=-1.3mm]y.west);
  \draw[wire] (y.east) -- ++(1.0,0) node[right] {$y = f(a,b)$};
  % Управляющие входы
  \draw[wire, FpgaBlue] (2.6,-4.0) node[below] {$b$} -- (x1.south);
  \draw[wire, FpgaBlue] (2.6,-3.5) -| (1.9,-1.2) -| (x0.south);
  \draw[wire, FpgaBlue] (5.0,-4.0) node[below] {$a$} -- (y.south);
  % Рамка
  \draw[dashed, FpgaGray, rounded corners] (-1.9,0.6) rectangle (5.7,-4.6);
  \node[lbl, anchor=south west] at (-1.9,0.6) {LUT2, настроенный как элемент И};
  \node[lbl, anchor=south, align=center, text=FpgaGray] at (0,0.3) {конфигурация};
\end{tikzpicture}
\caption{Устройство LUT2: четыре бита памяти и дерево мультиплексоров. Входы
$a$ и $b$ выбирают одну из ячеек}
\label{fig:lut2}
\end{figure}

В ПЛИС Gowin, установленной на Tang Nano 20K, используются
\term{LUT4} --- таблицы на 4 входа. Сколько различных функций может
реализовать один LUT4? В памяти 16 бит, каждый может быть 0 или 1:
\[
  N = 2^{2^4} = 2^{16} = 65\,536 \text{ функций.}
\]
Среди них и И, и ИЛИ-НЕ, и мажоритарный элемент, и полусумматор,
и вообще всё, что вы можете придумать для четырёх входов.

\begin{figure}[H]
\centering
\begin{tikzpicture}
\begin{axis}[
  ybar, bar width=14mm, width=0.7\textwidth, height=6cm,
  ymode=log, log origin=infty, ymin=1, ymax=1e6,
  symbolic x coords={LUT1, LUT2, LUT3, LUT4},
  xtick=data, ylabel={Число возможных функций},
  nodes near coords, nodes near coords style={font=\sffamily\small},
  point meta=explicit symbolic,
  label style={font=\sffamily\small}, tick label style={font=\sffamily\small},
  grid=major, grid style={FpgaGray!20}, axis lines*=left,
]
  \addplot[fill=FpgaBlue!70, draw=FpgaBlue] coordinates {
    (LUT1,4) [4] (LUT2,16) [16] (LUT3,256) [256] (LUT4,65536) [65\,536]};
\end{axis}
\end{tikzpicture}
\caption{Число функций, которые умеет реализовывать LUT с $n$ входами ($2^{2^n}$)}
\end{figure}

\begin{warnbox}[А если входов больше четырёх?]
Функция от 5, 6 и более переменных раскладывается на несколько LUT4,
соединённых между собой. Синтезатор делает это автоматически, а
специальные мультиплексоры MUX2 рядом с LUT помогают собрать из двух LUT4
один LUT5, из двух LUT5 --- LUT6 и т.\,д. Вы увидите их в отчёте синтеза
под именами \code{MUX2\_LUT5}, \code{MUX2\_LUT6}.
\end{warnbox}

\section{Триггер: у схемы появляется память}

Одних LUT мало. Из главы~\ref{ch:comb} вы помните, что есть
\emph{комбинационные} схемы (выход зависит только от текущих входов) и
\emph{последовательностные}, у которых есть память. Для памяти рядом с
каждым LUT стоит \term{D-триггер}.

Вспомним (глава~\ref{ch:ff}): D-триггер запоминает значение входа $D$ в
момент \emph{переднего фронта} тактового сигнала (перехода $0\to1$) и держит
его на выходе $Q$ до следующего фронта.

\begin{figure}[H]
\centering
\begin{tikztimingtable}[timing/slope=0.1, timing/coldist=2pt, xscale=2.6, yscale=1.4,
  timing/font=\sffamily\small, timing/rowdist=3.2ex]
  CLK & 8{1L 1H} \\
  D   & 2L 4H 4L 2H 2L 2H \\
  Q   & 3L 4H 4L 2H 2L 1H \\
\extracode
  \begin{pgfonlayer}{background}
    \foreach \t in {1,3,...,15} \draw[FpgaOrange, dashed, thin] (\t,0.6) -- (\t,-5.0);
  \end{pgfonlayer}
  \node[lbl, text=FpgaOrange, anchor=south] at (8,0.6) {пунктир --- передние фронты CLK};
\end{tikztimingtable}
\caption{Временная диаграмма D-триггера: выход $Q$ меняется только по
переднему фронту CLK и принимает значение, которое было на $D$ в этот момент}
\label{fig:dff}
\end{figure}

\section{Логическая ячейка CFU}

LUT и триггеры объединяются в \term{логические ячейки}. У Gowin такая ячейка
называется \term{CFU} (Configurable Function Unit). Внутри CFU есть четыре
«ломтика» CLS (Configurable Logic Slice); в каждом по два LUT4, а в трёх из них
ещё и по два триггера. Итого в одной CFU: 8 LUT4 и 6 триггеров.

Проверим: в ПЛИС GW2AR-18 содержится $2592$ ячеек CFU, т.\,е.
$2592\cdot 8 = 20\,736$ LUT4 и $2592\cdot 6 = 15\,552$ триггеров. Эти числа и
указаны в документации на микросхему.

\begin{figure}[H]
\centering
\begin{tikzpicture}[
  lut/.style={draw=FpgaBlue, fill=FpgaLight, thick, minimum width=16mm, minimum height=13mm, font=\sffamily\small, align=center},
  ff/.style={draw=FpgaGreen, fill=green!8, thick, minimum width=11mm, minimum height=13mm, font=\sffamily\small},
  mx/.style={draw=FpgaDark, fill=orange!15, thick, trapezium, trapezium angle=70, shape border rotate=270, minimum width=10mm, minimum height=5mm, font=\sffamily\tiny},
]
  \node[lut] (lut) at (0,0) {LUT4\\\scriptsize 16 бит};
  \foreach \i/\n in {1/A,2/B,3/C,4/D} {
    \draw[wire] ([yshift={(2.5-\i)*2.6mm}]lut.west) -- ++(-1,0) node[left, font=\sffamily\small] {\n};
  }
  \node[ff] (ff) at (3.6,0) {D\quad Q};
  \draw[wire] (lut.east) -- node[above, lbl] {F} (ff.west);
  \draw[wire] ([xshift=-3mm]ff.south) -- ++(0,-0.8) node[below left=0mm, lbl] {CLK};
  \draw[wire] ([xshift=3mm]ff.south) -- ++(0,-0.8) node[below right=0mm, lbl] {CE, SR};
  \node[mx] (m) at (6.1,0) {MUX};
  \draw[wire] (ff.east) -- ([yshift=-1.2mm]m.west);
  \draw[wire] (1.6,0) |- (5.0,1.1) |- ([yshift=1.2mm]m.west);
  \fill (1.6,0) circle (1.3pt);
  \draw[wire] (m.east) -- ++(1,0) node[right, font=\sffamily\small] {Выход};
  \draw[wire, FpgaGray] (m.south) -- ++(0,-0.7) node[below, lbl] {бит конфигурации};
  \node[lbl, align=center] at (3.1,1.45) {комбинационный выход (в обход триггера)};
  \draw[dashed, FpgaGray, rounded corners] (-2.0,-1.9) rectangle (8.8,1.9);
  \node[lbl, anchor=north west] at (-2.0,1.9) {половина CLS};
\end{tikzpicture}
\caption{Базовый элемент логической ячейки: LUT4, D-триггер и мультиплексор,
который выбирает, пойдёт ли на выход сигнал с триггера или прямо с LUT}
\label{fig:cls}
\end{figure}

\section{Трассировочные ресурсы}

Тысячи ячеек нужно соединять друг с другом. Для этого между ними проложены
\term{трассировочные каналы} --- пучки проводников разной длины (короткие для
соседей, длинные для дальних блоков). На пересечениях стоят
\term{коммутационные матрицы}: транзисторные ключи, каждым из которых
управляет свой бит конфигурации.

\begin{figure}[H]
\centering
\begin{tikzpicture}
  \foreach \i in {0,1,2,3} {
    \draw[wire, FpgaGray] (\i*0.5,-0.5) -- (\i*0.5,2.5);
    \draw[wire, FpgaGray] (-0.5,\i*0.5+0.25) -- (2.0,\i*0.5+0.25);
  }
  \foreach \x in {0,1,2,3} \foreach \y in {0,1,2,3}
    \draw[FpgaGray] (\x*0.5,\y*0.5+0.25) circle (0.08);
  \foreach \x/\y in {0/1, 2/3, 3/0}
    \fill[FpgaOrange] (\x*0.5,\y*0.5+0.25) circle (0.1);
  \draw[very thick, FpgaOrange] (-0.5,0.75) -- (0,0.75) -- (0,2.5);
  \draw[very thick, FpgaOrange] (1.0,-0.5) -- (1.0,1.75) -- (2.0,1.75);
  \draw[very thick, FpgaOrange] (1.5,-0.5) -- (1.5,0.25) -- (2.0,0.25);
  \node[lbl, align=left, anchor=west] at (2.6,1.9) {\textcolor{FpgaOrange}{$\bullet$} ключ замкнут (бит = 1)};
  \node[lbl, align=left, anchor=west] at (2.6,1.3) {$\circ$ ключ разомкнут (бит = 0)};
  \node[lbl, align=left, anchor=west] at (2.6,0.5) {Оранжевые линии --- три\\ реальных «провода» схемы};
\end{tikzpicture}
\caption{Коммутационная матрица: замкнутые ключи образуют нужные соединения}
\end{figure}

\begin{infobox}[Любопытный факт]
Большая часть площади кристалла ПЛИС занята вовсе не логикой, а
трассировкой и битами конфигурации, которые ею управляют. Поэтому ПЛИС
медленнее и дороже ASIC той же сложности: за гибкость приходится платить.
\end{infobox}

\section{Специализированные блоки}

Кроме универсальной логики, в ПЛИС есть готовые «кирпичи» для частых задач.
Сделать то же самое на LUT можно, но получится медленнее и займёт много места.

\begin{description}
  \item[IOB (блоки ввода-вывода)] связывают внутреннюю логику с ножками
    микросхемы. В IOB можно выбрать стандарт напряжения (LVCMOS 3,3~В,
    1,8~В, LVDS\ldots), силу тока, подтяжку к питанию или земле.
  \item[BSRAM (блочная память)] --- готовые блоки статической памяти по 18~Кбит.
    В GW2AR-18 их 46, всего 828~Кбит. Используются для буферов, FIFO, таблиц.
  \item[DSP (цифровая обработка сигналов)] --- аппаратные умножители
    $18\times18$ бит с сумматорами. Их 48 штук; они умножают за один такт.
  \item[PLL (фазовая автоподстройка частоты)] --- умеют из входной частоты
    27~МГц получить почти любую другую: 50, 100, 126, 371~МГц\ldots\
    В GW2AR-18 два PLL.
  \item[Глобальная тактовая сеть] --- специальные провода, которые разносят
    тактовый сигнал по всему кристаллу с минимальным разбросом по времени.
\end{description}

\begin{figure}[H]
\centering
\begin{tikzpicture}
\begin{axis}[
  xbar, bar width=4.5mm, width=0.78\textwidth, height=6.5cm,
  xmode=log, log origin=infty, xmin=1, xmax=2e5, xlabel={Количество (логарифмическая шкала)},
  symbolic y coords={PLL, {Умножители 18x18}, {Блоки BSRAM}, {Триггеры (FF)}, {LUT4}},
  ytick=data, y tick label style={font=\sffamily\small},
  nodes near coords, nodes near coords style={font=\sffamily\small, anchor=west},
  point meta=explicit symbolic,
  label style={font=\sffamily\small}, x tick label style={font=\small, /pgf/number format/1000 sep={\,}},
  axis lines*=left, grid=major, grid style={FpgaGray!20}, enlarge y limits=0.15,
]
  \addplot[fill=FpgaBlue!70, draw=FpgaBlue] coordinates {
    (20736,{LUT4}) [20\,736] (15552,{Триггеры (FF)}) [15\,552]
    (46,{Блоки BSRAM}) [46] (48,{Умножители 18x18}) [48] (2,PLL) [2]};
\end{axis}
\end{tikzpicture}
\caption{Ресурсы ПЛИС GW2AR-18 на плате Tang Nano 20K}
\end{figure}

\section{Где хранится конфигурация}

Память конфигурации ПЛИС Gowin сделана по технологии SRAM, как
оперативная память компьютера: при выключении питания она стирается. Поэтому
у ПЛИС два режима загрузки:
\begin{itemize}
  \item \textbf{в SRAM} --- быстро, удобно для отладки, но до выключения питания;
  \item \textbf{во Flash} --- файл конфигурации записывается во флеш-память,
        и при каждом включении ПЛИС сама загружает из неё свою схему за
        десятки миллисекунд.
\end{itemize}

\begin{ideabox}
Внутри ПЛИС: \textbf{LUT} (любая логическая функция), \textbf{триггеры}
(память на 1 бит), \textbf{трассировка} (провода и ключи), \textbf{IOB}
(ножки), а также готовые блоки \textbf{памяти}, \textbf{умножителей} и
\textbf{PLL}. Конфигурация --- это миллионы битов, которые задают
содержимое каждого LUT и состояние каждого ключа.
\end{ideabox}

\chapter{Плата Sipeed Tang Nano 20K}

\section{Общий вид}

Tang Nano 20K --- маленькая (примерно $22\times54$~мм) и недорогая
отладочная плата компании Sipeed. На ней установлена ПЛИС \textbf{Gowin
GW2AR-LV18QN88C8/I7}, а также всё необходимое, чтобы сразу начать
эксперименты: кнопки, светодиоды, разъём HDMI, слот для карты памяти и встроенный
программатор с разъёмом USB-C.

\begin{figure}[H]
\centering
\begin{tikzpicture}[scale=1.35,
  callout/.style={font=\sffamily\scriptsize, text=FpgaDark, align=left},
  cl/.style={FpgaDark, thin, -{Circle[length=1.3mm]}}]
  % Плата
  \fill[BoardGreen, rounded corners=4pt] (0,0) rectangle (10.8,4.4);
  % Отверстия
  \foreach \p in {(0.3,0.3),(10.5,0.3),(0.3,4.1),(10.5,4.1)} \fill[white] \p circle (0.12);
  % Гребёнки выводов
  \foreach \i in {0,...,19} {
    \fill[yellow!70!orange] (0.8+\i*0.48,0.2) circle (0.09);
    \fill[yellow!70!orange] (0.8+\i*0.48,4.2) circle (0.09);
  }
  % USB-C
  \fill[gray!60, rounded corners=1pt] (-0.35,1.6) rectangle (0.55,2.8);
  % BL616
  \fill[ChipBlack] (1.1,1.7) rectangle (2.1,2.7);
  \node[white, font=\tiny\sffamily] at (1.6,2.2) {BL616};
  % Flash
  \fill[ChipBlack] (2.5,0.7) rectangle (3.2,1.2);
  % Генератор
  \fill[gray!40] (2.5,3.1) rectangle (3.2,3.6);
  % ПЛИС
  \fill[ChipBlack] (3.8,1.1) rectangle (6.0,3.3);
  \foreach \i in {0,...,9} {
    \fill[gray!60] (3.95+\i*0.2,1.0) rectangle ++(0.08,0.1);
    \fill[gray!60] (3.95+\i*0.2,3.3) rectangle ++(0.08,0.1);
    \fill[gray!60] (3.7,1.25+\i*0.2) rectangle ++(0.1,0.08);
    \fill[gray!60] (6.0,1.25+\i*0.2) rectangle ++(0.1,0.08);
  }
  \node[white, font=\tiny\sffamily, align=center] at (4.9,2.2) {GOWIN\\GW2AR-18\\\scalebox{0.7}{+ SDRAM 64 Мбит}};
  % Светодиоды
  \foreach \i in {0,...,5} \fill[FpgaOrange] (6.6+\i*0.3,3.45) rectangle ++(0.18,0.14);
  % RGB
  \fill[white] (6.7,2.6) rectangle (7.1,3.0);
  % Кнопки
  \fill[gray!30] (6.6,0.6) rectangle (7.2,1.2); \fill[gray!80] (6.9,0.9) circle (0.17);
  \fill[gray!30] (7.6,0.6) rectangle (8.2,1.2); \fill[gray!80] (7.9,0.9) circle (0.17);
  % HDMI
  \fill[gray!55] (9.6,1.2) rectangle (11.0,3.2);
  % TF
  \fill[gray!35] (7.6,1.6) rectangle (9.2,2.8);
  % FPC
  \fill[gray!20] (2.3,-0.1) rectangle (3.5,0.4);

  % Подписи
  \draw[cl] (0.1,4.9) node[callout, above] {USB-C: питание,\\загрузка, UART} -- (0.1,2.4);
  \draw[cl] (1.6,5.1) node[callout, above] {BL616:\\программатор} -- (1.6,2.5);
  \draw[cl] (2.85,4.9) node[callout, above] {генератор\\27 МГц} -- (2.85,3.4);
  \draw[cl] (4.9,4.9) node[callout, above] {ПЛИС GW2AR-18\\со встроенной SDRAM} -- (4.9,3.0);
  \draw[cl] (7.3,4.9) node[callout, above] {6 светодиодов\\(led[0..5])} -- (7.3,3.5);
  \draw[cl] (10.3,4.9) node[callout, above] {HDMI} -- (10.3,3.0);
  \draw[cl] (6.9,-0.7) node[callout, below] {кнопка S1} -- (6.9,0.9);
  \draw[cl] (7.9,-1.3) node[callout, below] {кнопка S2} -- (7.9,0.9);
  \draw[cl] (8.4,-0.7) node[callout, below right=-1mm] {слот\\TF-карты} -- (8.4,2.0);
  \draw[cl] (6.2,-1.3) node[callout, below left=-1mm] {RGB-светодиод\\WS2812} -- (6.9,2.8);
  \draw[cl] (2.85,-0.7) node[callout, below] {Flash 64 Мбит\\(конфигурация)} -- (2.85,0.95);
  \draw[cl] (0.9,-0.7) node[callout, below] {выводы\\GPIO} -- (0.8,0.2);
\end{tikzpicture}
\caption{Расположение основных элементов на плате Tang Nano 20K (схематично)}
\label{fig:board}
\end{figure}

\section{Характеристики}

\begin{table}[H]
\centering
\caption{Основные характеристики Tang Nano 20K}
\label{tab:specs}
\small
\begin{tabularx}{0.95\textwidth}{>{\bfseries}l X}
\toprule
ПЛИС                 & Gowin GW2AR-LV18QN88C8/I7 (семейство GW2AR-18C, корпус QFN88) \\
Логика               & 20\,736 LUT4, 15\,552 триггеров \\
Распределённая память (SSRAM) & 41\,472 бит \\
Блочная память (BSRAM) & 828 Кбит (46 блоков по 18 Кбит) \\
Встроенная SDRAM     & 64 Мбит (8 Мбайт), шина 32 бита --- в том же корпусе, что и ПЛИС \\
Умножители           & 48 штук $18\times18$ \\
PLL                  & 2 \\
Тактовый генератор   & кварцевый, 27 МГц (плюс программируемый генератор MS5351) \\
Flash                & 64 Мбит, SPI --- хранит конфигурацию ПЛИС \\
Программатор         & BL616: USB-JTAG и USB-UART через один разъём USB-C \\
Периферия            & HDMI, TF-карта, 6 светодиодов, 2 кнопки, RGB-светодиод WS2812, \\
                     & разъём для RGB-дисплея, усилитель звука \\
Напряжение на выводах & 3,3 В (LVCMOS33) \\
\bottomrule
\end{tabularx}
\end{table}

\begin{infobox}[Что значит «R» в названии GW2AR]
Буква R означает, что в корпус микросхемы, помимо кристалла ПЛИС, установлен ещё
один кристалл --- память SDRAM. Такая сборка называется SiP
(System in Package, «система в корпусе»). Провода между ПЛИС и SDRAM проложены
внутри корпуса, поэтому на плате их не видно.
\end{infobox}

\section{Структурная схема}

\begin{figure}[H]
\centering
\begin{tikzpicture}[node distance=5mm and 12mm]
  \node[hblock, minimum width=40mm, minimum height=42mm] (fpga) {ПЛИС\\GW2AR-18\\[2mm]\scriptsize 20\,736 LUT4};
  \node[oblock, below=3mm of fpga.south, anchor=north, minimum width=40mm] (sdram) {SDRAM 64 Мбит\\\scriptsize (внутри корпуса)};
  \draw[bus, <->] (fpga.south) -- (sdram.north);
  \begin{scope}[on background layer]
    \node[draw=FpgaGray, dashed, rounded corners, fit=(fpga)(sdram), inner sep=2.5mm, label={[lbl]below:корпус QFN88}] {};
  \end{scope}

  \node[block, left=30mm of fpga.north west, anchor=north east] (osc) {Генератор\\27 МГц};
  \node[block, below=of osc] (bl) {BL616\\USB-JTAG/UART};
  \node[block, below=of bl] (flash) {Flash\\64 Мбит};
  \node[gblock, left=8mm of bl] (usb) {USB-C};
  \draw[arr] (osc.east) -- node[above, lbl] {clk (вывод 4)} (osc.east -| fpga.west);
  \draw[arr, <->] (bl.east) -- node[above, lbl] {JTAG, UART} (bl.east -| fpga.west);
  \draw[arr, <->] (flash.east) -- node[above, lbl] {SPI} (flash.east -| fpga.west);
  \draw[arr, <->] (usb) -- (bl);

  \node[block, right=20mm of fpga.north east, anchor=north west] (led) {6 светодиодов};
  \node[block, below=of led] (btn) {Кнопки S1, S2};
  \node[block, below=of btn] (hdmi) {HDMI};
  \node[block, below=of hdmi] (tf) {TF-карта};
  \node[block, below=of tf] (gpio) {Выводы GPIO};
  \draw[arr] (led.west -| fpga.east) -- (led.west);
  \draw[arr] (btn.west) -- (btn.west -| fpga.east);
  \draw[arr] (hdmi.west -| fpga.east) -- node[above, lbl] {TMDS} (hdmi.west);
  \draw[arr, <->] (tf.west -| fpga.east) -- node[above, lbl] {SPI/SD} (tf.west);
  \draw[arr, <->] (gpio.west) -- (gpio.west -| fpga.east);
\end{tikzpicture}
\caption{Структурная схема платы}
\end{figure}

\section{Выводы, которые понадобятся в примерах}

Каждая ножка ПЛИС имеет номер. Чтобы связать сигнал в коде
(например, \code{led[0]}) с физической ножкой, используют файл ограничений
\code{.cst} (Constraints). Номера выводов для Tang Nano 20K приведены в
табл.~\ref{tab:pins}.

\begin{table}[H]
\centering
\caption{Выводы ПЛИС, используемые в примерах}
\label{tab:pins}
\begin{tabular}{l c l}
\toprule
\textbf{Сигнал} & \textbf{Вывод ПЛИС} & \textbf{Примечание} \\
\midrule
\code{clk}     & 4  & тактовый генератор 27 МГц \\
\code{btn[0]}  & 88 & кнопка S1, нажата $\to$ 1 \\
\code{btn[1]}  & 87 & кнопка S2, нажата $\to$ 1 \\
\code{led[0]}  & 15 & \multirow{6}{*}{\parbox{5.5cm}{светодиоды, горят при \textbf{0}\\ (активный уровень --- низкий)}} \\
\code{led[1]}  & 16 & \\
\code{led[2]}  & 17 & \\
\code{led[3]}  & 18 & \\
\code{led[4]}  & 19 & \\
\code{led[5]}  & 20 & \\
\code{uart\_tx} & 69 & передача в компьютер через BL616 \\
\code{uart\_rx} & 70 & приём из компьютера \\
\code{ws2812}  & 79 & RGB-светодиод \\
\bottomrule
\end{tabular}
\end{table}

\begin{warnbox}[Светодиоды горят от нуля!]
Светодиоды на плате подключены анодом к питанию $+3{,}3$~В, а катодом --- через
резистор к ножке ПЛИС. Ток течёт (и светодиод горит), когда на ножке
\textbf{логический 0}. Поэтому во всех примерах перед выходом стоит
инверсия: \code{assign led = \textasciitilde{}value;}
\end{warnbox}

\begin{figure}[H]
\centering
\begin{circuitikz}[scale=0.9, transform shape]
  \draw (0,4) node[vcc]{$+3{,}3$ В} -- (0,3.4) to[leD, l=LED] (0,2) to[R, l=$R$] (0,0.4) -- (1.8,0.4);
  \draw (1.8,-0.4) rectangle (4.2,1.2);
  \node[font=\sffamily\small] at (3,0.4) {ПЛИС};
  \node[font=\sffamily\scriptsize, anchor=east] at (1.75,0.65) {вывод 15};
  \node[font=\sffamily\small, align=left, anchor=west] at (5,2.4) {на выводе \textbf{0}: ток есть $\to$ горит};
  \node[font=\sffamily\small, align=left, anchor=west] at (5,1.6) {на выводе \textbf{1}: $U_\text{R}=0$ $\to$ не горит};
\end{circuitikz}
\caption{Подключение светодиода с активным низким уровнем}
\end{figure}

\begin{infobox}[Кнопки]
В примерах мы считаем, что нажатая кнопка даёт на выводе логическую 1. Если на вашей
ревизии платы окажется наоборот, достаточно поставить инверсию на входе:
\code{wire pressed = \textasciitilde{}btn[0];}. Проверить это легко с помощью первого примера
(глава~\ref{ch:examples}).
\end{infobox}

\chapter{От текста до работающей схемы}

\section{Маршрут проектирования}

Для программиста путь от программы до работающей прошивки выглядит так:
«исходный код $\to$ компилятор $\to$ машинные команды». У разработчика ПЛИС
шагов больше, и каждый решает свою задачу.

\begin{figure}[H]
\centering
\begin{tikzpicture}[node distance=4mm and 11mm,
  fstep/.style={block, minimum width=26mm, minimum height=15mm},
  file/.style={draw=FpgaGray, fill=white, font=\ttfamily\scriptsize, minimum height=5mm, inner sep=2pt},
]
  \node[fstep] (hdl) {\textbf{1. Описание}\\\scriptsize Verilog (.v)};
  \node[fstep, right=of hdl] (syn) {\textbf{2. Синтез}\\\scriptsize Yosys / GowinSynth};
  \node[fstep, right=of syn] (pnr) {\textbf{3. Размещение}\\\textbf{и трассировка}\\\scriptsize nextpnr / Gowin PnR};
  \node[fstep, right=of pnr] (bit) {\textbf{4. Битстрим}\\\scriptsize gowin\_pack};
  \node[hblock, minimum width=26mm, minimum height=15mm, below=12mm of bit] (load) {\textbf{5. Загрузка}\\\scriptsize openFPGALoader\\\scriptsize Gowin Programmer};
  \node[oblock, minimum width=26mm, minimum height=15mm, left=of load] (board) {\textbf{Tang Nano 20K}};
  \node[gblock, minimum width=26mm, minimum height=15mm, below=12mm of hdl] (sim) {\textbf{Симуляция}\\\scriptsize Icarus Verilog\\\scriptsize GTKWave};
  \node[file, above=7mm of pnr] (cst) {pins.cst};
  \draw[arr] (hdl) -- (syn);
  \draw[arr] (syn) -- node[below, lbl] {netlist} (pnr);
  \draw[arr] (pnr) -- (bit);
  \draw[arr] (bit) -- node[right, lbl] {.fs} (load);
  \draw[arr] (load) -- node[above, lbl] {USB} (board);
  \draw[arr, <->, dashed] (hdl) -- node[right, lbl] {проверка} (sim);
  \draw[arr] (cst) -- (pnr);
  \node[lbl, right=1mm of cst] {номера ножек};
\end{tikzpicture}
\caption{Маршрут проектирования для ПЛИС}
\label{fig:flow}
\end{figure}

\begin{enumerate}
  \item \textbf{Описание.} Вы пишете, \emph{какой должна быть схема}, на языке
        описания аппаратуры (HDL, Hardware Description Language). Мы будем
        использовать самый распространённый из них --- \textbf{Verilog}.
  \item \textbf{Синтез.} Программа-синтезатор читает текст и превращает его в
        \term{список соединений} (netlist): «вот такие LUT, вот такие триггеры,
        вот так соединены». При этом логика упрощается, как при минимизации
        по картам Карно, только автоматически.
  \item \textbf{Размещение и трассировка} (Place \& Route). Каждый LUT и
        триггер из списка получает конкретное место на кристалле, а соединения
        прокладываются по трассировочным каналам. Здесь же используется файл
        \code{.cst}, чтобы порты схемы попали на нужные ножки.
  \item \textbf{Битстрим.} Результат упаковывается в файл конфигурации
        \code{.fs}: те самые биты для каждого LUT и каждого ключа.
  \item \textbf{Загрузка.} Файл по USB передаётся в ПЛИС (в SRAM) или во
        Flash-память на плате.
\end{enumerate}

\section{Что делает синтезатор}

Рассмотрим, как синтезатор превращает текст в «железо». Пусть мы написали:
\begin{vcode}
assign y = (a & b) | ~c;
\end{vcode}

Человек нарисовал бы схему из трёх элементов: И, НЕ и ИЛИ. Но у ПЛИС нет
отдельных элементов --- есть LUT. Функция зависит от трёх переменных,
поэтому синтезатор вычисляет её таблицу истинности и записывает её в
\emph{один} LUT3 (или LUT4 с одним неиспользуемым входом).

\begin{figure}[H]
\centering
\begin{tikzpicture}
  \begin{scope}
    \node[gAND] (and) at (1.3,0.9) {};
    \node[gNOT] (not) at (1.1,-0.5) {};
    \node[gOR] (or) at (3.0,0.2) {};
    \draw[wire] (and.input 1) -- ++(-0.7,0) node[left] {$a$};
    \draw[wire] (and.input 2) -- ++(-0.7,0) node[left] {$b$};
    \draw[wire] (not.input) -- ++(-0.5,0) node[left] {$c$};
    \draw[wire] (and.output) -- ++(0.3,0) |- (or.input 1);
    \draw[wire] (not.output) -- ++(0.5,0) |- (or.input 2);
    \draw[wire] (or.output) -- ++(0.5,0) node[right] {$y$};
    \node[lbl] at (1.8,-1.5) {как задумал человек};
  \end{scope}
  \draw[-{Stealth[length=4mm]}, line width=2pt, FpgaOrange] (4.9,0.2) -- (6.1,0.2) node[midway, above=1mm, lbl] {синтез};
  \begin{scope}[xshift=7cm]
    \node[draw=FpgaBlue, fill=FpgaLight, thick, minimum width=18mm, minimum height=20mm, font=\sffamily\small, align=center] (lut) at (1.5,0.2) {LUT3};
    \draw[wire] ([yshift=5mm]lut.west) -- ++(-0.7,0) node[left] {$a$};
    \draw[wire] (lut.west) -- ++(-0.7,0) node[left] {$b$};
    \draw[wire] ([yshift=-5mm]lut.west) -- ++(-0.7,0) node[left] {$c$};
    \draw[wire] (lut.east) -- ++(0.7,0) node[right] {$y$};
    \node[lbl] at (1.5,-1.5) {как это реализует ПЛИС};
  \end{scope}
  \begin{scope}[xshift=11.2cm, yshift=1.6cm]
    \matrix[matrix of nodes, nodes={font=\ttfamily\scriptsize, inner sep=1.2pt},
            row 1/.style={nodes={font=\sffamily\bfseries\scriptsize}}] (tt) {
      a & b & c & y \\ 0&0&0&1 \\ 0&0&1&0 \\ 0&1&0&1 \\ 0&1&1&0 \\ 1&0&0&1 \\ 1&0&1&0 \\ 1&1&0&1 \\ 1&1&1&1 \\};
    \draw[FpgaGray] (tt-1-1.south west) -- (tt-1-4.south east);
    \node[lbl, below=0mm of tt] {память LUT};
  \end{scope}
\end{tikzpicture}
\caption{Синтез: три логических элемента превращаются в один LUT3 с
таблицей истинности внутри}
\end{figure}

\begin{ideabox}
Синтезатору всё равно, как вы нарисовали бы схему на бумаге. Важно лишь,
\emph{какую функцию} вы описали. Поэтому на Verilog описывают
\emph{поведение} схемы, а подбором элементов занимается программа.
\end{ideabox}

\section{Инструменты}

Для Tang Nano 20K существуют два набора инструментов, и оба бесплатны.

\subsection{Gowin EDA (официальная среда)}

Среду разработки \textbf{Gowin EDA Education} можно скачать с сайта
производителя (\url{https://www.gowinsemi.com}). Порядок работы:
\begin{enumerate}
  \item \textbf{File $\to$ New $\to$ FPGA Design Project}, указать имя проекта.
  \item Выбрать микросхему: серия \textbf{GW2AR}, корпус \textbf{QFN88},
        скорость \textbf{C8/I7}, устройство \textbf{GW2AR-18C}
        (полное имя \code{GW2AR-LV18QN88C8/I7}).
  \item \textbf{File $\to$ New $\to$ Verilog File} --- создать файл модуля
        и написать код.
  \item \textbf{File $\to$ New $\to$ Physical Constraints File} --- создать
        файл \code{.cst} с номерами выводов.
  \item Во вкладке \textbf{Process} запустить \textbf{Synthesize}, затем
        \textbf{Place \& Route}.
  \item Открыть \textbf{Programmer}, выбрать режим \textbf{SRAM Program}
        (для проверки) или \textbf{External Flash Mode} (чтобы схема
        сохранялась после выключения) и нажать \textbf{Program}.
\end{enumerate}

\subsection{Открытые инструменты}

Существует полностью открытый набор программ, работающий в Linux, Windows и macOS
(удобнее всего установить сборку \textbf{OSS CAD Suite}):
\begin{itemize}
  \item \textbf{Yosys} --- синтезатор;
  \item \textbf{nextpnr-himbaechel} --- размещение и трассировка;
  \item \textbf{Apicula} (\code{gowin\_pack}) --- упаковка битстрима для Gowin;
  \item \textbf{openFPGALoader} --- загрузка в плату;
  \item \textbf{Icarus Verilog} и \textbf{GTKWave} --- симуляция и
        просмотр временных диаграмм.
\end{itemize}

Все команды собраны в файл \code{Makefile} в папке с примерами:
\begin{vcode}[bash]
# 1. Синтез
yosys -p "read_verilog blink.v; synth_gowin -top blink -json blink.json"
# 2. Размещение и трассировка
nextpnr-himbaechel --json blink.json --write pnr.json \
    --device GW2AR-LV18QN88C8/I7 --vopt family=GW2A-18C \
    --vopt cst=blink.cst --freq 27
# 3. Битстрим
gowin_pack -d GW2A-18C -o blink.fs pnr.json
# 4. Загрузка: в SRAM ...
openFPGALoader -b tangnano20k blink.fs
# ... или во Flash
openFPGALoader -b tangnano20k -f blink.fs

# То же самое одной командой:
make PROJECT=02_blink TOP=blink load
\end{vcode}

\section{Файл ограничений .cst}
\label{sec:cst}

Код на Verilog ничего не знает о плате: в нём есть только имена портов
(\code{clk}, \code{led}, \code{btn}\ldots). Файл \code{.cst} (Physical
Constraints, «физические ограничения») сообщает трассировщику, к какой ножке
микросхемы подключить каждый порт и как эту ножку настроить.

\begin{figure}[H]
\centering
\begin{tikzpicture}[node distance=4mm and 20mm]
  \node[block, minimum width=30mm, minimum height=22mm, align=left, font=\ttfamily\scriptsize] (v) {module blink (\\\ \ input\ \ clk,\\\ \ output [5:0] led\\);};
  \node[oblock, right=of v, minimum width=36mm, minimum height=22mm, align=left, font=\ttfamily\scriptsize] (c) {IO\_LOC "clk" 4;\\IO\_LOC "led[0]" 15;\\\ldots\\IO\_LOC "led[5]" 20;};
  \node[draw=FpgaDark, fill=ChipBlack, text=white, minimum width=24mm, minimum height=22mm, right=of c, font=\sffamily\small, align=center] (chip) {ПЛИС\\\scriptsize выводы 4, 15..20};
  \draw[arr] (v) -- node[above, lbl] {имена} (c);
  \draw[arr] (c) -- node[above, lbl] {номера ножек} (chip);
  \node[lbl, below=1mm of v] {blink.v};
  \node[lbl, below=1mm of c] {blink.cst};
\end{tikzpicture}
\caption{Файл \code{.cst} связывает имена портов модуля с ножками микросхемы}
\end{figure}

В файле встречаются всего две команды:
\begin{description}
  \item[\code{IO\_LOC "порт" номер;}] --- подключить порт к выводу с указанным
    номером. Номера выводов Tang Nano 20K приведены в табл.~\ref{tab:pins}.
  \item[\code{IO\_PORT "порт" параметр=значение ...;}] --- задать
    электрические параметры вывода.
\end{description}

\begin{table}[H]
\centering
\caption{Основные параметры команды \code{IO\_PORT}}
\label{tab:ioport}
\small
\begin{tabularx}{\textwidth}{l l X}
\toprule
\textbf{Параметр} & \textbf{Значения} & \textbf{Смысл} \\
\midrule
\code{IO\_TYPE} & \code{LVCMOS33}, \code{LVCMOS18}, \ldots & стандарт логических уровней; на Tang Nano 20K почти везде \code{LVCMOS33} (3,3~В), см. главу~\ref{ch:signals} \\
\code{PULL\_MODE} & \code{UP}, \code{DOWN}, \code{NONE}, \code{KEEPER} & встроенный резистор: подтяжка к питанию, к земле, без подтяжки или «удержание» последнего уровня \\
\code{DRIVE} & \code{4}, \code{8}, \code{12}, \code{16}, \code{24} & максимальный выходной ток в миллиамперах (только для выходов) \\
\bottomrule
\end{tabularx}
\end{table}

Правила записи:
\begin{itemize}
  \item каждая команда заканчивается точкой с запятой;
  \item комментарии начинаются с \code{//}, как в Verilog;
  \item имена портов пишутся в кавычках и должны \emph{в точности} совпадать с
        именами в модуле (регистр букв тоже важен);
  \item у шины каждый бит описывается отдельной строкой: \code{"led[0]"},
        \code{"led[1]"}, \ldots
\end{itemize}

Зачем нужна подтяжка? Вход, к которому ничего не подключено (например, вывод
кнопки, пока она отпущена), «висит в воздухе» и ловит наводки: ПЛИС будет
читать то 0, то 1. Резистор \code{PULL\_MODE=DOWN} слабо притягивает такой вход
к земле, и пока кнопка отпущена, на нём гарантированно 0.

\begin{figure}[H]
\centering
\begin{circuitikz}[scale=0.85, transform shape]
  \draw (0,3.2) node[vcc] {$+3{,}3$ В} -- (0,2.8) to[push button, l=S1] (0,1.4) -- (2.4,1.4);
  \draw (1.3,1.4) to[R, l_=$R_\text{подтяжки}$, *-] (1.3,-0.2) node[ground] {};
  \draw (2.4,0.8) rectangle (4.6,2.0);
  \node[font=\sffamily\small] at (3.5,1.4) {ПЛИС};
  \node[font=\sffamily\scriptsize, anchor=south west] at (2.4,1.4) {88};
  \node[font=\sffamily\small, align=left, anchor=west] at (5.2,1.9) {кнопка отпущена: резистор тянет вход к 0};
  \node[font=\sffamily\small, align=left, anchor=west] at (5.2,1.0) {кнопка нажата: вход соединён с 3,3~В $\to$ 1};
\end{circuitikz}
\caption{\code{PULL\_MODE=DOWN}: резистор подтяжки включается внутри ПЛИС,
паять его не нужно}
\end{figure}

\subsection{Файлы .cst в примерах}

У каждого проекта из главы~\ref{ch:examples} есть свой файл ограничений
рядом с кодом, например \code{code/02\_blink/blink.cst}. В нём описаны ровно
те порты, которые есть в модуле, и каждая группа строк прокомментирована.
Вот полный файл для мигалки:

\vfile[c]{\codepath/02_blink/blink.cst}{code/02\_blink/blink.cst}

Кроме того, в \code{code/common/tang\_nano\_20k.cst} собрана сводная
распиновка всех выводов, которые используются в книге. Это справочник: не
подключайте его к проекту целиком, а копируйте из него нужные строки.

\begin{warnbox}[Частые ошибки]
Мы проверили поведение открытого трассировщика \code{nextpnr} на примерах:
\begin{itemize}
  \item \textbf{Порт модуля не описан в \code{.cst}} --- ошибка
        \code{Unconstrained IO}, прошивка не собирается. Так даже лучше: ошибка
        видна сразу.
  \item \textbf{В \code{.cst} описан порт, которого нет в модуле} --- строка
        просто пропускается с сообщением \code{Cell ... not found}, \emph{без
        ошибки}. Обычно это опечатка в имени: тогда настоящий порт остаётся
        неописанным, и срабатывает предыдущая ошибка. Увидев
        \code{Unconstrained IO}, первым делом ищите опечатку в \code{.cst}.
  \item \textbf{Неверный номер вывода} --- ошибка \code{Pin ... not found}.
\end{itemize}
\end{warnbox}

\subsection{Тактовая частота}

Трассировщику нужно знать, на какой частоте будет работать схема, чтобы
проверить, успевают ли сигналы пройти критический путь
(глава~\ref{ch:sync}). Для \code{nextpnr} частота задаётся ключом
\code{--freq 27} (в мегагерцах), и в журнале появляется строка вида
\begin{vcode}[text]
Info: Max frequency for clock 'clk': 256.21 MHz (PASS at 27.00 MHz)
\end{vcode}
В Gowin EDA частоту задают в отдельном файле временных ограничений \code{.sdc}:
\begin{vcode}[text]
create_clock -name clk -period 37.037 [get_ports {clk}]
\end{vcode}
Здесь 37,037~нс --- период сигнала 27~МГц.

\chapter{Основы языка Verilog}

\section{Код описывает схему, а не последовательность действий}

Главное, что нужно понять перед изучением Verilog: \textbf{это не язык
программирования в привычном смысле}. Строчки кода на Verilog не выполняются
одна за другой. Каждая строчка описывает кусок схемы, и все куски
работают \emph{одновременно}, как элементы на макетной плате.

\begin{figure}[H]
\centering
\begin{tikzpicture}
  \node[draw=FpgaBlue, fill=CodeBack, thick, rounded corners, align=left, font=\ttfamily\small, inner sep=3mm] (code) at (0,0) {
  \textcolor{CodeKeyword}{\textbf{assign}} x = a \& b;\\
  \textcolor{CodeKeyword}{\textbf{assign}} y = c | d;\\
  \textcolor{CodeKeyword}{\textbf{assign}} z = x \^{} y;};
  \node[lbl, above=1mm of code] {код на Verilog};
  \begin{scope}[xshift=5.3cm, yshift=-0.2cm]
    \node[gAND] (g1) at (0,0.8) {};
    \node[gOR] (g2) at (0,-0.8) {};
    \node[gXOR] (g3) at (2.2,0) {};
    \draw[wire] (g1.input 1) -- ++(-0.6,0) node[left] {$a$};
    \draw[wire] (g1.input 2) -- ++(-0.6,0) node[left] {$b$};
    \draw[wire] (g2.input 1) -- ++(-0.6,0) node[left] {$c$};
    \draw[wire] (g2.input 2) -- ++(-0.6,0) node[left] {$d$};
    \draw[wire] (g1.output) -- node[above, lbl] {x} ++(0.4,0) |- (g3.input 1);
    \draw[wire] (g2.output) -- node[below, lbl] {y} ++(0.4,0) |- (g3.input 2);
    \draw[wire] (g3.output) -- ++(0.6,0) node[right] {$z$};
    \node[lbl] at (1.1,1.8) {схема, которую он описывает};
  \end{scope}
  \draw[arr, FpgaOrange] (code.east) -- (4.0,0);
\end{tikzpicture}
\caption{Три строчки кода --- три логических элемента, работающих одновременно}
\label{fig:code_to_schem}
\end{figure}

Если поменять строчки местами, схема не изменится: провод \code{x} всё равно
соединит выход элемента И со входом «Исключающего ИЛИ».

\section{Модуль}

Любая схема на Verilog оформляется как \term{модуль} (\code{module}). Модуль
похож на микросхему: у него есть имя, входы, выходы и «начинка».

\begin{figure}[H]
\centering
\begin{minipage}{0.52\textwidth}
\begin{vcode}
module half_adder (   // имя модуля
    input  wire a,    // входы
    input  wire b,
    output wire s,    // выходы
    output wire c
);
    assign s = a ^ b; // сумма
    assign c = a & b; // перенос
endmodule
\end{vcode}
\end{minipage}\hfill
\begin{minipage}{0.44\textwidth}
\centering
\begin{tikzpicture}
  \node[draw=FpgaBlue, fill=FpgaLight, very thick, minimum width=26mm, minimum height=24mm, font=\sffamily\small, align=center] (m) {half\_adder};
  \draw[wire] ([yshift=5mm]m.west) -- ++(-0.9,0) node[left] {a};
  \draw[wire] ([yshift=-5mm]m.west) -- ++(-0.9,0) node[left] {b};
  \draw[wire] ([yshift=5mm]m.east) -- ++(0.9,0) node[right] {s};
  \draw[wire] ([yshift=-5mm]m.east) -- ++(0.9,0) node[right] {c};
  \node[lbl, below=1mm of m] {модуль = «микросхема»};
\end{tikzpicture}
\end{minipage}
\caption{Модуль полусумматора и его «корпус»}
\end{figure}

Разберём ключевые слова:
\begin{itemize}
  \item \code{module имя ( ... ); ... endmodule} --- начало и конец описания;
  \item \code{input}, \code{output}, \code{inout} --- направление порта;
  \item \code{wire} --- провод (соединение);
  \item \code{assign} --- «подключить к проводу слева выход схемы справа».
\end{itemize}

\section{Сигналы: wire и reg}

\begin{description}
  \item[wire] --- провод. Не хранит значение, а лишь передаёт то, что на него
    подаёт какая-то схема. Присваивается через \code{assign}.
  \item[reg] --- переменная, которой присваивают значение внутри блока
    \code{always}. Несмотря на название, \code{reg} становится триггером
    \emph{не всегда}: это зависит от того, в каком \code{always} он описан (см.
    раздел~\ref{sec:always}).
\end{description}

\subsection{Шины (многоразрядные сигналы)}

Несколько проводов можно объединить в \term{шину}:
\begin{vcode}
wire [7:0] data;      // 8 проводов: data[7], data[6], ..., data[0]
reg  [5:0] counter;   // 6-разрядный регистр
wire       bit3 = data[3];      // один провод из шины
wire [3:0] low  = data[3:0];    // младшие 4 провода
\end{vcode}

\begin{figure}[H]
\centering
\begin{tikzpicture}
  \foreach \i in {0,...,7} {
    \pgfmathtruncatemacro{\n}{7-\i}
    \node[draw=FpgaBlue, fill=\ifnum\n<4 orange!20\else FpgaLight\fi, minimum width=9mm, minimum height=8mm, font=\ttfamily\small] (b\i) at (\i*0.9,0) {\n};
  }
  \node[left=2mm of b0, font=\ttfamily\small] {data =};
  \draw[decorate, decoration={brace, amplitude=5pt, mirror}, thick] (b4.south west) -- (b7.south east) node[midway, below=6pt, lbl] {data[3:0]};
  \node[above=1mm of b0, lbl] {старший (MSB)};
  \node[above=1mm of b7, lbl] {младший (LSB)};
\end{tikzpicture}
\caption{Шина \code{[7:0]}: нумерация разрядов справа налево}
\end{figure}

\subsection{Запись чисел}

Числа в Verilog записываются в формате \code{<разрядность>'<система><значение>}:

\begin{table}[H]
\centering
\caption{Запись чисел}
\begin{tabular}{l l l}
\toprule
\textbf{Запись} & \textbf{Значение} & \textbf{Пояснение} \\
\midrule
\code{4'b1010}        & 10  & 4 бита, двоичная система \\
\code{8'hFF}          & 255 & 8 бит, шестнадцатеричная \\
\code{6'd42}          & 42  & 6 бит, десятичная \\
\code{1'b0}, \code{1'b1} & 0, 1 & один бит \\
\code{27\_000\_000}   & 27000000 & подчёркивания игнорируются, их ставят для удобства чтения \\
\code{4'bx01z}        & ---  & \code{x} --- неизвестно, \code{z} --- отключено (высокий импеданс) \\
\bottomrule
\end{tabular}
\end{table}

\section{Операторы}

\begin{table}[H]
\centering
\caption{Основные операторы Verilog}
\label{tab:ops}
\small
\begin{tabularx}{\textwidth}{l l X}
\toprule
\textbf{Группа} & \textbf{Операторы} & \textbf{Пример и смысл} \\
\midrule
Побитовые & \code{\textasciitilde{} \& | \^{} \textasciitilde\&  \textasciitilde|} &
  \code{a \& b} --- И над каждой парой битов; \code{\textasciitilde{}a} --- НЕ \\
Логические & \code{! \&\& ||} & \code{(x > 3) \&\& en} --- результат 1 бит (истина/ложь) \\
Свёртка & \code{\&a  |a  \^{}a} & \code{\&a} --- И всех битов шины (1, если все единицы) \\
Арифметика & \code{+ - * / \%} & \code{cnt + 1'b1} --- сумматор \\
Сравнение & \code{== != < > <= >=} & \code{cnt == 25} --- компаратор \\
Сдвиги & \code{<< >>} & \code{x << 1} --- сдвиг влево на 1 \\
Склейка & \code{\{a, b\}} & \code{\{4'b0, x\}} --- склеить шины в одну \\
Повтор & \code{\{n\{a\}\}} & \code{\{6\{1'b1\}\}} = \code{6'b111111} \\
Условие & \code{c ? a : b} & мультиплексор: если \code{c}, то \code{a}, иначе \code{b} \\
\bottomrule
\end{tabularx}
\end{table}

Каждый оператор --- это кусок аппаратуры. \code{a + b} --- это сумматор,
\code{a == b} --- компаратор, \code{c ? a : b} --- мультиплексор.

\begin{figure}[H]
\centering
\begin{tikzpicture}
  \node[draw=FpgaDark, fill=orange!15, thick, trapezium, trapezium angle=65, shape border rotate=270, minimum width=20mm, minimum height=10mm, font=\sffamily\small] (mux) at (0,0) {MUX};
  \draw[wire] ([yshift=3mm]mux.west) -- ++(-1,0) node[left] {a};
  \node[font=\tiny, anchor=west] at ([yshift=3mm, xshift=1mm]mux.west) {1};
  \draw[wire] ([yshift=-3mm]mux.west) -- ++(-1,0) node[left] {b};
  \node[font=\tiny, anchor=west] at ([yshift=-3mm, xshift=1mm]mux.west) {0};
  \draw[wire] (mux.east) -- ++(1,0) node[right] {y};
  \draw[wire, FpgaBlue] (mux.south) -- ++(0,-0.6) node[below] {c};
  \node[font=\ttfamily\small, anchor=west] at (3,0) {assign y = c ? a : b;};
\end{tikzpicture}
\caption{Тернарный оператор \code{?:} --- это мультиплексор}
\end{figure}

\section{Блоки always}
\label{sec:always}

Описывать сложную логику одними \code{assign} неудобно. Для этого есть блок
\code{always}. Он бывает двух видов, и важно их не путать.

\subsection{Комбинационный: always @(*)}

Звёздочка означает: «пересчитывай выход при любом изменении любого входа». Такой
блок описывает \textbf{комбинационную схему без памяти}.
\begin{vcode}
reg [1:0] y;
always @(*) begin
    case (sel)
        2'd0: y = a;
        2'd1: y = b;
        2'd2: y = c;
        default: y = 2'b00;
    endcase
end
\end{vcode}

\subsection{Последовательностный: always @(posedge clk)}

Этот блок срабатывает только по \textbf{переднему фронту тактового сигнала}.
Все \code{reg}, которым присваивают значение внутри него, становятся
\textbf{триггерами}.
\begin{vcode}
reg [3:0] cnt = 0;
always @(posedge clk) begin
    cnt <= cnt + 1'b1;   // каждый такт прибавляем единицу
end
\end{vcode}

\begin{figure}[H]
\centering
\begin{tikzpicture}
  \node[draw=FpgaGreen, fill=green!8, thick, minimum width=18mm, minimum height=18mm, font=\sffamily\small, align=center] (reg) at (4.2,0) {регистр\\\scriptsize 4 бита};
  \node[font=\sffamily\scriptsize, anchor=west] at (reg.west) {D};
  \node[font=\sffamily\scriptsize, anchor=east] at (reg.east) {Q};
  \draw[thick] ([xshift=-1.2mm]reg.south) -- ([yshift=1.8mm]reg.south) -- ([xshift=1.2mm]reg.south);
  \draw[wire] (reg.south) -- ++(0,-0.8) node[below] {clk};
  \node[draw=FpgaBlue, fill=FpgaLight, thick, circle, minimum size=11mm, font=\sffamily\bfseries] (add) at (1.2,0) {$+$};
  \draw[bus, -{Stealth}] (add.east) -- node[above, lbl] {cnt+1} (reg.west);
  \draw[bus, -{Stealth}] (reg.east) -- ++(1.6,0) node[right, text=FpgaDark] {cnt};
  \draw[bus, -{Stealth}] ($(reg.east)+(0.8,0)$) -- ++(0,1.5) -| (add.north);
  \fill[FpgaBlue] ($(reg.east)+(0.8,0)$) circle (1.6pt);
  \draw[wire] (add.south) -- ++(0,-0.7) node[below] {1};
\end{tikzpicture}
\caption{Счётчик: сумматор и регистр с обратной связью. Код \code{cnt <= cnt + 1}
описывает именно эту схему}
\label{fig:counter}
\end{figure}

\begin{figure}[H]
\centering
\begin{tikztimingtable}[timing/slope=0.1, xscale=2.2, yscale=1.4, timing/font=\sffamily\small, timing/rowdist=3.2ex]
  clk & 9{1L 1H} \\
  cnt & 1D{0} 2D{1} 2D{2} 2D{3} 2D{4} 2D{5} 2D{6} 2D{7} 2D{8} 1D{9} \\
\end{tikztimingtable}
\caption{Временная диаграмма счётчика: значение увеличивается по каждому
переднему фронту}
\end{figure}

\subsection{Блокирующее (=) и неблокирующее (<=) присваивание}

\begin{warnbox}[Золотое правило]
\begin{itemize}
  \item В \code{always @(posedge clk)} используйте только \code{<=}.
  \item В \code{always @(*)} используйте только \code{=}.
\end{itemize}
\end{warnbox}

Почему? Неблокирующее присваивание \code{<=} означает: «все правые части
вычисляются по значениям \emph{до} фронта, и все триггеры обновляются
одновременно». Именно так ведут себя настоящие триггеры. Посмотрим на
пример, в котором два регистра обмениваются значениями:

\begin{figure}[H]
\centering
\begin{minipage}{0.47\textwidth}
\begin{vcode}
always @(posedge clk) begin
    a <= b;
    b <= a;   // обмен работает!
end
\end{vcode}
\centering
\begin{tikzpicture}[ffn/.style={draw=FpgaGreen, fill=green!8, thick, minimum width=10mm, minimum height=10mm, font=\sffamily\small}]
  \node[ffn] (A) at (0,0) {a};
  \node[ffn] (B) at (2.4,0) {b};
  \draw[arr] (A.north) to[bend left=35] (B.north);
  \draw[arr] (B.south) to[bend left=35] (A.south);
\end{tikzpicture}
\end{minipage}\hfill
\begin{minipage}{0.47\textwidth}
\begin{vcode}
always @(posedge clk) begin
    a = b;
    b = a;    // ошибка: b = b
end
\end{vcode}
\centering
\begin{tikzpicture}[ffn/.style={draw=FpgaRed, fill=red!5, thick, minimum width=10mm, minimum height=10mm, font=\sffamily\small}]
  \node[ffn] (A) at (0,0) {a};
  \node[ffn] (B) at (2.4,0) {b};
  \draw[arr] (B.north) to[bend right=35] (A.north);
  \draw[arr] (B.east) to[out=0, in=-60, looseness=6] (B.south east);
\end{tikzpicture}
\end{minipage}
\caption{Слева --- правильный обмен (\code{<=}), справа --- ошибка (\code{=}):
к моменту второй строки \code{a} уже равно \code{b}}
\end{figure}

\section{Условия: if и case}

Внутри \code{always} можно использовать условия. Они тоже превращаются в
мультиплексоры.
\begin{vcode}
always @(posedge clk) begin
    if (reset)
        cnt <= 0;               // синхронный сброс
    else if (enable)
        cnt <= cnt + 1'b1;      // счёт
    // иначе cnt сохраняет значение
end
\end{vcode}

\begin{warnbox}[Защёлки (latch)]
Если в \code{always @(*)} при каком-то условии переменной ничего не присвоено,
синтезатор вынужден «запомнить» старое значение и создаёт \term{защёлку} ---
нежелательный элемент памяти без такта. Всегда пишите \code{else} и
\code{default}!
\end{warnbox}

\section{Параметры и константы}

\begin{vcode}
localparam integer CLK_HZ = 27_000_000;  // константа внутри модуля
localparam [1:0]   IDLE = 2'd0, RUN = 2'd1;

module blink #(parameter integer HZ = 1) (...);  // параметр можно
                                                 // менять снаружи
\end{vcode}

\section{Иерархия: модуль внутри модуля}

Большие схемы собирают из маленьких модулей, как устройство из микросхем.
Подключение одного модуля внутри другого называется \term{созданием экземпляра}
(instantiation).

\begin{figure}[H]
\centering
\begin{minipage}{0.55\textwidth}
\begin{vcode}
module full_adder (
    input  wire a, b, cin,
    output wire s, cout
);
    wire s1, c1, c2;
    half_adder ha1 (.a(a),  .b(b),
                    .s(s1), .c(c1));
    half_adder ha2 (.a(s1), .b(cin),
                    .s(s),  .c(c2));
    assign cout = c1 | c2;
endmodule
\end{vcode}
\end{minipage}\hfill
\begin{minipage}{0.42\textwidth}
\centering
\begin{tikzpicture}[scale=0.85, transform shape,
  ha/.style={draw=FpgaBlue, fill=FpgaLight, thick, minimum width=13mm, minimum height=13mm, font=\sffamily\scriptsize, align=center}]
  \node[ha] (h1) at (0,0.6) {ha1};
  \node[ha] (h2) at (2.3,-0.2) {ha2};
  \node[draw, circle, thick, minimum size=6mm, font=\scriptsize] (or) at (4.2,-1.3) {$\ge1$};
  \draw[wire] ([yshift=3mm]h1.west) -- ++(-0.6,0) node[left] {a};
  \draw[wire] ([yshift=-3mm]h1.west) -- ++(-0.6,0) node[left] {b};
  \draw[wire] ([yshift=3mm]h1.east) -- node[above, lbl] {s1} ++(0.4,0) |- ([yshift=3mm]h2.west);
  \draw[wire] (-1.0,-0.7) node[left] {cin} -| ([xshift=-3mm]h2.west) -- ([yshift=-3mm]h2.west);
  \draw[wire] ([yshift=3mm]h2.east) -- ++(0.8,0) node[right] {s};
  \draw[wire] ([yshift=-3mm]h1.east) -- node[below, lbl] {c1} ++(0.25,0) |- ([yshift=0mm]or.west);
  \draw[wire] ([yshift=-3mm]h2.east) -- node[below, lbl] {c2} ++(0.3,0) |- ([yshift=1.5mm]or.north west);
  \draw[wire] (or.east) -- ++(0.4,0) node[right] {cout};
  \node[draw=FpgaGray, dashed, fit=(h1)(h2)(or), inner sep=5mm, label={[lbl]above:full\_adder}] {};
\end{tikzpicture}
\end{minipage}
\caption{Полный сумматор из двух полусумматоров}
\end{figure}

Запись \code{.a(s1)} означает: «вход \code{a} экземпляра соединить с проводом
\code{s1}».

\begin{ideabox}
\begin{itemize}
  \item \code{assign} и \code{always @(*)} описывают \textbf{комбинационную}
        логику (элементы И, ИЛИ, мультиплексоры, сумматоры);
  \item \code{always @(posedge clk)} описывает \textbf{триггеры} и всё, что
        перед ними;
  \item все блоки работают \textbf{одновременно}.
\end{itemize}
\end{ideabox}

\chapter{Практика: шесть проектов}
\label{ch:examples}

В этой главе мы шаг за шагом соберём шесть схем, от простых логических
элементов до конечного автомата и ШИМ. Каждый проект лежит в своей папке
внутри \code{code/}; файл ограничений для всех проектов общий:
\code{code/common/tang\_nano\_20k.cst}.

\begin{figure}[H]
\centering
\begin{tikzpicture}
\begin{axis}[
  ybar stacked, bar width=10mm, width=0.9\textwidth, height=6.2cm,
  symbolic x coords={Логика, Мигалка, {Бег. огонь}, Кнопка, Светофор, ШИМ},
  xtick=data, ylabel={Число ячеек},
  legend style={at={(0.5,1.03)}, anchor=south, legend columns=3, font=\sffamily\small, draw=none},
  label style={font=\sffamily\small}, tick label style={font=\sffamily\small},
  axis lines*=left, ymin=0, grid=major, grid style={FpgaGray!20},
]
  \addplot[fill=FpgaBlue!70, draw=FpgaBlue] coordinates {(Логика,5) (Мигалка,114) ({Бег. огонь},60) (Кнопка,21) (Светофор,32) (ШИМ,70)};
  \addplot[fill=FpgaGreen!60, draw=FpgaGreen] coordinates {(Логика,0) (Мигалка,26) ({Бег. огонь},29) (Кнопка,31) (Светофор,30) (ШИМ,34)};
  \addplot[fill=FpgaOrange!60, draw=FpgaOrange] coordinates {(Логика,0) (Мигалка,25) ({Бег. огонь},22) (Кнопка,25) (Светофор,35) (ШИМ,49)};
  \legend{LUT1--LUT4, триггеры, ячейки ALU (перенос)}
\end{axis}
\end{tikzpicture}
\caption{Сколько ресурсов занимает каждый проект (по отчёту синтезатора Yosys).
Даже самый «тяжёлый» пример использует меньше 1\,\% ПЛИС}
\label{fig:usage}
\end{figure}

% =============================================================================
\section{Проект 1. Логические элементы на кнопках}

\textbf{Цель:} убедиться, что элементы И, ИЛИ, НЕ, И-НЕ, ИЛИ-НЕ и
«Исключающее ИЛИ» действительно можно «собрать» кодом. Кнопки S1 и S2 будут
входами $a$ и $b$, шесть светодиодов --- выходами шести функций.

\begin{figure}[H]
\centering
\begin{tikzpicture}[yscale=0.95]
  \coordinate (A) at (0,0); \coordinate (B) at (0,-0.4);
  \node[left, font=\sffamily\small] at (-0.1,-2.4) {\begin{tabular}{r}S1 $\to a$\\S2 $\to b$\end{tabular}};
  \draw[very thick, FpgaBlue] (0,0.5) -- (0,-5.4) node[below, lbl] {$a$};
  \draw[very thick, FpgaRed]  (0.4,0.5) -- (0.4,-5.4) node[below, lbl] {$b$};
  \foreach \i/\g/\t in {0/gAND/И, 1/gOR/ИЛИ, 2/gNOT/НЕ, 3/gNAND/И-НЕ, 4/gNOR/ИЛИ-НЕ, 5/gXOR/Искл. ИЛИ} {
    \pgfmathsetmacro{\y}{-\i*1.0}
    \ifnum\i=2
      \node[gNOT] (g\i) at (2.4,\y) {};
      \draw[wire] (g\i.input) -- (0,\y); \fill[FpgaBlue] (0,\y) circle (1.5pt);
    \else
      \node[\g, minimum height=7mm, minimum width=10mm] (g\i) at (2.4,\y) {};
      \draw[wire] (g\i.input 1) -- (0,{\y+0.12}); \fill[FpgaBlue] (0,{\y+0.12}) circle (1.5pt);
      \draw[wire] (g\i.input 2) -- (0.4,{\y-0.12}); \fill[FpgaRed] (0.4,{\y-0.12}) circle (1.5pt);
    \fi
    \node[gNOT, minimum width=6mm, minimum height=4mm] (n\i) at (4.3,\y) {};
    \draw[wire] (g\i.output) -- (n\i.input);
    \draw[wire] (n\i.output) -- (5.4,\y) node[right, font=\sffamily\small] {led[\i]\quad (\t)};
  }
  \node[lbl, align=center] at (4.3,0.7) {инверсия\\для LED};
\end{tikzpicture}
\caption{Схема проекта 1}
\end{figure}

\vfile{\codepath/01_gates/gates.v}{code/01\_gates/gates.v}

\begin{table}[H]
\centering
\caption{Какие светодиоды должны гореть ($\bullet$ --- горит)}
\label{tab:gates}
\begin{tabular}{cc|cccccc}
\toprule
S1 ($a$) & S2 ($b$) & led0 И & led1 ИЛИ & led2 НЕ $a$ & led3 И-НЕ & led4 ИЛИ-НЕ & led5 XOR \\
\midrule
0 & 0 &           &           & $\bullet$ & $\bullet$ & $\bullet$ &           \\
0 & 1 &           & $\bullet$ & $\bullet$ & $\bullet$ &           & $\bullet$ \\
1 & 0 &           & $\bullet$ &           & $\bullet$ &           & $\bullet$ \\
1 & 1 & $\bullet$ & $\bullet$ &           &           &           &           \\
\bottomrule
\end{tabular}
\end{table}

\begin{infobox}[Что получилось после синтеза]
Yosys сообщает, что схема заняла всего \textbf{5 LUT2}, а не 6 элементов
и ещё 6 инверторов. Синтезатор объединил каждую функцию с инверсией на выходе в
один LUT (например, $\overline{a\cdot b}$ стало одной таблицей). А для
\code{led[2]} получилось $\overline{\overline{a}} = a$ --- светодиод подключён
к кнопке просто проводом, и LUT не понадобился вовсе!
\end{infobox}

\begin{taskbox}
\begin{enumerate}
  \item Загрузите проект и по очереди проверьте все строки
        табл.~\ref{tab:gates}.
  \item Замените функцию на \code{led[5]} мажоритарной функцией от трёх
        сигналов: $a$, $b$ и $a \oplus b$. Что получится и почему?
\end{enumerate}
\end{taskbox}

% =============================================================================
\section{Проект 2. Мигающий светодиод}

\textbf{Цель:} познакомиться с тактовым сигналом и счётчиком. Это
«Hello, world!» мира ПЛИС.

Генератор выдаёт \SI{27}{МГц}, то есть $27\,000\,000$ импульсов в секунду.
Человеческий глаз такого мигания не увидит. Нужно \term{поделить частоту}:
считаем такты и каждые $13\,500\,000$ тактов (0,5~с) переключаем
светодиод. Какой разрядности нужен счётчик?
\[
  2^{23} = 8\,388\,608 < 13\,500\,000 < 2^{24} = 16\,777\,216
  \quad\Rightarrow\quad \text{достаточно 24 бит (мы взяли 25 с запасом).}
\]

\begin{figure}[H]
\centering
\begin{tikzpicture}[node distance=10mm]
  \node[hblock] (clk) {Генератор\\27 МГц};
  \node[block, right=of clk, minimum width=30mm] (cnt) {Счётчик\\\scriptsize 0 \ldots 13\,499\,999};
  \node[block, right=of cnt] (cmp) {Сравнение\\\scriptsize cnt == HALF$-$1};
  \node[gblock, right=of cmp] (tff) {T-триггер\\\scriptsize state};
  \node[oblock, below=8mm of tff] (led) {6 LED};
  \draw[arr] (clk) -- (cnt); \draw[bus, -{Stealth}] (cnt) -- (cmp); \draw[arr] (cmp) -- (tff);
  \draw[arr] (tff) -- node[right, lbl] {\textasciitilde state} (led);
  \draw[arr, dashed] (cmp.south) |- node[pos=0.75, below, lbl] {сброс в 0} ($(cnt.south)+(0,-0.6)$) -- (cnt.south);
\end{tikzpicture}
\caption{Структура делителя частоты}
\end{figure}

\vfile{\codepath/02_blink/blink.v}{code/02\_blink/blink.v}

\begin{figure}[H]
\centering
\begin{tikztimingtable}[timing/slope=0.05, xscale=1.6, yscale=1.4, timing/font=\sffamily\small, timing/rowdist=3.3ex]
  clk   & 6{1L 1H} 3S 6{1L 1H} \\
  cnt   & 2D{0} 2D{1} 2D{2} 2D{\ldots} 2D{H$-$2} 2D{H$-$1} 3S 2D{0} 2D{1} 2D{2} 2D{\ldots} 2D{H$-$2} 2D{H$-$1} \\
  state & 12L 3S 12H \\
\extracode
  \node[lbl, text=FpgaOrange] at (13.5,0.7) {$\approx$0,5 с};
\end{tikztimingtable}
\caption{Временная диаграмма мигалки (H = HALF)}
\end{figure}

\begin{figure}[H]
\centering
\begin{tikzpicture}
\begin{axis}[
  width=0.85\textwidth, height=4.5cm, xlabel={Время, с}, ylabel={led},
  ymin=-0.2, ymax=1.3, xmin=0, xmax=3, ytick={0,1}, yticklabels={горит (0), не горит (1)},
  label style={font=\sffamily\small}, tick label style={font=\sffamily\small},
  grid=major, grid style={FpgaGray!20}, axis lines*=left,
]
  \addplot[very thick, FpgaBlue, const plot] coordinates {(0,0) (0.5,1) (1,0) (1.5,1) (2,0) (2.5,1) (3,1)};
\end{axis}
\end{tikzpicture}
\caption{Сигнал на ножке светодиода: период 1 с, частота 1 Гц}
\end{figure}

\begin{taskbox}
\begin{enumerate}
  \item Сделайте, чтобы светодиод мигал 5 раз в секунду.
  \item Сделайте, чтобы \code{led[0]} мигал с частотой 1 Гц, \code{led[1]} --- 2 Гц,
        \code{led[2]} --- 4 Гц\ldots\ Подсказка: используйте разные биты одного
        счётчика. Какой бит 25-разрядного счётчика меняется с частотой около 1 Гц?
\end{enumerate}
\end{taskbox}

% =============================================================================
\section{Проект 3. Бегущий огонь}

\textbf{Цель:} изучить сдвиговый регистр и сигнал разрешения (\code{tick}).

Шесть триггеров соединены в кольцо: по каждому импульсу \code{tick} значение
каждого триггера переходит в соседний, а значение последнего --- в первый.
Если в начале в кольце одна единица, она будет «бегать» по кругу.

\begin{figure}[H]
\centering
\begin{tikzpicture}[ffn/.style={draw=FpgaGreen, fill=green!8, thick, minimum width=13mm, minimum height=11mm, font=\sffamily\small, align=center}]
  \foreach \i in {0,...,5} {
    \node[ffn] (r\i) at (\i*2.2,0) {D\enskip Q\\[-1mm]\scriptsize ring[\i]};
    \draw[wire] (r\i.north) -- ++(0,0.6) node[above, font=\sffamily\scriptsize] {led[\i]};
  }
  \foreach \i [evaluate=\i as \j using int(\i+1)] in {0,...,4}
    \draw[arr] (r\i.east) -- (r\j.west);
  \draw[arr] (r5.east) -- ++(0.4,0) |- ($(r0.south)+(0,-0.7)$) -| ($(r0.west)+(-0.4,0)$) -- (r0.west);
  \draw[wire, FpgaBlue] ($(r0.south)+(0,-1.3)$) node[left, font=\sffamily\scriptsize] {clk, tick} -- ($(r5.south)+(0,-1.3)$);
  \foreach \i in {0,...,5} \draw[wire, FpgaBlue] (r\i.south) -- ++(0,-1.3);
\end{tikzpicture}
\caption{Кольцевой сдвиговый регистр из шести D-триггеров}
\end{figure}

\vfile{\codepath/03_running_lights/running_lights.v}{code/03\_running\_lights/running\_lights.v}

\begin{figure}[H]
\centering
\begin{tikztimingtable}[timing/slope=0.1, xscale=1.35, yscale=1.35, timing/font=\sffamily\small, timing/rowdist=3.1ex]
  tick    & 2L 1H 3L 1H 3L 1H 3L 1H 3L 1H 3L 1H 3L 1H 1L \\
  ring[0] & 3H 20L 4H 1L \\
  ring[1] & 3L 4H 21L \\
  ring[2] & 7L 4H 17L \\
  ring[3] & 11L 4H 13L \\
  ring[4] & 15L 4H 9L \\
  ring[5] & 19L 4H 5L \\
\end{tikztimingtable}
\caption{Единица «переезжает» в соседний разряд после каждого импульса \code{tick}
(масштаб по времени условный)}
\end{figure}

\begin{infobox}[Почему tick, а не отдельный медленный такт?]
Можно было бы подать на триггеры самодельный медленный такт, взятый прямо с
делителя. Но в ПЛИС так делать не принято: у такого сигнала нет доступа к
глобальной тактовой сети, и схема может работать нестабильно. Правильный приём ---
\textbf{все триггеры тактируются одним \code{clk}}, а редкие события
задаются \term{сигналом разрешения} длиной в один такт.
\end{infobox}

\begin{taskbox}
Сделайте огонь, который бегает туда и обратно («Кнайт Райдер»). Подсказка:
нужен дополнительный регистр направления \code{dir}.
\end{taskbox}

% =============================================================================
\section{Проект 4. Кнопка: синхронизация, дребезг, фронт}

\textbf{Цель:} научиться правильно работать с кнопками. Задача простая:
считать нажатия S1 и показывать число в двоичном виде на светодиодах. Но
на пути три препятствия.

\subsection{Препятствие 1: дребезг контактов}

Механические контакты кнопки при нажатии несколько миллисекунд «прыгают», и
вместо одного перехода $0\to1$ получается целая серия. Схема, работающая
на частоте 27~МГц, посчитает каждый «прыжок» как отдельное нажатие.

\begin{figure}[H]
\centering
\begin{tikzpicture}
\begin{axis}[
  width=0.9\textwidth, height=4.6cm, xlabel={Время, мс}, ylabel={Уровень},
  xmin=0, xmax=30, ymin=-0.2, ymax=1.5, ytick={0,1},
  label style={font=\sffamily\small}, tick label style={font=\sffamily\small},
  legend style={font=\sffamily\small, at={(0.5,1.08)}, anchor=south, legend columns=2, draw=none},
  axis lines*=left,
]
  \addplot[thick, FpgaRed, const plot] coordinates {
    (0,0) (5,1) (5.4,0) (5.9,1) (6.2,0) (6.9,1) (7.1,0) (7.8,1) (8.4,0) (8.6,1)
    (22,1) (22.3,0) (22.7,1) (23.2,0) (23.5,1) (24.0,0) (30,0)};
  \addlegendentry{кнопка (сырой сигнал)}
  \addplot[very thick, FpgaGreen, const plot] coordinates {(0,-0.05) (18.6,-0.05) (18.6,1.1) (30,1.1)};
  \addlegendentry{после антидребезга}
  \draw[FpgaOrange, <->, thick] (axis cs:8.6,1.25) -- (axis cs:18.6,1.25) node[midway, above, font=\sffamily\small] {10 мс стабильно};
\end{axis}
\end{tikzpicture}
\caption{Дребезг контактов и его подавление: новое значение принимается, только
если оно не меняется 10 мс}
\end{figure}

\subsection{Препятствие 2: метастабильность}

Кнопка нажимается в случайный момент, никак не связанный с тактом \code{clk}.
Если вход триггера меняется точно в момент фронта, триггер может на
какое-то время «зависнуть» между 0 и 1 (\term{метастабильное состояние}). Защита
--- \term{синхронизатор} из двух последовательно включённых триггеров: даже
если первый «задумается», ко второму фронту он почти наверняка успокоится.

\begin{figure}[H]
\centering
\begin{tikzpicture}[ffn/.style={draw=FpgaGreen, fill=green!8, thick, minimum width=13mm, minimum height=11mm, font=\sffamily\small}]
  \node[font=\sffamily\small] (in) at (0,0) {btn};
  \node[ffn] (f1) at (2.2,0) {D\enskip Q};
  \node[ffn] (f2) at (4.6,0) {D\enskip Q};
  \node[block] (db) at (7.6,0) {антидребезг\\\scriptsize 10 мс};
  \node[block] (ed) at (10.8,0) {детектор\\фронта};
  \draw[arr] (in) -- (f1); \draw[arr] (f1) -- node[above, lbl] {sync0} (f2);
  \draw[arr] (f2) -- node[above, lbl] {sync1} (db); \draw[arr] (db) -- node[above, lbl] {s1} (ed);
  \draw[arr] (ed) -- ++(1.6,0) node[right, font=\sffamily\small] {press};
  \draw[wire, FpgaBlue] (1.2,-1.2) node[left, font=\sffamily\scriptsize] {clk} -- (10.8,-1.2);
  \foreach \n in {f1,f2,db,ed} \draw[wire, FpgaBlue] (\n.south) -- (\n.south |- 0,-1.2);
\end{tikzpicture}
\caption{Цепочка обработки сигнала кнопки}
\end{figure}

\subsection{Препятствие 3: нажатие длится миллионы тактов}

Если писать \code{if (s1) count <= count + 1;}, счётчик будет прибавлять
единицу \emph{каждый такт}, пока кнопка нажата, --- миллионы раз.
Нужен \term{детектор переднего фронта}: сравниваем текущее значение с
задержанным на такт. Импульс \code{press = s1 \& \textasciitilde{}s1\_prev} длится
ровно один такт.

\begin{figure}[H]
\centering
\begin{tikztimingtable}[timing/slope=0.1, xscale=2.4, yscale=1.4, timing/font=\sffamily\small, timing/rowdist=3.2ex]
  clk      & 8{1L 1H} \\
  s1       & 3L 10H 3L \\
  s1\_prev & 5L 10H 1L \\
  press    & 3L 2H 11L \\
  count    & 5D{5} 11D{6} \\
\end{tikztimingtable}
\caption{Детектор фронта: \code{press} равен 1 только один такт после нажатия}
\end{figure}

\vfile{\codepath/04_button_edge/button_counter.v}{code/04\_button\_edge/button\_counter.v}

\begin{taskbox}
\begin{enumerate}
  \item Сделайте, чтобы S2 не сбрасывал счётчик, а вычитал единицу. Для этого
        S2 тоже понадобятся антидребезг и детектор фронта. Удобно вынести их
        в отдельный модуль \code{debounce} и создать два экземпляра.
  \item Что будет, если убрать антидребезг? Попробуйте и посчитайте, сколько
        лишних «нажатий» набегает.
\end{enumerate}
\end{taskbox}

% =============================================================================
\section{Проект 5. Светофор: конечный автомат}

\textbf{Цель:} освоить главный инструмент проектирования управляющей логики ---
\term{конечный автомат} (FSM, Finite State Machine).

Автомат находится в одном из конечного числа \emph{состояний} и по
определённым условиям \emph{переходит} в другие. Светофор --- классический
пример:

\begin{figure}[H]
\centering
\begin{tikzpicture}[
  st/.style={circle, draw=FpgaDark, very thick, minimum size=22mm, font=\sffamily\small, align=center},
  tr/.style={-{Stealth[length=3mm]}, very thick, FpgaDark},
]
  \node[st, fill=red!20]    (r)  at (0,2.6)  {RED\\\scriptsize 4 с};
  \node[st, fill=orange!25] (ry) at (5.2,2.6) {RED\_YL\\\scriptsize 1 с};
  \node[st, fill=green!20]  (g)  at (5.2,-1.6) {GREEN\\\scriptsize 4 с};
  \node[st, fill=yellow!30] (y)  at (0,-1.6)  {YELLOW\\\scriptsize 1 с};
  \draw[tr] (r)  -- node[above, lbl] {таймер истёк} (ry);
  \draw[tr] (ry) -- node[right, lbl] {таймер истёк} (g);
  \draw[tr] (g)  -- node[below, lbl] {таймер истёк} (y);
  \draw[tr] (y)  -- node[left, lbl] {таймер истёк} (r);
  \draw[tr, FpgaRed, dashed] (-2.6,4.2) node[left, lbl, text=FpgaRed] {btn[0] (сброс)} -- (r);
  % Лампы
  \foreach \s/\cr/\cy/\cg in {r/red/gray!30/gray!30, ry/red/yellow/gray!30, g/gray!30/gray!30/green!70!black, y/gray!30/yellow/gray!30} {
    \begin{scope}[shift={($(\s)+(1.7,0.5)$)}]
    \fill[ChipBlack, rounded corners=1pt] (-0.22,-1.0) rectangle (0.22,0.35);
    \fill[\cr] (0,0.15) circle (0.13); \fill[\cy] (0,-0.3) circle (0.13); \fill[\cg] (0,-0.75) circle (0.13);
    \end{scope}
  }
\end{tikzpicture}
\caption{Граф переходов автомата «Светофор»}
\end{figure}

Автомат на Verilog описывается тремя частями:
\begin{enumerate}
  \item \textbf{регистр состояния} (\code{state}) --- триггеры, хранящие
        номер текущего состояния;
  \item \textbf{логика переходов} --- определяет следующее состояние;
  \item \textbf{выходная логика} --- что зажигать в каждом состоянии.
\end{enumerate}

\begin{figure}[H]
\centering
\begin{tikzpicture}[node distance=12mm]
  \node[block, minimum width=30mm, minimum height=14mm] (next) {Логика\\переходов};
  \node[gblock, right=of next, minimum height=14mm] (st) {Регистр\\state};
  \node[block, right=of st, minimum width=30mm, minimum height=14mm] (out) {Выходная\\логика};
  \draw[arr] (next) -- (st); \draw[arr] (st) -- (out);
  \draw[arr] (out.east) -- ++(1,0) node[right, font=\sffamily\small] {led};
  \draw[arr] ($(st.east)+(0.5,0)$) -- ++(0,-1.2) -| (next.south);
  \draw[arr] ($(next.west)+(-1.2,0)$) node[left, font=\sffamily\small] {timer, btn} -- (next.west);
  \draw[wire, FpgaBlue] (st.north) -- ++(0,0.5) node[above, font=\sffamily\scriptsize] {clk};
\end{tikzpicture}
\caption{Структура автомата Мура: выход зависит только от состояния}
\end{figure}

\vfile{\codepath/05_traffic_light/traffic_light.v}{code/05\_traffic\_light/traffic\_light.v}

\begin{figure}[H]
\centering
\begin{tikztimingtable}[timing/slope=0.05, xscale=1.2, yscale=1.4, timing/font=\sffamily\small, timing/rowdist=3.3ex]
  state   & 8D{RED} 2D{R+Y} 8D{GREEN} 2D{YEL} 8D{RED} 2D{R+Y} \\
  красный & 10H 10L 10H \\
  жёлтый  & 8L 2H 8L 2H 8L 2H \\
  зелёный & 10L 8H 12L \\
\extracode
  \draw[FpgaGray] (0,-6.2) -- (30,-6.2);
  \foreach \t/\l in {0/0, 8/4, 10/5, 18/9, 20/10, 28/14, 30/15} {
    \draw[FpgaGray] (\t,-6.1) -- (\t,-6.3) node[below, font=\sffamily\scriptsize] {\l\,с};
  }
\end{tikztimingtable}
\caption{Работа светофора во времени (1 = лампа горит)}
\end{figure}

\begin{taskbox}
Добавьте состояние «мигающий зелёный» (3 раза по 0,5 с) между GREEN и
YELLOW. Нарисуйте новый граф переходов перед тем, как писать код!
\end{taskbox}

% =============================================================================
\section{Проект 6. «Дышащий» светодиод: ШИМ}

\textbf{Цель:} научиться управлять яркостью. У цифровой ножки только два
состояния: 0 и 1. Но если очень быстро переключать светодиод, глаз увидит
\emph{среднюю} яркость. Это \term{широтно-импульсная модуляция} (ШИМ, англ. PWM):
частота постоянна, а меняется доля времени, когда сигнал равен 1,
--- \term{коэффициент заполнения} $D$:
\[
  D = \frac{t_\text{вкл}}{T}\cdot 100\,\%, \qquad
  U_\text{ср} = D \cdot U_\text{пит}.
\]

\begin{figure}[H]
\centering
\begin{tikzpicture}
\begin{groupplot}[
  group style={group size=1 by 3, vertical sep=13mm, xticklabels at=edge bottom},
  width=0.85\textwidth, height=3cm, xmin=0, xmax=4, ymin=-0.15, ymax=1.3,
  ytick={0,1}, axis lines*=left, tick label style={font=\sffamily\small},
  label style={font=\sffamily\small},
]
\nextgroupplot[title={$D=25\,\%$ --- тускло}, title style={font=\sffamily\small}]
  \addplot[very thick, FpgaBlue, const plot] coordinates {(0,1) (0.25,0) (1,1) (1.25,0) (2,1) (2.25,0) (3,1) (3.25,0) (4,0)};
  \addplot[thick, dashed, FpgaOrange, domain=0:4] {0.25};
\nextgroupplot[title={$D=50\,\%$}, title style={font=\sffamily\small}]
  \addplot[very thick, FpgaBlue, const plot] coordinates {(0,1) (0.5,0) (1,1) (1.5,0) (2,1) (2.5,0) (3,1) (3.5,0) (4,0)};
  \addplot[thick, dashed, FpgaOrange, domain=0:4] {0.5};
\nextgroupplot[title={$D=90\,\%$ --- ярко}, title style={font=\sffamily\small}, xlabel={Время, периоды ШИМ}]
  \addplot[very thick, FpgaBlue, const plot] coordinates {(0,1) (0.9,0) (1,1) (1.9,0) (2,1) (2.9,0) (3,1) (3.9,0) (4,0)};
  \addplot[thick, dashed, FpgaOrange, domain=0:4] {0.9};
\end{groupplot}
\end{tikzpicture}
\caption{ШИМ с разным коэффициентом заполнения. Пунктир --- среднее значение,
которое «видит» глаз}
\end{figure}

Как получить ШИМ на ПЛИС? Нужен свободно бегущий счётчик и компаратор:
пока счётчик меньше заданного числа \code{duty}, выход равен 1.

\begin{figure}[H]
\centering
\begin{tikzpicture}
\begin{axis}[
  width=0.9\textwidth, height=5.2cm, xmin=0, xmax=768, ymin=0, ymax=300,
  xlabel={Такты}, ylabel={Значение},
  label style={font=\sffamily\small}, tick label style={font=\sffamily\small},
  legend style={font=\sffamily\small, at={(0.5,1.02)}, anchor=south, legend columns=3, draw=none},
  axis lines*=left,
]
  \addplot[thick, FpgaBlue] coordinates {(0,0) (255,255) (256,0) (511,255) (512,0) (767,255)};
  \addlegendentry{pwm\_cnt}
  \addplot[thick, dashed, FpgaRed] coordinates {(0,96) (768,96)};
  \addlegendentry{duty = 96}
  \addplot[very thick, FpgaGreen, const plot] coordinates {(0,285) (96,265) (256,285) (352,265) (512,285) (608,265) (768,265)};
  \addlegendentry{pwm (выход)}
\end{axis}
\end{tikzpicture}
\caption{Счётчик-«пила» и компаратор: выход \code{pwm} равен 1, пока
\code{pwm\_cnt < duty}}
\end{figure}

Чтобы светодиод «дышал», будем медленно увеличивать \code{duty} от 0 до 255,
а затем уменьшать обратно.

\begin{figure}[H]
\centering
\begin{tikzpicture}
\begin{axis}[
  width=0.85\textwidth, height=4.6cm, xmin=0, xmax=6, ymin=0, ymax=270,
  xlabel={Время, с}, ylabel={duty},
  label style={font=\sffamily\small}, tick label style={font=\sffamily\small},
  axis lines*=left, grid=major, grid style={FpgaGray!20},
]
  \addplot[very thick, FpgaPurple] coordinates {(0,0) (1,255) (2,0) (3,255) (4,0) (5,255) (6,0)};
\end{axis}
\end{tikzpicture}
\caption{Изменение коэффициента заполнения: «вдох» за 1 с и «выдох» за 1 с}
\end{figure}

\vfile{\codepath/06_pwm_breath/pwm_breath.v}{code/06\_pwm\_breath/pwm\_breath.v}

\begin{infobox}[Почему светодиод кажется «неравномерным»]
Глаз воспринимает яркость нелинейно: разница между $D=5\,\%$ и $10\,\%$ заметна
гораздо сильнее, чем между $90\,\%$ и $95\,\%$. Чтобы «дыхание» выглядело
плавным, в настоящих устройствах используют \term{гамма-коррекцию}: таблицу,
которая переводит линейную яркость в нужное значение \code{duty}. Такую таблицу
удобно хранить в блочной памяти BSRAM.
\end{infobox}

\begin{taskbox}
\begin{enumerate}
  \item Сделайте, чтобы шесть светодиодов «дышали» со сдвигом по фазе
        («волна»).
  \item Добавьте управление: кнопка S1 увеличивает яркость на 1/8, S2 ---
        уменьшает (используйте модуль антидребезга из проекта 4).
\end{enumerate}
\end{taskbox}

\chapter{Симуляция: проверяем схему до загрузки}

\section{Зачем нужна симуляция}

Отладка на плате напоминает поиск неисправности с завязанными глазами:
видно только шесть светодиодов, а внутри тысячи сигналов, и меняются они
27 миллионов раз в секунду. Поэтому профессионалы \textbf{сначала
проверяют схему в симуляторе} --- программе, которая вычисляет значения всех
сигналов во времени и рисует временные диаграммы.

\begin{figure}[H]
\centering
\begin{tikzpicture}
\begin{axis}[
  width=0.8\textwidth, height=5.4cm, ybar, bar width=12mm,
  symbolic x coords={Симуляция, {Синтез + трассировка}, {Загрузка в плату}},
  xtick=data, ylabel={Время одной итерации, с}, ymin=0, ymax=80,
  nodes near coords, nodes near coords style={font=\sffamily\small},
  label style={font=\sffamily\small}, tick label style={font=\sffamily\small},
  axis lines*=left, grid=major, grid style={FpgaGray!20},
]
  \addplot[fill=FpgaGreen!60, draw=FpgaGreen] coordinates {(Симуляция,1) ({Синтез + трассировка},40) ({Загрузка в плату},5)};
\end{axis}
\end{tikzpicture}
\caption{Типичное время одного цикла «исправил --- проверил» для небольшого
проекта. В симуляторе видны все внутренние сигналы, а на плате --- только
выходы}
\end{figure}

\section{Тестбенч}

Для симуляции пишут ещё один модуль на Verilog --- \term{тестбенч} (testbench,
«испытательный стенд»). У него нет портов: он сам создаёт входные сигналы,
подключает к ним проверяемый модуль (DUT --- Device Under Test) и
следит за выходами.

\begin{figure}[H]
\centering
\begin{tikzpicture}
  \node[gblock, minimum width=34mm, minimum height=20mm] (dut) {DUT\\\scriptsize проверяемый модуль\\\scriptsize (gates, blink\ldots)};
  \node[block, left=18mm of dut, minimum height=20mm] (gen) {Генератор\\воздействий\\\scriptsize clk, кнопки};
  \node[block, right=18mm of dut, minimum height=20mm] (mon) {Проверка\\\scriptsize \$display,\\\scriptsize VCD-файл};
  \draw[arr] (gen) -- node[above, lbl] {входы} (dut);
  \draw[arr] (dut) -- node[above, lbl] {выходы} (mon);
  \node[draw=FpgaBlue, dashed, thick, rounded corners, fit=(gen)(dut)(mon), inner sep=4mm, label={[font=\sffamily\small\bfseries, text=FpgaBlue]above:модуль tb (тестбенч)}] {};
\end{tikzpicture}
\caption{Устройство тестбенча}
\end{figure}

В тестбенче разрешены конструкции, которые \emph{нельзя} синтезировать в
«железо», но удобно использовать для моделирования:
\begin{itemize}
  \item \code{\#10;} --- подождать 10 единиц времени;
  \item \code{initial begin ... end} --- выполнить действия один раз с начала
        моделирования;
  \item \code{\$display(...)}, \code{\$monitor(...)} --- печать в консоль;
  \item \code{\$dumpfile}, \code{\$dumpvars} --- запись всех сигналов в файл
        \code{.vcd} для просмотра в GTKWave;
  \item \code{\$finish} --- завершить моделирование.
\end{itemize}

\section{Пример 1: таблица истинности проекта 1}

\vfile{\codepath/01_gates/tb_gates.v}{code/01\_gates/tb\_gates.v}

Запуск в терминале:
\begin{vcode}[bash]
iverilog -o tb_gates tb_gates.v gates.v   # компиляция
vvp tb_gates                              # моделирование
\end{vcode}

Результат (это настоящий вывод симулятора):
\begin{vcode}[text]
 a b | AND OR NOT NAND NOR XOR
 0 0 |  0   0   1    1    1   0
 1 0 |  0   1   0    1    0   1
 0 1 |  0   1   1    1    0   1
 1 1 |  1   1   0    0    0   0
\end{vcode}

Сравните с таблицей~\ref{tab:gates_all} из главы~\ref{ch:gates} --- всё совпадает.

\section{Пример 2: мигалка с уменьшенной частотой}

Проверять мигалку на реальной частоте долго: чтобы дождаться одного
переключения, нужно промоделировать $13{,}5$ млн тактов. Поэтому в модуле
\code{blink} частота вынесена в \code{parameter}, и тестбенч задаёт
\code{CLK\_HZ = 20}: светодиод переключается каждые 10 тактов.

\vfile{\codepath/02_blink/tb_blink.v}{code/02\_blink/tb\_blink.v}

Фрагмент вывода:
\begin{vcode}[text]
t =    0.0 нс  cnt =  0  led = 111111
t =   18.5 нс  cnt =  1  led = 111111
...
t =  314.5 нс  cnt =  9  led = 111111
t =  351.5 нс  cnt =  0  led = 000000
t =  388.5 нс  cnt =  1  led = 000000
\end{vcode}

После \code{cnt = 9} (это \code{HALF - 1}) счётчик сбрасывается, и
светодиоды загораются (\code{000000}). Схема работает как задумано.

\section{Временные диаграммы в GTKWave}

Команда \code{gtkwave tb\_blink.vcd} откроет окно, похожее на
рис.~\ref{fig:gtkwave}. Слева выбирают сигналы, справа видны их диаграммы.

\begin{figure}[H]
\centering
\begin{tikzpicture}
  \fill[gray!15] (0,0) rectangle (15,5.4);
  \fill[gray!35] (0,5.0) rectangle (15,5.4);
  \node[font=\sffamily\scriptsize, anchor=west] at (0.1,5.2) {GTKWave --- tb\_blink.vcd};
  \fill[white] (0.2,0.2) rectangle (2.8,4.8);
  \node[font=\sffamily\scriptsize\bfseries, anchor=west] at (0.3,4.55) {Signals};
  \foreach \n/\y in {clk/3.9, dut.cnt[24:0]/2.9, led[5:0]/1.9} \node[font=\ttfamily\scriptsize, anchor=west] at (0.3,\y) {\n};
  \fill[black] (3.0,0.2) rectangle (14.8,4.8);
  \begin{scope}[shift={(3.2,0)}, x=0.46cm]
    \foreach \i in {0,...,11} {
      \draw[green!80!white, thick] (\i*2,3.6) -- (\i*2,4.2) -- (\i*2+1,4.2) -- (\i*2+1,3.6) -- (\i*2+2,3.6);
    }
    \foreach \i/\v in {0/0,1/1,2/2,3/3,4/4,5/5,6/6,7/7,8/8,9/9,10/0,11/1} {
      \draw[green!80!white, thick] (\i*2+0.15,2.6) -- (\i*2+1.85,2.6) -- (\i*2+2,2.9) -- (\i*2+1.85,3.2) -- (\i*2+0.15,3.2) -- (\i*2,2.9) -- cycle;
      \node[white, font=\ttfamily\tiny] at (\i*2+1,2.9) {\v};
    }
    \draw[green!80!white, thick] (0,1.6) -- (20,1.6) -- (20.1,1.9) -- (20,2.2) -- (0,2.2);
    \node[white, font=\ttfamily\tiny] at (10,1.9) {3F};
    \draw[green!80!white, thick] (20.1,1.9) -- (20.2,2.2) -- (24,2.2); \draw[green!80!white, thick] (20.2,1.6) -- (24,1.6);
    \node[white, font=\ttfamily\tiny] at (22,1.9) {00};
    \draw[yellow, dashed] (20.1,0.3) -- (20.1,4.7);
    \node[yellow, font=\sffamily\tiny, anchor=west] at (20.2,0.6) {351,5 нс};
  \end{scope}
\end{tikzpicture}
\caption{Так выглядит результат симуляции в GTKWave (схематично)}
\label{fig:gtkwave}
\end{figure}

\begin{warnbox}[Симуляция не заменяет плату]
Симулятор проверяет \emph{логику}, но не знает о дребезге кнопок, задержках
в проводах и ошибках в файле \code{.cst}. Поэтому последний шаг всегда
проверка на реальной плате.
\end{warnbox}

\chapter{Память SDRAM на 64 Мбит}

\section{Зачем ПЛИС внешняя память}

Триггеры и блочная память BSRAM очень быстрые, но их мало. Для
изображения, звука или большого массива данных нужна память большего объёма.
Именно для этого в корпус GW2AR-18 встроена микросхема SDRAM на
\textbf{64 Мбит = 8 Мбайт}.

\begin{figure}[H]
\centering
\begin{tikzpicture}
\begin{axis}[
  xbar, bar width=5mm, width=0.8\textwidth, height=5.8cm,
  xmode=log, log origin=infty, xmin=1e3, xmax=3e8,
  xlabel={Объём, бит (логарифмическая шкала)},
  symbolic y coords={Триггеры, SSRAM, BSRAM, SDRAM},
  ytick=data, y tick label style={font=\sffamily\small},
  nodes near coords, nodes near coords style={font=\sffamily\small, anchor=west},
  point meta=explicit symbolic,
  label style={font=\sffamily\small}, x tick label style={font=\small},
  axis lines*=left, grid=major, grid style={FpgaGray!20}, enlarge y limits=0.2,
]
  \addplot[fill=FpgaOrange!70, draw=FpgaOrange] coordinates {
    (15552,Триггеры) [$\approx$1,9 Кбайт] (41472,SSRAM) [$\approx$5 Кбайт]
    (847872,BSRAM) [$\approx$103 Кбайт] (67108864,SDRAM) [8 Мбайт]};
\end{axis}
\end{tikzpicture}
\caption{Сколько данных можно хранить в разных видах памяти Tang Nano 20K}
\end{figure}

Например, один кадр изображения $640\times480$ точек по 16 бит на точку
занимает
\[
  640 \cdot 480 \cdot 16 = 4\,915\,200 \text{ бит} \approx 600 \text{ Кбайт},
\]
что в шесть раз больше всей блочной памяти ПЛИС, но легко помещается в
SDRAM: туда влезает больше десятка таких кадров.

\section{Как устроена DRAM}

Ячейка динамической памяти (DRAM) состоит всего из \textbf{одного
транзистора и одного конденсатора}. Заряженный конденсатор хранит 1,
разряженный --- 0. Поэтому DRAM такая плотная и дешёвая. Но у неё есть два
неудобства:
\begin{itemize}
  \item конденсатор постепенно \textbf{разряжается}, поэтому каждую строку
        нужно периодически перезаписывать (\term{регенерация}, refresh) --- все
        строки за 64~мс;
  \item чтение разрушает заряд, поэтому строку сначала копируют в
        быстрый буфер (\term{открывают}), а по окончании работы записывают
        обратно (\term{закрывают}).
\end{itemize}

\begin{figure}[H]
\centering
\begin{circuitikz}[scale=0.9, transform shape]
  \draw (3,1.5) node[nmos] (T) {};
  \draw (0,0 |- T.gate) node[left, font=\sffamily\small] {строка (WL)} -- (T.gate);
  \draw (T.drain) -- (T.drain |- 0,3.4) node[above, font=\sffamily\small] {столбец (BL)};
  \draw (T.source) to[C, l_=$C$] ++(0,-1.6) node[ground]{};
  \node[font=\sffamily\small, align=left, anchor=west] at (4.6,1.2) {1 бит = 1 транзистор + 1 конденсатор\\[1mm]
  \textcolor{FpgaRed}{заряд утекает} $\Rightarrow$ нужна регенерация};
\end{circuitikz}
\caption{Ячейка DRAM (упрощённо)}
\end{figure}

\section{Организация памяти}

SDRAM в GW2AR-18 имеет 32-разрядную шину данных и организована как
$4 \text{ банка} \times 2048 \text{ строк} \times 256 \text{ столбцов}
\times 32 \text{ бита}$:
\[
  4 \cdot 2048 \cdot 256 \cdot 32 = 67\,108\,864 \text{ бит} = 64 \text{ Мбит}.
\]

\begin{figure}[H]
\centering
\begin{tikzpicture}[x=1cm, y=1cm]
  \foreach \b/\c in {3/FpgaBlue!15, 2/FpgaBlue!25, 1/FpgaBlue!35, 0/FpgaBlue!50} {
    \begin{scope}[shift={(\b*0.45,\b*0.45)}]
      \fill[\c, draw=FpgaBlue] (0,0) rectangle (5,3.2);
      \node[font=\sffamily\scriptsize, anchor=north east] at (5,3.2) {банк \b};
    \end{scope}
  }
  \foreach \y in {0.4,0.8,...,3.0} \draw[FpgaBlue!40] (0,\y) -- (5,\y);
  \foreach \x in {0.5,1.0,...,4.6} \draw[FpgaBlue!40] (\x,0) -- (\x,3.2);
  \fill[FpgaOrange!80] (0,1.6) rectangle (5,2.0);
  \fill[FpgaRed] (2.5,1.6) rectangle (3.0,2.0);
  \draw[decorate, decoration={brace, amplitude=5pt}, thick] (-0.2,0) -- (-0.2,3.2) node[midway, left=6pt, font=\sffamily\small, align=right] {2048 строк\\\scriptsize адрес A[10:0]};
  \draw[decorate, decoration={brace, amplitude=5pt, mirror}, thick] (0,-0.2) -- (5,-0.2) node[midway, below=6pt, font=\sffamily\small, align=center] {256 столбцов, адрес A[7:0]};
  \node[font=\sffamily\small, align=left, anchor=west] at (7.4,2.6) {\textcolor{FpgaOrange}{\rule{3mm}{3mm}} открытая строка};
  \node[font=\sffamily\small, align=left, anchor=west] at (7.4,1.9) {\textcolor{FpgaRed}{\rule{3mm}{3mm}} одно 32-битное слово};
  \node[font=\sffamily\small, align=left, anchor=west] at (7.4,1.2) {банк выбирается BA[1:0]};
\end{tikzpicture}
\caption{Организация SDRAM: банк, строка, столбец}
\end{figure}

\section{Команды SDRAM}

В отличие от BSRAM, в SDRAM нельзя просто подать адрес и получить данные.
Микросхема управляется \term{командами}, которые кодируются комбинацией
сигналов \code{CS\#}, \code{RAS\#}, \code{CAS\#}, \code{WE\#} (решётка
означает активный низкий уровень) по фронту тактового сигнала.

\begin{table}[H]
\centering
\caption{Основные команды SDRAM}
\begin{tabular}{l c c c c l}
\toprule
\textbf{Команда} & \textbf{CS\#} & \textbf{RAS\#} & \textbf{CAS\#} & \textbf{WE\#} & \textbf{Действие} \\
\midrule
NOP        & 0 & 1 & 1 & 1 & ничего не делать \\
ACTIVE     & 0 & 0 & 1 & 1 & открыть строку в банке \\
READ       & 0 & 1 & 0 & 1 & прочитать из открытой строки \\
WRITE      & 0 & 1 & 0 & 0 & записать в открытую строку \\
PRECHARGE  & 0 & 0 & 1 & 0 & закрыть строку \\
AUTO REFRESH & 0 & 0 & 0 & 1 & регенерация \\
LOAD MODE  & 0 & 0 & 0 & 0 & настройка (CAS latency, длина пакета) \\
\bottomrule
\end{tabular}
\end{table}

\begin{figure}[H]
\centering
\begin{tikztimingtable}[timing/slope=0.1, xscale=1.9, yscale=1.4, timing/font=\sffamily\small, timing/rowdist=3.3ex]
  CLK      & 9{1L 1H} \\
  Команда  & 2D{ACTIVE} 2D{NOP} 2D{READ} 2D{NOP} 2D{NOP} 2D{PRE} 6D{NOP} \\
  Адрес    & 2D{строка} 2U 2D{столб.} 12U \\
  DQ[31:0] & 10Z 2D{D0} 2D{D1} 4Z \\
\extracode
  \draw[FpgaOrange, <->, thick] (1,-6.6) -- (5,-6.6) node[midway, below, font=\sffamily\scriptsize] {$t_\text{RCD}$};
  \draw[FpgaPurple, <->, thick] (5,-6.6) -- (9,-6.6) node[midway, below, font=\sffamily\scriptsize] {CAS latency = 2};
\end{tikztimingtable}
\caption{Чтение из SDRAM: открыть строку (ACTIVE), подождать $t_\text{RCD}$,
выдать команду READ, через CAS latency тактов получить данные}
\label{fig:sdram_read}
\end{figure}

\section{Контроллер SDRAM}

Помнить обо всех командах и задержках при каждом обращении к памяти неудобно.
Поэтому между логикой пользователя и SDRAM ставят \term{контроллер} ---
модуль (обычно конечный автомат, как в проекте 5), который:
\begin{itemize}
  \item после включения питания инициализирует микросхему;
  \item превращает простые запросы «прочитай слово по адресу X» в
        последовательность команд ACTIVE $\to$ READ $\to$ PRECHARGE;
  \item сам вовремя выполняет регенерацию.
\end{itemize}

\begin{figure}[H]
\centering
\begin{tikzpicture}[node distance=24mm]
  \node[gblock, minimum height=24mm, minimum width=24mm] (user) {Ваша логика\\\scriptsize (видео, звук\ldots)};
  \node[hblock, right=of user, minimum height=24mm, minimum width=26mm] (ctrl) {Контроллер\\SDRAM\\\scriptsize конечный автомат};
  \node[oblock, right=of ctrl, minimum height=24mm, minimum width=20mm] (mem) {SDRAM\\64 Мбит};
  \draw[arr] ([yshift=6mm]user.east) -- node[above, lbl] {addr, wr, rd} ([yshift=6mm]ctrl.west);
  \draw[arr] ([yshift=0mm]user.east) -- node[above, lbl] {din[31:0]} ([yshift=0mm]ctrl.west);
  \draw[arr] ([yshift=-6mm]ctrl.west) -- node[below, lbl] {dout, busy} ([yshift=-6mm]user.east);
  \draw[arr] ([yshift=6mm]ctrl.east) -- node[above, lbl] {команды} ([yshift=6mm]mem.west);
  \draw[arr] ([yshift=0mm]ctrl.east) -- node[above, lbl] {A, BA} ([yshift=0mm]mem.west);
  \draw[arr, <->] ([yshift=-6mm]ctrl.east) -- node[below, lbl] {DQ[31:0]} ([yshift=-6mm]mem.west);
  \node[lbl, below=2mm of user] {простые запросы};
  \node[lbl, below=2mm of mem] {сложный протокол};
\end{tikzpicture}
\caption{Контроллер скрывает сложность протокола SDRAM}
\end{figure}

\subsection{Порты встроенной памяти}

В проекте Gowin встроенная SDRAM подключается к модулю верхнего уровня через
порты со специальными (зарезервированными) именами. Номера выводов в
\code{.cst} для них указывать не нужно: соединение находится внутри корпуса.

\begin{vcode}
module top (
    input  wire        clk,
    // --- Встроенная SDRAM (имена изменять нельзя!) ---
    output wire        O_sdram_clk,
    output wire        O_sdram_cke,
    output wire        O_sdram_cs_n,    // CS#
    output wire        O_sdram_ras_n,   // RAS#
    output wire        O_sdram_cas_n,   // CAS#
    output wire        O_sdram_wen_n,   // WE#
    output wire [3:0]  O_sdram_dqm,     // маска байтов
    output wire [10:0] O_sdram_addr,    // адрес строки/столбца
    output wire [1:0]  O_sdram_ba,      // номер банка
    inout  wire [31:0] IO_sdram_dq      // данные
);
\end{vcode}

\begin{infobox}[Где взять контроллер]
\begin{itemize}
  \item В Gowin EDA есть готовый IP-блок: \textbf{Tools $\to$ IP Core Generator
        $\to$ Memory Interface $\to$ SDRAM Controller}. Мастер сам
        сгенерирует модуль с простым интерфейсом.
  \item Существуют открытые контроллеры, написанные для Tang Nano 20K
        энтузиастами; их используют, например, проекты ретро-приставок NESTang и
        SNESTang.
  \item Хорошее упражнение повышенной сложности --- написать свой контроллер,
        следуя рис.~\ref{fig:sdram_read} и документации на микросхему.
\end{itemize}
\end{infobox}

\section{Пропускная способность}

При тактовой частоте SDRAM, например, 100~МГц и шине 32 бита пиковая
скорость передачи равна
\[
  B = 100 \cdot 10^6 \ \text{Гц} \cdot 32 \ \text{бит} = 3{,}2 \ \text{Гбит/с} = 400 \ \text{Мбайт/с}.
\]
Реальная скорость ниже из-за открытия и закрытия строк и регенерации.
Чем длиннее пакеты последовательных адресов, тем ближе скорость к пиковой.

\begin{figure}[H]
\centering
\begin{tikzpicture}
\begin{axis}[
  width=0.82\textwidth, height=5.6cm,
  xlabel={Длина пакета, слов}, ylabel={Эффективная скорость, Мбайт/с},
  xmin=1, xmax=256, xmode=log, log basis x=2, ymin=0, ymax=420,
  xtick={1,2,4,8,16,32,64,128,256}, xticklabels={1,2,4,8,16,32,64,128,256},
  label style={font=\sffamily\small}, tick label style={font=\sffamily\small},
  grid=major, grid style={FpgaGray!20}, axis lines*=left,
]
  \addplot[very thick, FpgaBlue, samples at={1,2,4,8,16,32,64,128,256}, mark=*] {400*x/(x+6)};
  \addplot[thick, dashed, FpgaRed, domain=1:256] {400};
  \node[font=\sffamily\small, text=FpgaRed, anchor=north east] at (axis cs:256,395) {пиковая 400 Мбайт/с};
\end{axis}
\end{tikzpicture}
\caption{Оценка скорости обмена при 100 МГц: на каждый пакет приходится около
6 «служебных» тактов (ACTIVE, $t_\text{RCD}$, CAS latency, PRECHARGE)}
\end{figure}

\begin{ideabox}
SDRAM даёт много памяти (8 Мбайт), но требует контроллера: строки нужно
открывать и закрывать, а память регулярно регенерировать. Самая высокая
скорость получается, когда данные читаются и пишутся длинными
последовательными пакетами.
\end{ideabox}

\chapter{Задания для самостоятельной работы}

Задания расположены по возрастанию сложности. Звёздочками отмечена
сложность: \textcolor{FpgaOrange}{$\star$} --- разминка,
\textcolor{FpgaOrange}{$\star\star$} --- нужно подумать,
\textcolor{FpgaOrange}{$\star\star\star$} --- мини-проект.

\begin{figure}[H]
\centering
\begin{tikzpicture}[node distance=5mm and 7mm,
  t/.style={block, minimum width=24mm, minimum height=11mm, font=\sffamily\scriptsize}]
  \node[t] (a) {1--3\\комбинационная\\логика};
  \node[t, right=of a] (b) {4--6\\счётчики\\и регистры};
  \node[t, right=of b] (c) {7--8\\автоматы};
  \node[oblock, minimum width=24mm, minimum height=11mm, font=\sffamily\scriptsize, right=of c] (d) {9--10\\интерфейсы};
  \node[hblock, minimum width=24mm, minimum height=11mm, font=\sffamily\scriptsize, right=of d] (e) {11--12\\проекты};
  \draw[arr] (a) -- (b); \draw[arr] (b) -- (c); \draw[arr] (c) -- (d); \draw[arr] (d) -- (e);
\end{tikzpicture}
\caption{Карта заданий}
\end{figure}

\begin{taskbox}[Задание 1 \quad \textcolor{FpgaOrange}{$\star$}]
Реализуйте на двух кнопках \textbf{полусумматор}: \code{led[0]} --- сумма,
\code{led[1]} --- перенос. Составьте таблицу истинности и проверьте её на плате.
\end{taskbox}

\begin{taskbox}[Задание 2 \quad \textcolor{FpgaOrange}{$\star$}]
Выразите функции И, ИЛИ и НЕ \textbf{только через И-НЕ} и опишите получившуюся
схему операторами \code{\textasciitilde{}(... \& ...)}. Сколько LUT займёт
результат по отчёту синтезатора? Объясните, почему.
\end{taskbox}

\begin{taskbox}[Задание 3 \quad \textcolor{FpgaOrange}{$\star\star$}]
Сделайте \textbf{шифратор}: при нажатии одной кнопки на светодиодах
отображается двоичный код 1, другой --- 2, обеих --- 3. Используйте
\code{always @(*)} и \code{case}. Не забудьте \code{default}!
\end{taskbox}

\begin{taskbox}[Задание 4 \quad \textcolor{FpgaOrange}{$\star$}]
Двоичный счётчик на 6 светодиодах, прибавляющий единицу 4 раза в секунду.
Через сколько секунд он переполнится?
\end{taskbox}

\begin{taskbox}[Задание 5 \quad \textcolor{FpgaOrange}{$\star\star$}]
Счётчик в \textbf{коде Грея}: соседние значения отличаются ровно одним битом.
Подсказка: \code{gray = bin \^{} (bin >> 1)}. Сравните мигание со
счётчиком из задания 4.
\end{taskbox}

\begin{taskbox}[Задание 6 \quad \textcolor{FpgaOrange}{$\star\star$}]
Бегущий огонь, у которого кнопка S1 меняет направление, а S2 --- скорость
(3 скорости по кругу).
\end{taskbox}

\begin{taskbox}[Задание 7 \quad \textcolor{FpgaOrange}{$\star\star$}]
\textbf{Кодовый замок}. Правильная последовательность нажатий:
S1, S1, S2, S1. После верного ввода загораются все светодиоды, после
ошибки --- мигает \code{led[5]}, и ввод начинается заново. Нарисуйте граф
переходов автомата.
\end{taskbox}

\begin{taskbox}[Задание 8 \quad \textcolor{FpgaOrange}{$\star\star$}]
\textbf{Секундомер}: S1 --- старт/стоп, S2 --- сброс. На светодиодах ---
число секунд в двоичном виде (до 63).
\end{taskbox}

\begin{taskbox}[Задание 9 \quad \textcolor{FpgaOrange}{$\star\star\star$}]
\textbf{Передатчик UART}. Раз в секунду отправляйте в компьютер символ
\code{'A'} на скорости 115\,200 бод (вывод 69). Кадр UART: стартовый бит (0),
8 бит данных младшим битом вперёд, стоповый бит (1). Длительность бита ---
$27\,000\,000 / 115\,200 \approx 234$ такта. Смотрите результат в любом
терминале (PuTTY, \code{screen}, Arduino IDE).

\centering
\begin{tikztimingtable}[timing/slope=0.1, xscale=1.35, yscale=1.4, timing/font=\sffamily\small]
  TX & 2H 2L 2H 2L 2L 2L 2L 2L 2H 2L 2H 2H \\
\extracode
  \foreach \x/\t in {3/старт, 5/D0, 7/D1, 9/D2, 11/D3, 13/D4, 15/D5, 17/D6, 19/D7, 21/стоп}
    \node[font=\sffamily\tiny] at (\x,1.3) {\t};
\end{tikztimingtable}
\par\small Передача символа 'A' = \code{0x41} = \code{0100\,0001}$_2$
\end{taskbox}

\begin{taskbox}[Задание 10 \quad \textcolor{FpgaOrange}{$\star\star\star$}]
\textbf{Приёмник UART}: принятый из компьютера байт показывайте на
светодиодах (младшие 6 бит). Вход --- вывод 70. Не забудьте синхронизатор!
\end{taskbox}

\begin{taskbox}[Задание 11 \quad \textcolor{FpgaOrange}{$\star\star\star$}]
\textbf{Игра «Реакция»}. Через случайное время (используйте
сдвиговый регистр с линейной обратной связью, LFSR) загорается светодиод.
Игрок должен как можно быстрее нажать S1. Время реакции в десятках миллисекунд
выводится на светодиоды или отправляется по UART.
\end{taskbox}

\begin{taskbox}[Задание 12 \quad \textcolor{FpgaOrange}{$\star\star\star$}]
\textbf{Генератор ШИМ для сервопривода}. Сервоприводы и регуляторы моторов
дронов управляются импульсами длительностью 1--2~мс с периодом 20~мс.
Сделайте генератор, в котором кнопки S1 и S2 меняют длительность импульса с
шагом 0,1~мс, и выведите сигнал на любой свободный вывод GPIO. Проверьте
осциллографом или с настоящим сервоприводом.
\end{taskbox}


\appendix
\chapter{Шпаргалка}

\section{Verilog на одной странице}

\begin{vcode}
module name #(parameter integer N = 8) (   // параметр
    input  wire         clk,
    input  wire [N-1:0] a,
    output wire         y,
    output reg  [N-1:0] q
);
    localparam [1:0] S0 = 2'd0;            // константа

    wire t = &a;                           // свёртка И
    assign y = t ? 1'b1 : 1'b0;            // мультиплексор

    always @(*) begin                      // комбинационная логика: =
        case (a[1:0])
            2'd0:    ...;
            default: ...;
        endcase
    end

    always @(posedge clk) begin            // триггеры: <=
        if (a == 0) q <= 0;
        else        q <= q + 1'b1;
    end

    other_module u1 (.in(a), .out(...));   // экземпляр модуля
endmodule
\end{vcode}

\section{Типичные ошибки}

\begin{longtable}{p{0.42\textwidth} p{0.52\textwidth}}
\toprule
\textbf{Симптом} & \textbf{Вероятная причина} \\
\midrule
\endhead
Светодиоды горят «наоборот» & Забыли, что LED горят от 0: нужна инверсия \\
Ничего не работает, схема «пустая» & Выход модуля ни к чему не подключён, и синтезатор удалил всю логику; неверные имена в \code{.cst} \\
Счётчик считает одно нажатие много раз & Нет антидребезга и/или детектора фронта \\
Предупреждение \emph{latch inferred} & В \code{always @(*)} не хватает \code{else} или \code{default} \\
Два регистра не обмениваются значениями & В \code{always @(posedge clk)} использовано \code{=} вместо \code{<=} \\
Число не помещается, счётчик не доходит до нужного значения & Мала разрядность регистра: проверьте $2^n > N$ \\
\emph{Multiple drivers} & Одному сигналу присваивают значение в двух разных \code{always} или \code{assign} \\
После выключения питания схема пропадает & Загрузили в SRAM; для постоянной работы загрузите во Flash \\
\bottomrule
\end{longtable}

\section{Выводы Tang Nano 20K}

\begin{table}[H]
\centering
\begin{tabular}{l c l}
\toprule
\textbf{Сигнал} & \textbf{Вывод} & \textbf{Назначение} \\
\midrule
\code{clk} & 4 & 27 МГц \\
\code{btn[0]}, \code{btn[1]} & 88, 87 & кнопки S1, S2 \\
\code{led[0..5]} & 15--20 & светодиоды (активный 0) \\
\code{uart\_tx}, \code{uart\_rx} & 69, 70 & UART через USB \\
\code{ws2812} & 79 & RGB-светодиод \\
\bottomrule
\end{tabular}
\end{table}

\section{Полезные ссылки}

\begin{itemize}
  \item Документация Sipeed: \url{https://wiki.sipeed.com/hardware/en/tang/tang-nano-20k/nano-20k.html}
  \item Примеры Sipeed: \url{https://github.com/sipeed/TangNano-20K-example}
  \item Среда Gowin EDA: \url{https://www.gowinsemi.com}
  \item Открытые инструменты (OSS CAD Suite): \url{https://github.com/YosysHQ/oss-cad-suite-build}
  \item Проект Apicula: \url{https://github.com/YosysHQ/apicula}
  \item Интерактивный учебник по Verilog HDLBits: \url{https://hdlbits.01xz.net}
\end{itemize}


\end{document}
